% Interaction lumière-matière : le photon — Physique 1ère TI — Cours signé OnBuch+
% Programme officiel MINESEC. Compiler avec : tectonic N13-interaction-lumiere-matiere-photon.tex
\newcommand{\DOCMATIERE}{Physique}
\newcommand{\DOCNIVEAU}{1ère TI}
\newcommand{\DOCMODULE}{Module 3 : Optique}
\newcommand{\DOCLECON}{N13}
\newcommand{\DOCTITRE}{Interaction lumière-matière : le photon}
% OnBuch+ — préambule « cours signés » ENRICHI (compilé avec tectonic / XeLaTeX)
% Système visuel « Pop » d'OnBuch+ : fond crème (jamais blanc), bordures encre
% épaisses, ombres « dures » décalées, titres Archivo Black en capitales.
% Version MATHÉMATIQUES : graphes (pgfplots/tikz), boîtes pédagogiques
% (définition, propriété, méthode, attention, le savais-tu, exercice résolu,
% exercice, TP, bilan), page de garde et signature OnBuch+ ancrée en bas.
%
% Macros attendues dans main.tex AVANT \input{preamble} :
%   \DOCMATIERE  \DOCNIVEAU  \DOCMODULE  \DOCLECON (numéro)  \DOCTITRE
\documentclass[11pt]{article}

\usepackage{fontspec}
\usepackage[french]{babel}
\usepackage[a4paper,margin=2cm,top=3.4cm,bottom=2.8cm,headheight=1.4cm,footskip=1.3cm]{geometry}
\usepackage{amsmath,amssymb,amsfonts}
\usepackage{mathtools} % \xmapsto, \coloneqq... souvent utilisés par les modèles
\usepackage{cancel} % simplifications barrées (bilans de forces)
% \abs{z} (module, valeur absolue) et \norm{u} (norme), utilisés par les modèles
\providecommand{\abs}{}\let\abs\relax\DeclarePairedDelimiter{\abs}{\lvert}{\rvert}
\providecommand{\norm}{}\let\norm\relax\DeclarePairedDelimiter{\norm}{\lVert}{\rVert}
\usepackage[table]{xcolor} % [table] : fournit aussi \rowcolors (alternance de lignes)
\usepackage{pagecolor}
\usepackage{tikz}
\usetikzlibrary{arrows.meta,positioning,calc,shapes.geometric,shapes.misc,shapes.arrows,decorations.pathmorphing,decorations.pathreplacing,decorations.markings,patterns,shadows,mindmap,trees,fit,backgrounds,matrix,intersections,angles,quotes}
\usepackage{pgfplots}
\usepackage[version=4]{mhchem} % équations nucléaires \ce{^{235}_{92}U}
\usepackage[european,siunitx]{circuitikz} % schémas électriques
\pgfplotsset{compat=1.18}
\usepgfplotslibrary{fillbetween,groupplots} % aires entre courbes (calcul intégral)
\usepackage{enumitem}
\usepackage{fancyhdr}
% [most] charge toutes les bibliothèques tcolorbox (listings, minted, raster,
% documentation...) et gonfle inutilement la mémoire TeX sur un document long
% (~300 boîtes) : on ne charge que ce qui est réellement utilisé.
\usepackage[skins,breakable]{tcolorbox}
\usepackage{graphicx}
\usepackage{colortbl}
\usepackage{booktabs}
\usepackage{tabularx}
\usepackage{diagbox} % case d'en-tête diagonale des tableaux à double entrée
% Colonnes extensibles centrée / gauche / droite (C, L, R), souvent utilisées par les modèles
\newcolumntype{C}{>{\centering\arraybackslash}X}
\newcolumntype{L}{>{\raggedright\arraybackslash}X}
\newcolumntype{R}{>{\raggedleft\arraybackslash}X}
\usepackage{array}
\usepackage{multicol}
\usepackage{siunitx}
% parse-numbers=false : accepte aussi \SI{2,5 \times 10^{-3}}{...}
\sisetup{locale=FR,output-decimal-marker={,},group-minimum-digits=4,parse-numbers=false}
\DeclareSIUnit\ohm{\ensuremath{\symup{\Omega}}} % U+2126 absent de la police
\DeclareSIUnit\division{div} % réglages d'oscilloscope (ms/div, V/div)
\DeclareSIUnit\tr{tr} % tours (vitesse de rotation, ex. \SI{1500}{\tr\per\minute})
\DeclareSIUnit\ch{ch} % cheval-vapeur (puissance moteur, unité usuelle hors SI)
\DeclareSIUnit\franccfa{FCFA} % franc CFA (coûts, contexte camerounais)
\DeclareSIUnit\dioptre{\ensuremath{\delta}} % vergence d'une lentille (dioptrie)
\usepackage{qrcode}
% \degree : commande fréquemment utilisée par les modèles pour un angle ou une
% température en dehors de \SI{}{} (non fournie par les paquets chargés).
\providecommand{\degree}{^{\circ}}
% \qed (fin de démonstration), utilisé par les modèles sans amsthm
\providecommand{\qed}{\hfill$\square$}
% \barre : double barre des valeurs interdites dans un tableau de variations (tkz-tab)
\providecommand{\barre}{\ensuremath{\|}}
% environnement proof (amsthm non chargé) utilisé par les modèles
\ifdefined\proof\else\newenvironment{proof}[1][Démonstration]{\par\noindent\textit{#1.}\ }{\qed\par}\fi
\usepackage{lastpage}
\usepackage{hyperref}
\hypersetup{hidelinks}

% ============================================================
% Palette « Pop » OnBuch+ — mêmes hex que la classe P (plus_ui.dart)
% ============================================================
\definecolor{poporange}{HTML}{F59321}
\definecolor{popdark}{HTML}{E07A0C}
\definecolor{popcream}{HTML}{FFF6E8}
\definecolor{popink}{HTML}{1C1714}
\definecolor{poppurple}{HTML}{7A5AE0}
\definecolor{popgreen}{HTML}{2E9E6A}
\definecolor{poppink}{HTML}{E0457B}
\definecolor{popblue}{HTML}{2F7DE1}
\definecolor{popgold}{HTML}{C98A12}
\definecolor{popmuted}{HTML}{8A7F76}
\definecolor{popline}{HTML}{EADFD2}
\definecolor{popgray}{HTML}{8A7F76}
\colorlet{popgrey}{popgray}
% Alias tolérants (le modèle nomme parfois les couleurs différemment)
\colorlet{popwhite}{white}
\colorlet{popblack}{popink}
\colorlet{popred}{poppink}
\colorlet{popyellow}{popgold}
% Teintes claires pour fonds de tableaux / zones
\colorlet{popblueL}{popblue!10!white}
\colorlet{poporangeL}{poporange!14!white}
\colorlet{popgreenL}{popgreen!12!white}
\colorlet{poppurpleL}{poppurple!12!white}
\colorlet{poppinkL}{poppink!12!white}
\colorlet{popgoldL}{popgold!14!white}

\pagecolor{popcream}

% ============================================================
% Typographie
% ============================================================
\defaultfontfeatures{Ligatures=TeX}
\setmainfont{PlusJakartaSans-Medium.ttf}[Path=../../../fonts/,BoldFont=PlusJakartaSans-Bold.ttf]
% Maths en sans-serif (Fira Math) avec chiffres et lettres droites de Plus
% Jakarta, pour que les formules se fondent dans le texte.
\usepackage[math-style=ISO,bold-style=ISO]{unicode-math}
\setmathfont{FiraMath-Regular.otf}[Path=../../../fonts/]
\setmathfont{PlusJakartaSans-Medium.ttf}[Path=../../../fonts/,range={up/num,up/{Latin,latin},\mathup}]
\usepackage{icomma} % virgule décimale sans espace en mode math
% Caractères mathématiques Unicode absents des polices de texte (Plus Jakarta,
% Archivo Black) : rendus par leur équivalent mathématique, en texte comme en
% formule — sinon ils s'affichent en carré vide (ex. « Arithmétique dans ℤ »).
\usepackage{newunicodechar}
\newunicodechar{ℝ}{\ensuremath{\mathbb{R}}}
\newunicodechar{ℤ}{\ensuremath{\mathbb{Z}}}
\newunicodechar{ℂ}{\ensuremath{\mathbb{C}}}
\newunicodechar{ℕ}{\ensuremath{\mathbb{N}}}
\newunicodechar{ℚ}{\ensuremath{\mathbb{Q}}}
\newunicodechar{𝔻}{\ensuremath{\mathbb{D}}}
\newunicodechar{α}{\ensuremath{\alpha}}
\newunicodechar{β}{\ensuremath{\beta}}
\newunicodechar{γ}{\ensuremath{\gamma}}
\newunicodechar{δ}{\ensuremath{\delta}}
\newunicodechar{ε}{\ensuremath{\varepsilon}}
\newunicodechar{θ}{\ensuremath{\theta}}
\newunicodechar{λ}{\ensuremath{\lambda}}
\newunicodechar{μ}{\ensuremath{\mu}}
\newunicodechar{σ}{\ensuremath{\sigma}}
\newunicodechar{τ}{\ensuremath{\tau}}
\newunicodechar{φ}{\ensuremath{\varphi}}
\newunicodechar{ψ}{\ensuremath{\psi}}
\newunicodechar{ω}{\ensuremath{\omega}}
\newunicodechar{Σ}{\ensuremath{\Sigma}}
% majuscules produites par \MakeUppercase dans les titres (α-aminés -> Α)
\newunicodechar{Α}{\ensuremath{\alpha}}
\newunicodechar{Β}{\ensuremath{\beta}}
\newunicodechar{Γ}{\ensuremath{\Gamma}}
\newunicodechar{Δ}{\ensuremath{\Delta}}
\newunicodechar{Ε}{\ensuremath{\varepsilon}}
\newunicodechar{Θ}{\ensuremath{\Theta}}
\newunicodechar{Λ}{\ensuremath{\Lambda}}
\newunicodechar{Μ}{\ensuremath{\mu}}
\newunicodechar{Τ}{\ensuremath{\tau}}
\newunicodechar{Φ}{\ensuremath{\Phi}}
\newunicodechar{Ψ}{\ensuremath{\Psi}}
\newunicodechar{Ω}{\ensuremath{\Omega}}
\newunicodechar{↦}{\ensuremath{\mapsto}}
\newunicodechar{∈}{\ensuremath{\in}}
\newunicodechar{∉}{\ensuremath{\notin}}
\newunicodechar{∧}{\ensuremath{\wedge}}
\newunicodechar{⇔}{\ensuremath{\Leftrightarrow}}
\newunicodechar{⟺}{\ensuremath{\Longleftrightarrow}}
\newunicodechar{⇒}{\ensuremath{\Rightarrow}}
\newunicodechar{⟹}{\ensuremath{\Longrightarrow}}
\newunicodechar{∘}{\ensuremath{\circ}}
\newunicodechar{∪}{\ensuremath{\cup}}
\newunicodechar{∩}{\ensuremath{\cap}}
\newunicodechar{≡}{\ensuremath{\equiv}}
\newunicodechar{⊂}{\ensuremath{\subset}}
\newunicodechar{⊥}{\ensuremath{\perp}}
\newunicodechar{∃}{\ensuremath{\exists}}
\newunicodechar{∀}{\ensuremath{\forall}}
\newunicodechar{∛}{\ensuremath{\sqrt[3]{\,}}}
\newunicodechar{∜}{\ensuremath{\sqrt[4]{\,}}}
\newunicodechar{⃗}{\ensuremath{{}^{\to}}}
\newunicodechar{ⁿ}{\ifmmode^{n}\else\textsuperscript{\ensuremath{n}}\fi}
\newunicodechar{ˣ}{\ifmmode^{x}\else\textsuperscript{\ensuremath{x}}\fi}
\newunicodechar{ᵇ}{\ifmmode^{b}\else\textsuperscript{\ensuremath{b}}\fi}
\newunicodechar{ʳ}{\ifmmode^{r}\else\textsuperscript{\ensuremath{r}}\fi}
\newunicodechar{ᵘ}{\ifmmode^{u}\else\textsuperscript{\ensuremath{u}}\fi}
\newunicodechar{ʸ}{\ifmmode^{y}\else\textsuperscript{\ensuremath{y}}\fi}
\newunicodechar{ᵃ}{\ifmmode^{a}\else\textsuperscript{\ensuremath{a}}\fi}
\newunicodechar{ᵉ}{\ifmmode^{e}\else\textsuperscript{\ensuremath{e}}\fi}
\newunicodechar{ᵏ}{\ifmmode^{k}\else\textsuperscript{\ensuremath{k}}\fi}
\newunicodechar{ᵖ}{\ifmmode^{p}\else\textsuperscript{\ensuremath{p}}\fi}
\newunicodechar{ᴺ}{\ifmmode^{N}\else\textsuperscript{\ensuremath{N}}\fi}
\newunicodechar{ᶜ}{\ifmmode^{c}\else\textsuperscript{\ensuremath{c}}\fi}
\newunicodechar{ʰ}{\ifmmode^{h}\else\textsuperscript{\ensuremath{h}}\fi}
\newunicodechar{ᵠ}{\ifmmode^{q}\else\textsuperscript{\ensuremath{q}}\fi}
\newunicodechar{ₙ}{\ifmmode_{n}\else\textsubscript{\ensuremath{n}}\fi}
\newunicodechar{ₐ}{\ifmmode_{a}\else\textsubscript{\ensuremath{a}}\fi}
\newunicodechar{ᵢ}{\ifmmode_{i}\else\textsubscript{\ensuremath{i}}\fi}
\newunicodechar{ᵣ}{\ifmmode_{r}\else\textsubscript{\ensuremath{r}}\fi}
\newunicodechar{ₖ}{\ifmmode_{k}\else\textsubscript{\ensuremath{k}}\fi}
\newunicodechar{ᵦ}{\ifmmode_{\beta}\else\textsubscript{\ensuremath{\beta}}\fi}
\newunicodechar{ₓ}{\ifmmode_{x}\else\textsubscript{\ensuremath{x}}\fi}
\newfontfamily\popdisplay{ArchivoBlack-Regular.ttf}[Path=../../../fonts/]
\newfontfamily\poplogo{SpaceGrotesk-Bold.ttf}[Path=../../../fonts/]
\newfontfamily\popsemi{PlusJakartaSans-SemiBold.ttf}[Path=../../../fonts/]
\newfontfamily\popxbold{PlusJakartaSans-ExtraBold.ttf}[Path=../../../fonts/]
\newfontfamily\popxboldls{PlusJakartaSans-ExtraBold.ttf}[Path=../../../fonts/,LetterSpace=3.0]

\setlist[enumerate]{leftmargin=1.7em,itemsep=5pt,topsep=4pt}
\setlist[itemize]{leftmargin=1.7em,itemsep=4pt,topsep=4pt,label={\color{poporange}\textbullet}}
\setlength{\parindent}{0pt}
\setlength{\parskip}{5pt}
\renewcommand{\arraystretch}{1.25}

% pgfplots : style Pop par défaut
\pgfplotsset{
  every axis/.append style={axis line style={popink,line width=1pt},tick style={popink},
    grid style={popline,line width=0.5pt},label style={font=\small},tick label style={font=\footnotesize}},
  popcurve/.style={poporange,line width=1.8pt,no markers,smooth},
}
% Même style disponible dans un \draw TikZ simple (hors pgfplots)
\tikzset{popcurve/.style={poporange,line width=1.8pt}}
\tikzset{popgrid/.style={popline,very thin}}
\colorlet{popcurve}{poporange} % le nom du style est parfois utilisé comme couleur

% ============================================================
% Pill / squiggle
% ============================================================
\newcommand{\poppill}[3]{%
  \tikz[baseline=-0.55ex]\node[draw=popink,line width=1.2pt,rounded corners=2.6pt,%
    fill=#2,inner xsep=6.5pt,inner ysep=3pt]{%
    \popxboldls\fontsize{7}{7}\selectfont\color{#3}\MakeUppercase{#1}};%
}
\newcommand{\popsquiggle}[1]{%
  \tikz[baseline=-0.4ex]{
    \draw[#1,line width=2.6pt,line cap=round]
      plot [smooth] coordinates {(0,0.09) (0.34,0.01) (0.68,0.21) (1.02,0.01) (1.36,0.10)};
  }%
}
% Mot-clé surligné (marqueur orange sous le texte)
\newcommand{\cle}[1]{{\popxbold\color{popdark}#1}}

% ============================================================
% En-tête / pied
% ============================================================
\pagestyle{fancy}
\fancyhf{}
\renewcommand{\headrulewidth}{0pt}
\renewcommand{\footrulewidth}{0pt}
\lhead{\raisebox{-0.15ex}{\poplogo\fontsize{13}{13}\selectfont\color{popink}OnBuch\color{poporange}+}}
\rhead{\poppill{\DOCMATIERE\ \textbullet\ \DOCNIVEAU\ \textbullet\ Leçon \DOCLECON}{white}{popink}}
\lfoot{\popxbold\fontsize{7.6}{7.6}\selectfont\color{popmuted}\MakeUppercase{Cours signé OnBuch+}\ \textbullet\ {\popsemi\DOCTITRE}}
\rfoot{\tikz[baseline=-0.6ex]\node[rounded corners=8.5pt,fill=popink,minimum height=17pt,minimum width=17pt,inner xsep=5pt,inner sep=0]{\popxbold\fontsize{8}{8}\selectfont\color{popcream}\thepage\,/\,\pageref*{LastPage}};}

% ============================================================
% Titres
% ============================================================
\newcommand{\chaptitre}[2]{%
  \begin{tcolorbox}[enhanced,colback=poporange,colframe=popink,boxrule=2.6pt,arc=11pt,
      drop shadow=popink,left=15pt,right=15pt,top=12pt,bottom=11pt,before skip=3pt,after skip=18pt]
    \poppill{#2}{popink}{popcream}\par\vspace{9pt}
    {\raggedright\hyphenpenalty=10000\exhyphenpenalty=10000\popdisplay\fontsize{22}{24}\selectfont\color{popcream}\MakeUppercase{#1}\par}
  \end{tcolorbox}
}
% Section numérotée : pastille orange avec le numéro + titre Archivo
\newcounter{popsec}
\newcommand{\coursec}[1]{%
  \stepcounter{popsec}\setcounter{popsub}{0}%
  \par\vspace{16pt}\noindent
  \begin{minipage}{\linewidth}
  \tikz[baseline=-0.7ex]\node[draw=popink,line width=1.6pt,fill=poporange,rounded corners=4pt,minimum size=22pt,inner sep=0]{\popdisplay\fontsize{12}{12}\selectfont\color{popcream}\thepopsec};%
  \hspace{8pt}{\popdisplay\fontsize{15}{17}\selectfont\color{popink}\MakeUppercase{#1}}\par
  \vspace{4pt}\noindent{\color{poporange}\rule{36pt}{3.6pt}}
  \end{minipage}\par\nopagebreak\vspace{8pt}
}
\newcounter{popsub}
\newcommand{\courssub}[1]{%
  \stepcounter{popsub}%
  \par\vspace{9pt}\noindent{\popxbold\fontsize{11.5}{13}\selectfont\color{popink}{\color{poporange}\thepopsec.\thepopsub}\hspace{6pt}#1}\par\nopagebreak\vspace{4pt}
}

% ============================================================
% Boîtes pédagogiques — même recette : fond blanc, bordure épaisse colorée,
% ombre dure encre, badge plein. Titre optionnel [#1].
% ============================================================
\tcbset{popbase/.style={breakable,enhanced,colback=white,boxrule=2.3pt,arc=9pt,
  shadow={3pt}{-3pt}{0pt}{popink},left=14pt,right=14pt,top=11pt,bottom=11pt,before skip=11pt,after skip=15pt,
  fontupper=\popsemi\fontsize{10.6}{14.5}\selectfont\color{popink},
  fontlower=\popsemi\fontsize{10.6}{14.5}\selectfont\color{popink}}}
% Coche dessinée (le glyphe ✓ n'existe pas dans les polices du document)
\newcommand{\popcheck}[1]{\tikz[baseline=-0.5ex]\draw[#1,line width=1.6pt,line cap=round,line join=round](-0.13,0.0)--(-0.04,-0.09)--(0.13,0.1);}
\providecommand{\coche}{\popcheck{popgreen}}
\providecommand{\resultat}[1]{\par\smallskip\noindent\fcolorbox{popgreen}{popgreenL}{\parbox{\dimexpr\linewidth-2\fboxsep-2\fboxrule\relax}{\textbf{Résultat.} #1}}\par\smallskip}
\newcommand{\popboxhead}[3]{\poppill{#1}{#2}{white}\ifx\relax#3\relax\else\hspace{6pt}{\popxbold\fontsize{10.6}{12}\selectfont\color{popink}#3}\fi\par\vspace{7pt}}

\newtcolorbox{aretenir}[1][]{popbase,colframe=popgreen,before upper={\popboxhead{À retenir}{popgreen}{#1}}}
\newtcolorbox{exemplebox}[1][]{popbase,colframe=poporange,before upper={\popboxhead{Exemple}{poporange}{#1}}}
\newtcolorbox{definition}[1][]{popbase,colframe=popblue,before upper={\popboxhead{Définition}{popblue}{#1}}}
\newtcolorbox{propriete}[1][]{popbase,colframe=poppurple,before upper={\popboxhead{Propriété}{poppurple}{#1}}}
\newtcolorbox{methode}[1][]{popbase,colframe=popgold,before upper={\popboxhead{Méthode}{popgold}{#1}}}
\newtcolorbox{attention}[1][]{popbase,colframe=poppink,before upper={\popboxhead{Attention}{poppink}{#1}}}
\newtcolorbox{experience}[1][]{popbase,colframe=popblue,colback=popblueL,before upper={\popboxhead{Expérience}{popblue}{#1}}}
% « Le savais-tu ? » : carte violette pleine (contexte, culture, Cameroun)
\newtcolorbox{savaistu}[1][]{popbase,colback=poppurple,colframe=popink,
  fontupper=\popsemi\fontsize{10.4}{14.2}\selectfont\color{white},
  before upper={\poppill{Le savais-tu\,?}{white}{poppurple}\ifx\relax#1\relax\else\hspace{6pt}{\popxbold\color{white}#1}\fi\par\vspace{7pt}}}
% Exercice résolu : énoncé en haut, \tcblower puis la solution sur fond crème
\newcounter{popexo}\newcounter{popexr}
\newtcolorbox{exoresolu}[1][]{popbase,colframe=poporange,colbacklower=popcream,
  segmentation style={popink,line width=1.2pt,dashed},
  before upper={\stepcounter{popexr}\popboxhead{Exercice résolu \thepopexr}{poporange}{#1}},
  before lower={\poppill{Solution}{popink}{popcream}\par\vspace{6pt}}}
% Exercice (non résolu en place) — la correction est dans la partie Corrigés
\newtcolorbox{exercice}[1][]{popbase,colframe=popink,
  before upper={\stepcounter{popexo}\popboxhead{Exercice \thepopexo}{popink}{#1}}}
% Corrigé
\newtcolorbox{corrige}[1][]{popbase,colframe=popgreen,colback=popgreenL,
  before upper={\popboxhead{Corrigé}{popgreen}{#1}}}
% Objectifs / prérequis / bilan
\newtcolorbox{objectifs}{popbase,colframe=popink,colback=poporangeL,
  before upper={\popboxhead{Objectifs de la leçon}{popink}{}}}
\newtcolorbox{prerequis}{popbase,colframe=popink,colback=popblueL,
  before upper={\popboxhead{Prérequis}{popblue}{}}}
\newtcolorbox{bilan}{popbase,colframe=popink,colback=white,boxrule=2.8pt,
  before upper={\popboxhead{Fiche bilan}{poporange}{L'essentiel à connaître pour l'examen}}}
% Figure encadrée avec légende
\newtcolorbox{popfigure}[1][]{enhanced,breakable=false,colback=white,colframe=popline,boxrule=1.4pt,arc=8pt,
  left=10pt,right=10pt,top=10pt,bottom=8pt,before skip=10pt,after skip=12pt,halign=center,
  lower separated=false,halign lower=center,
  fontlower=\popsemi\fontsize{9}{11.5}\selectfont\color{popmuted},
  #1}
\newcommand{\legende}[1]{\par\vspace{6pt}{\popsemi\fontsize{9}{11.5}\selectfont\color{popmuted}#1}}

% ============================================================
% Signature OnBuch+
% ============================================================
\newcommand{\poplogoinline}[1]{{\poplogo\fontsize{#1}{#1}\selectfont\color{popink}OnBuch\color{poporange}+}}
\newcommand{\signatureonbuch}{%
  \begin{tcolorbox}[enhanced,colback=popink,colframe=popink,boxrule=2.2pt,arc=10pt,
      drop shadow=poporange,left=14pt,right=14pt,top=10pt,bottom=10pt,before skip=0pt,after skip=0pt]
    \tikz[baseline=-0.7ex]\node[circle,fill=popgreen,draw=popcream,line width=1.3pt,minimum size=22pt,inner sep=0]{\popcheck{white}};%
    \hspace{9pt}\begin{minipage}[c]{0.62\linewidth}
      {\popxbold\fontsize{12}{13}\selectfont\color{popcream}Signé\ \textbullet\ {\poplogo OnBuch\color{poporange}+}}\par\vspace{2pt}
      {\raggedright\popsemi\fontsize{8}{10.5}\selectfont\color{popline}\MakeUppercase{Équipe pédagogique OnBuch+}\\\MakeUppercase{Conforme au programme officiel MINESEC}\par}
    \end{minipage}\hfill
    {\poplogo\fontsize{22}{22}\selectfont\color{popcream}OnBuch\color{poporange}+}
  \end{tcolorbox}%
}

% ============================================================
% Page de garde
% #1 titre · #2 matière · #3 niveau · #4 module · #5 n° leçon · #6 durée ·
% #7 description / objectifs (texte ou itemize)
% ============================================================
\newcommand{\pagegarde}[7]{%
  \thispagestyle{empty}
  \newgeometry{a4paper,margin=1.7cm,top=1.6cm,bottom=1.4cm}%
  \begin{tikzpicture}[remember picture,overlay]
    \fill[poporange!18] ([xshift=-3.2cm,yshift=-3.6cm]current page.north east) circle (4.6cm);
    \fill[poppurple!14] ([xshift=2.4cm,yshift=7.6cm]current page.south west) circle (3.8cm);
  \end{tikzpicture}%
  {\poplogo\fontsize{30}{30}\selectfont\color{popink}OnBuch\color{poporange}+}\hfill
  \raisebox{0.6ex}{\poppill{Cours signé \textbullet\ Premium}{popink}{popcream}}\par
  \vspace{1pt}\popsquiggle{poppurple}\par
  \vspace{8pt}
  \begin{tcolorbox}[enhanced,colback=poporange,colframe=popink,boxrule=2.8pt,arc=13pt,
      drop shadow=popink,left=18pt,right=18pt,top=12pt,bottom=13pt,after skip=10pt]
    \poppill{#2 \textbullet\ #3}{popink}{popcream}\hspace{5pt}\poppill{#4}{white}{popink}\par\vspace{8pt}
    {\popdisplay\fontsize{14}{14}\selectfont\color{popink}LEÇON #5}\par\vspace{4pt}
    {\raggedright\hyphenpenalty=10000\exhyphenpenalty=10000\popdisplay\fontsize{25}{27}\selectfont\color{popcream}\MakeUppercase{#1}\par}
    \vspace{8pt}
    {\popxbold\fontsize{9.5}{9.5}\selectfont\color{popink}\MakeUppercase{Durée conseillée : #6}}
  \end{tcolorbox}
  \begin{tcolorbox}[enhanced,colback=white,colframe=popink,boxrule=2.2pt,arc=10pt,
      drop shadow=popink,left=16pt,right=16pt,top=10pt,bottom=10pt,after skip=8pt]
    \poppill{Dans cette leçon}{poppurple}{white}\par\vspace{5pt}
    \popsemi\fontsize{9.3}{12.6}\selectfont\color{popink}#7
  \end{tcolorbox}
  \noindent\begin{minipage}[t]{0.487\linewidth}
    \begin{tcolorbox}[enhanced,colback=popgreen,colframe=popink,boxrule=2.2pt,arc=10pt,
        drop shadow=popink,left=11pt,right=11pt,top=10pt,bottom=10pt,height=3.1cm,valign=center]
      \tcbox[colback=white,colframe=popink,boxrule=1.3pt,arc=5pt,boxsep=0pt,nobeforeafter,
        left=3pt,right=3pt,top=3pt,bottom=3pt,valign=center]{\qrcode[height=1.9cm]{https://play.google.com/store/apps/details?id=com.luvvix.onbuch&hl=fr}}%
      \hspace{8pt}\begin{minipage}[c]{0.5\linewidth}
        {\popdisplay\fontsize{11}{12.5}\selectfont\color{white}\MakeUppercase{Télécharge OnBuch}}\par\vspace{4pt}
        {\raggedright\popsemi\fontsize{8.4}{11}\selectfont\color{white}Gratuit \textbullet\ Google Play\\Tuteur IA, annales, concours, résultats\par}
      \end{minipage}
    \end{tcolorbox}
  \end{minipage}\hfill
  \begin{minipage}[t]{0.487\linewidth}
    \begin{tcolorbox}[enhanced,colback=poppurple,colframe=popink,boxrule=2.2pt,arc=10pt,
        drop shadow=popink,left=11pt,right=11pt,top=10pt,bottom=10pt,height=3.1cm,valign=center]
      \tikz[baseline=-0.7ex]\node[circle,fill=white,draw=popink,line width=1.3pt,minimum size=1.6cm,inner sep=0]{\popdisplay\fontsize{24}{24}\selectfont\color{poppurple}+};%
      \hspace{8pt}\begin{minipage}[c]{0.6\linewidth}
        {\popdisplay\fontsize{11}{12.5}\selectfont\color{white}\MakeUppercase{Passe à OnBuch+}}\par\vspace{4pt}
        {\raggedright\popsemi\fontsize{8.4}{11}\selectfont\color{white}Tous les cours signés, Tuteur IA illimité, fiches PDF — onglet OnBuch+ de l'app\par}
      \end{minipage}
    \end{tcolorbox}
  \end{minipage}
  \vspace{8pt}
  \popcontenu
  \vfill
  \signatureonbuch
  \restoregeometry
  \newpage
}

% Rangée « Ce cours contient » (page de garde)
\newcommand{\poptile}[3]{% #1 couleur #2 gros texte #3 légende
  \begin{tcolorbox}[enhanced,colback=#1,colframe=popink,boxrule=2pt,arc=9pt,shadow={2.5pt}{-2.5pt}{0pt}{popink},
      left=6pt,right=6pt,top=9pt,bottom=9pt,halign=center,valign=center,height=1.75cm,width=\linewidth,nobeforeafter]
    {\popdisplay\fontsize{12.5}{13}\selectfont\color{white}\MakeUppercase{#2}}\par\vspace{3pt}
    {\centering\popsemi\fontsize{7.8}{9.5}\selectfont\color{white}#3\par}
  \end{tcolorbox}}
\newcommand{\popcontenu}{%
  \par\noindent{\popxboldls\fontsize{7.5}{7.5}\selectfont\color{popmuted}CE COURS SIGNÉ CONTIENT}\par\vspace{7pt}
  \noindent\begin{minipage}[t]{0.235\linewidth}\poptile{popblue}{Cours}{complet, illustré, pas à pas}\end{minipage}\hfill
  \begin{minipage}[t]{0.235\linewidth}\poptile{poporange}{Méthodes}{et exercices résolus}\end{minipage}\hfill
  \begin{minipage}[t]{0.235\linewidth}\poptile{poppink}{Exercices}{type examen + corrigés}\end{minipage}\hfill
  \begin{minipage}[t]{0.235\linewidth}\poptile{popgreen}{Fiche bilan}{l'essentiel à retenir}\end{minipage}\par}

% Sommaire
\newtcolorbox{sommaire}{enhanced,colback=white,colframe=popink,boxrule=2.2pt,arc=10pt,
  drop shadow=popink,left=18pt,right=18pt,top=13pt,bottom=13pt,before skip=6pt,after skip=20pt,
  before upper={\poppill{Sommaire}{poppurple}{white}\par\vspace{9pt}\popsemi\fontsize{10.6}{14.5}\selectfont\color{popink}}}

% Signature de fin de cours — poussée tout en bas de la dernière page
\newcommand{\findecours}{%
  \par\vfill\nopagebreak
  \begin{center}\popsquiggle{poporange}\end{center}\vspace{4pt}
  \signatureonbuch
}
\AtBeginDocument{\def\llbracket{\mathopen{[\mkern-3.5mu[}}\def\rrbracket{\mathclose{]\mkern-3.5mu]}}} % intervalles d'entiers (glyphe ⟦ absent de la police)
% Glyphes absents de Fira Math : redéfinis à partir de glyphes disponibles
\AtBeginDocument{%
  \def\setminus{\mathbin{\backslash}}\let\smallsetminus\setminus
  \def\vdots{\mathord{\vbox{\baselineskip3.2pt\lineskiplimit0pt\kern4pt\hbox{.}\hbox{.}\hbox{.}}}}%
  \def\ddots{\mathinner{\mkern1mu\raise6.5pt\hbox{.}\mkern2mu\raise3.6pt\hbox{.}\mkern2mu\raise0.7pt\hbox{.}\mkern1mu}}%
  \def\triangle{\mathord{\scalebox{0.9}{$\symup{\Delta}$}}}%
  \def\nabla{\mathord{\raisebox{\depth}{\scalebox{1}[-1]{$\symup{\Delta}$}}}}%
  \def\checkmark{\popcheck{popdark}}%
  \def\star{\mathbin{\ast}}\def\bigstar{\mathord{\ast}}%
  \def\diamond{\mathbin{\scalebox{0.6}{$\Diamond$}}}%
  \def\bigcup{\mathop{\mathchoice{\raisebox{-0.3em}{\scalebox{1.7}{$\cup$}}}{\scalebox{1.25}{$\cup$}}{\cup}{\cup}}\displaylimits}%
  \def\bigcap{\mathop{\mathchoice{\raisebox{-0.3em}{\scalebox{1.7}{$\cap$}}}{\scalebox{1.25}{$\cap$}}{\cap}{\cap}}\displaylimits}%
  \def\lesseqqgtr{\mathrel{\lessgtr}}\def\bigoplus{\mathop{\oplus}\displaylimits}%
}
\providecommand{\ssi}{si et seulement si} % abréviation fréquente
\AtBeginDocument{\providecommand{\wideparen}{\overparen}} % arc de cercle

%% FIGFIX-BEGIN v5 — correctifs globaux des figures (docs/onbuch-generation/tools/figfix)
\usepackage{adjustbox}\usepackage{etoolbox}\tcbuselibrary{raster}
% toute figure TikZ/pgfplots plus large que la ligne est réduite à la largeur disponible
\BeforeBeginEnvironment{tikzpicture}{\begin{adjustbox}{max width=\linewidth}}
\AfterEndEnvironment{tikzpicture}{\end{adjustbox}}
% légendes pgfplots lisibles par-dessus les courbes
\pgfplotsset{every axis/.append style={legend style={fill=white,fill opacity=0.92,text opacity=1,draw=black!15,inner sep=3pt}}}
% étiquettes d'arêtes (syntaxe edge["..."]) avec fond blanc
\tikzset{every edge quotes/.append style={fill=white,inner sep=1.2pt}}
%% FIGFIX-END
\begin{document}

\pagegarde{Interaction lumière-matière : le photon}{Physique}{1ère TI}{Module 3 : Optique}{N13}{7 h}{Cette leçon introduit le modèle corpusculaire de la lumière et la quantification de l'énergie pour expliquer comment la matière émet, absorbe et diffuse la lumière. Elle fournit les outils (E = hν, λ = c/ν, diagrammes de niveaux d'énergie) pour interpréter les spectres et déterminer la température d'un corps chaud par analyse spectrale.
\vspace{4pt}
\begin{multicols}{2}\small
\begin{itemize}
\item À la fin de cette leçon, je sais décrire les trois modes d'interaction lumière-matière : émission, absorption et diffusion.
\item À la fin de cette leçon, je sais expliquer pourquoi les niveaux d'énergie d'un atome sont quantifiés et représenter un diagramme de niveaux d'énergie.
\item À la fin de cette leçon, je sais définir le photon et citer ses caractéristiques (masse nulle, charge nulle, vitesse c).
\item À la fin de cette leçon, je sais utiliser les relations E = hν et λ = c/ν pour calculer l'énergie d'un photon en joules et en électronvolts.
\item À la fin de cette leçon, je sais exploiter un diagramme de niveaux d'énergie pour interpréter l'émission ou l'absorption d'un photon (ΔE = hν).
\item À la fin de cette leçon, je sais décrire le profil spectral d'une source lumineuse et distinguer spectre continu et spectre de raies.
\end{itemize}\end{multicols}}

\begin{sommaire}
\begin{multicols}{2}
\begin{enumerate}
\item Interaction lumière-matière : émission, absorption, diffusion
\item Quantification des niveaux d'énergie de la matière
\item Le photon et son énergie : E = hν
\item Échanges d'énergie lumière-matière : ΔE = hν
\item Profil spectral d'une source lumineuse
\item Le spectre solaire : forme et raies
\item Activité d'intégration
\item Méthodes et exercices résolus
\item Exercices
\item Corrigés des exercices
\item Fiche bilan
\end{enumerate}
\end{multicols}
\end{sommaire}

\begin{objectifs}
\begin{itemize}
\item À la fin de cette leçon, je sais décrire les trois modes d'interaction lumière-matière : émission, absorption et diffusion.
\item À la fin de cette leçon, je sais expliquer pourquoi les niveaux d'énergie d'un atome sont quantifiés et représenter un diagramme de niveaux d'énergie.
\item À la fin de cette leçon, je sais définir le photon et citer ses caractéristiques (masse nulle, charge nulle, vitesse c).
\item À la fin de cette leçon, je sais utiliser les relations E = hν et λ = c/ν pour calculer l'énergie d'un photon en joules et en électronvolts.
\item À la fin de cette leçon, je sais exploiter un diagramme de niveaux d'énergie pour interpréter l'émission ou l'absorption d'un photon (ΔE = hν).
\item À la fin de cette leçon, je sais décrire le profil spectral d'une source lumineuse et distinguer spectre continu et spectre de raies.
\item À la fin de cette leçon, je sais expliquer les caractéristiques du spectre solaire (fond continu, raies d'absorption de Fraunhofer).
\item À la fin de cette leçon, je sais déterminer la température d'un corps chaud à partir de la longueur d'onde du maximum d'émission (loi de Wien).
\end{itemize}
\end{objectifs}

\begin{prerequis}
\begin{itemize}
\item Modèle ondulatoire de la lumière : longueur d'onde, fréquence, célérité c = 3,00 × 10\^{}8 m/s (classe de Seconde / début de Première)
\item Spectre de la lumière blanche et dispersion par un prisme (domaines du visible : 400 nm à 800 nm)
\item Notion d'énergie et conversions d'unités (joule, préfixes nano, micro)
\item Écriture scientifique et calculs avec les puissances de 10
\item Notion de spectre d'émission observé au collège (lampe à incandescence, tube spectral)
\end{itemize}
\end{prerequis}

\newpage

\coursec{Interaction lumière-matière : émission, absorption, diffusion}

La nuit, à Yaoundé, les lampadaires au sodium baignent les rues d'une lumière jaune orangée, tandis que les lampes LED des boutiques émettent une lumière blanche éclatante. Pourquoi ces couleurs si différentes ? Au marché Mokolo, le forgeron juge la température de son fer simplement en observant sa couleur : rouge sombre, puis orange, puis blanc. Et comment les astronomes du monde entier déterminent-ils la température et la composition du Soleil sans jamais y toucher ? Toutes ces questions trouvent leur réponse dans un grain de lumière extraordinairement petit : le \cle{photon}.

Avant de découvrir ce grain de lumière, il faut comprendre \emph{ce qui se passe} lorsque la lumière rencontre la matière. C'est l'objet de cette première section : décrire les trois grandes interactions possibles entre la lumière et la matière — l'\cle{émission}, l'\cle{absorption} et la \cle{diffusion} — et montrer qu'elles obéissent toutes à un même principe fondamental : la conservation de l'énergie.

\courssub{Les deux visages de la lumière : onde ou grain ?}

\textbf{Ce que le modèle ondulatoire explique très bien.}

Depuis le chapitre sur les ondes, tu sais que la lumière se comporte comme une \cle{onde électromagnétique} caractérisée par sa fréquence $\nu$ (en hertz) et sa longueur d'onde dans le vide $\lambda$ (en mètres), reliées par la relation fondamentale :
\[ \lambda = \frac{c}{\nu} \]
où $c = \SI{3,00e8}{m.s^{-1}}$ est la célérité de la lumière dans le vide. Vérifions l'homogénéité : $[\lambda] = \dfrac{[c]}{[\nu]} = \dfrac{\text{m}\cdot\text{s}^{-1}}{\text{s}^{-1}} = \text{m}$. La relation est bien cohérente.

Ce modèle ondulatoire rend parfaitement compte de la \emph{propagation} de la lumière : la diffraction (étalement d'un faisceau traversant une ouverture étroite), les interférences (franges lumineuses obtenues avec deux fentes fines), la réflexion et la réfraction. Autant de phénomènes que tu as étudiés en classe de Seconde et en Première.

\textbf{Ce que le modèle ondulatoire n'explique pas.}

Mais dès qu'il s'agit des \emph{échanges d'énergie} entre la lumière et la matière, le modèle ondulatoire échoue. Trois faits expérimentaux historiques l'ont montré :

\begin{itemize}
\item la couleur de la lumière émise par un corps chauffé ne dépend que de sa température, jamais de sa composition chimique — le modèle ondulatoire classique prédisait au contraire une catastrophe (le corps chaud devrait émettre une énergie infinie dans l'ultraviolet !) ;
\item l'effet photoélectrique : une lumière même très intense ne peut arracher des électrons à un métal si sa fréquence est trop faible, alors qu'une lumière faible mais de fréquence suffisante y parvient instantanément ;
\item les gaz émettent et absorbent la lumière uniquement à des fréquences bien précises (raies spectrales), ce qui est incompréhensible avec une onde continue.
\end{itemize}

Pour résoudre ces énigmes, Planck (1900) puis Einstein (1905) ont proposé une idée révolutionnaire : la lumière échange son énergie avec la matière \emph{par paquets}, par grains, appelés \cle{photons}. C'est le \cle{modèle corpusculaire} de la lumière, que nous étudierons en détail dans la section 3. Retenons dès maintenant l'idée directrice :

\begin{aretenir}[Deux modèles complémentaires]
La lumière possède une \emph{double nature} :
\begin{itemize}
\item pour décrire sa \textbf{propagation} (diffraction, interférences), on utilise le \textbf{modèle ondulatoire} (fréquence $\nu$, longueur d'onde $\lambda$) ;
\item pour décrire ses \textbf{échanges d'énergie} avec la matière (émission, absorption), on utilise le \textbf{modèle corpusculaire} (le photon, grain d'énergie $E = h\nu$).
\end{itemize}
Les deux modèles ne se contredisent pas : ils se complètent. La fréquence $\nu$ est le pont entre les deux.
\end{aretenir}

\courssub{L'émission : la matière libère de la lumière}

\textbf{Observation.} Un filament de lampe à incandescence porté à haute température brille. Un tube fluorescent de bureau s'allume lorsqu'on y établit une décharge électrique. Une LED de torche éclaire dès qu'un courant la traverse. Dans les trois cas, la matière a reçu de l'énergie (chaleur, décharge, courant électrique), puis elle la \emph{restitue sous forme de lumière}.

\begin{definition}[Émission de lumière]
L'\cle{émission} est le phénomène par lequel la matière, préalablement portée dans un \textbf{état excité} (état d'énergie élevée), revient dans un état de \textbf{moindre énergie} en libérant de la lumière.

L'énergie perdue par la matière est emportée par la lumière émise. On distingue selon le mode d'excitation :
\begin{itemize}
\item l'\emph{incandescence} : excitation par la chaleur (filament de lampe, fer du forgeron, filament d'une bougie) ;
\item la \emph{luminescence} : excitation par décharge électrique (tube fluorescent, enseignes au néon) ou par courant dans un semi-conducteur (LED).
\end{itemize}
\end{definition}

\begin{exemplebox}[Trois sources de lumière, trois couleurs]
\begin{itemize}
\item \textbf{Lampe à incandescence} : le filament de tungstène, chauffé vers $\SI{2700}{\celsius}$, émet un spectre continu dominé par le rouge-orangé : lumière « chaude ».
\item \textbf{Lampadaire au sodium} (rues de Yaoundé) : la décharge électrique dans la vapeur de sodium excite les atomes, qui retombent en émettant une lumière quasi monochromatique jaune-orangée ($\lambda \approx \SI{589}{nm}$).
\item \textbf{LED blanche} : une LED bleue excite un luminophore qui réémet un spectre large ; la superposition donne une lumière blanche éclatante, avec très peu de chaleur perdue.
\end{itemize}
La couleur de la lumière émise dépend donc directement des \emph{niveaux d'énergie} de la matière émettrice — c'est tout l'objet de la section 2.
\end{exemplebox}

\courssub{L'absorption : la matière capte la lumière}

\textbf{Observation.} Laisse une feuille de papier noir et une feuille blanche au soleil de midi à Douala : la noire devient brûlante, la blanche reste tiède. La feuille noire a \emph{capté} l'énergie lumineuse.

\begin{definition}[Absorption de lumière]
L'\cle{absorption} est le phénomène par lequel la matière \textbf{capture l'énergie} d'une lumière qui la traverse ou l'éclaire. L'énergie lumineuse absorbée est convertie en une autre forme d'énergie :
\begin{itemize}
\item énergie thermique (agitation des atomes : le corps s'échauffe) ;
\item énergie chimique (photosynthèse : la chlorophylle absorbe le rouge et le bleu pour fabriquer la matière végétale) ;
\item excitation des atomes (qui pourront ensuite réémettre).
\end{itemize}
Un corps qui absorbe toutes les radiations visibles paraît \textbf{noir}. Un filtre coloré absorbe certaines couleurs et laisse passer les autres : un filtre rouge absorbe le bleu et le vert, transmet le rouge.
\end{definition}

\begin{exemplebox}[La chlorophylle, championne de l'absorption]
Les feuilles des plantes nous paraissent vertes précisément parce que la chlorophylle \emph{n'absorbe pas} le vert : elle absorbe fortement le rouge ($\lambda \approx \SI{660}{nm}$) et le bleu ($\lambda \approx \SI{430}{nm}$) pour alimenter la photosynthèse, et renvoie le vert vers nos yeux. C'est grâce à cette absorption sélective que les plantations de cacao du Centre-Cameroun transforment la lumière du Soleil en énergie chimique.
\end{exemplebox}

\courssub{La diffusion : la lumière repart dans toutes les directions}

\textbf{Observation.} Dans une salle de classe poussiéreuse, le faisceau du vidéoprojecteur est visible de côté, alors même qu'on n'est pas dans son axe. Pourquoi ? Parce que les grains de poussière \emph{renvoient la lumière dans toutes les directions}.

\begin{definition}[Diffusion de la lumière]
La \cle{diffusion} est le phénomène par lequel la matière (particules, gouttelettes, molécules) \textbf{renvoie la lumière reçue dans toutes les directions de l'espace}, \textbf{sans changer sa fréquence} (donc sans changer sa couleur).

La lumière diffusée a la même fréquence $\nu$ que la lumière incidente : seule sa \emph{direction} change. L'efficacité de la diffusion dépend fortement de la taille des particules et de la longueur d'onde :
\begin{itemize}
\item particules très petites devant $\lambda$ (molécules de l'air) : diffusion d'autant plus intense que la longueur d'onde est courte — le bleu ($\lambda \approx \SI{450}{nm}$) est diffusé environ 5 fois plus que le rouge ($\lambda \approx \SI{700}{nm}$) ; c'est la \emph{diffusion Rayleigh} ;
\item particules de taille comparable ou supérieure à $\lambda$ (gouttelettes d'eau, poussières) : toutes les couleurs sont diffusées de façon comparable — la lumière blanche reste blanche.
\end{itemize}
\end{definition}

\begin{popfigure}
\begin{tikzpicture}[scale=0.95, every node/.style={font=\small}]
% Panneau 1 : émission
\begin{scope}[shift={(0,0)}]
\node[font=\bfseries\small, popblue] at (0,2.3) {Émission};
\shade[ball color=popblue!60] (0,0) circle (0.55);
\node[white] at (0,0) {atome};
\node[popdark, font=\scriptsize] at (0,-0.95) {état excité $\to$ état bas};
\draw[-{Stealth[length=3mm]}, poporange, very thick, decorate, decoration={snake, amplitude=1.2pt, segment length=5pt}] (0.6,0.25) -- (2.1,0.9);
\node[poporange, font=\scriptsize] at (1.6,1.25) {lumière};
\end{scope}
% Panneau 2 : absorption
\begin{scope}[shift={(5.2,0)}]
\node[font=\bfseries\small, popblue] at (0,2.3) {Absorption};
\shade[ball color=popblue!60] (0,0) circle (0.55);
\node[white] at (0,0) {atome};
\draw[-{Stealth[length=3mm]}, poporange, very thick, decorate, decoration={snake, amplitude=1.2pt, segment length=5pt}] (-2.1,0.9) -- (-0.6,0.25);
\node[poporange, font=\scriptsize] at (-1.6,1.25) {lumière};
\node[popdark, font=\scriptsize] at (0,-0.95) {l'atome gagne de l'énergie};
\end{scope}
% Panneau 3 : diffusion
\begin{scope}[shift={(10.4,0)}]
\node[font=\bfseries\small, popblue] at (0,2.3) {Diffusion};
\shade[ball color=poppurple!60] (0,0) circle (0.4);
\node[white, font=\scriptsize] at (0,0) {particule};
\draw[-{Stealth[length=3mm]}, poporange, very thick, decorate, decoration={snake, amplitude=1.2pt, segment length=5pt}] (-2.2,0) -- (-0.5,0);
\node[poporange, font=\scriptsize] at (-1.5,0.35) {lumière};
\foreach \a in {20,60,110,160,-40,-90} {
  \draw[-{Stealth[length=2.2mm]}, poppink, thick, decorate, decoration={snake, amplitude=0.9pt, segment length=4.5pt}] (\a:0.45) -- (\a:1.5);
}
\node[popdark, font=\scriptsize] at (0,-1.7) {même fréquence, toutes directions};
\end{scope}
\end{tikzpicture}
\legende{Les trois interactions lumière-matière : émission (l'atome excité libère de la lumière), absorption (l'atome capture la lumière), diffusion (la particule renvoie la lumière dans toutes les directions sans changer sa fréquence).}
\end{popfigure}

\begin{attention}[Diffusion $\neq$ réflexion]
Ne confonds pas \textbf{diffusion} et \textbf{réflexion} :
\begin{itemize}
\item la \textbf{réflexion} (miroir, surface polie) renvoie la lumière dans \emph{une seule direction}, donnée par les lois de Snell-Descartes ($i' = i$) : l'image est nette ;
\item la \textbf{diffusion} (mur blanc, nuage, poussière) renvoie la lumière dans \emph{toutes les directions} : aucune image nette ne se forme, mais l'objet devient visible depuis n'importe quel point de vue.
\end{itemize}
C'est grâce à la diffusion que tu vois cette page de cours sous tous les angles, alors qu'un miroir ne te renvoie une image que dans une direction précise. Autre piège : dans la diffusion étudiée ici, la fréquence (la couleur) de la lumière \textbf{ne change pas}.
\end{attention}

\begin{popfigure}
\begin{tikzpicture}[scale=1]
% Lampe torche
\draw[fill=popdark!80] (0,-0.35) rectangle (1.1,0.35);
\draw[fill=popdark!60] (1.1,-0.5) -- (1.6,-0.35) -- (1.6,0.35) -- (1.1,0.5) -- cycle;
\node[font=\scriptsize, popdark] at (0.55,-0.75) {lampe torche};
% Faisceau
\fill[popgold!25] (1.6,-0.35) -- (7.5,-1.1) -- (7.5,1.1) -- (1.6,0.35) -- cycle;
\draw[popgold, thick] (1.6,-0.35) -- (7.5,-1.1);
\draw[popgold, thick] (1.6,0.35) -- (7.5,1.1);
% Poussières
\foreach \p in {(2.5,0.2),(3.2,-0.3),(3.8,0.5),(4.4,-0.1),(5.0,0.35),(5.6,-0.5),(6.2,0.1),(6.8,0.55),(4.1,-0.6),(5.3,0.7)} {
  \fill[poppurple!80] \p circle (0.05);
}
% Rayons diffusés depuis deux grains
\foreach \a in {40,120,200,300} {
  \draw[-{Stealth[length=1.8mm]}, poppink, thin] (3.8,0.5) -- ++(\a:0.7);
  \draw[-{Stealth[length=1.8mm]}, poppink, thin] (5.6,-0.5) -- ++(\a:0.7);
}
% Oeil
\draw[popblue, thick] (4.6,2.4) circle (0.28);
\fill[popblue] (4.6,2.4) circle (0.1);
\draw[-{Stealth[length=2.5mm]}, poppink, thick] (5.0,0.35) -- (4.75,2.1);
\node[font=\scriptsize, popblue] at (5.6,2.4) {œil de l'observateur};
\node[font=\scriptsize, popdark] at (4.6,-1.6) {Les grains de poussière diffusent la lumière : le faisceau devient visible de côté.};
\end{tikzpicture}
\legende{Le faisceau d'une lampe torche est invisible de côté dans un air parfaitement pur ; les poussières en suspension diffusent une partie de la lumière vers l'œil de l'observateur, révélant le trajet du faisceau.}
\end{popfigure}

\begin{experience}[Couleurs du ciel : diffusion Rayleigh dans l'atmosphère]
\textbf{Matériel.} Un aquarium rectangulaire rempli d'eau, quelques gouttes de lait (ou savon) pour créer une suspension de particules fines ($\ll \lambda$), une lampe torche à lumière blanche (ou un projecteur à fente), un écran blanc.
\textbf{Protocole.}
\begin{enumerate}
\item Placer la lampe sur un côté de l'aquarium, l'écran sur le côté opposé.
\item Observer la couleur de la lumière \emph{transmise} sur l'écran et la couleur de la lumière \emph{diffusée} vue par le côté (perpendiculaire au faisceau).
\item Ajouter progressivement du lait jusqu'à ce que les effets soient nets.
\end{enumerate}
\textbf{Observations.}
\begin{itemize}
\item La lumière diffusée latéralement apparaît \textbf{bleue} (courtes longueurs d'onde fortement diffusées).
\item La lumière transmise sur l'écran apparaît \textbf{jaune-orangée puis rouge} (le bleu a été « retiré » par diffusion latérale).
\item Plus le trajet dans l'eau laiteuse est long (aquarium plus grand ou plus de lait), plus la lumière transmise est rouge.
\end{itemize}
\textbf{Interprétation.} Les particules de lait (tailles $\sim \SI{100}{nm}$) sont bien plus petites que $\lambda_{\text{visible}}$. La diffusion est de type Rayleigh : efficacité $\propto 1/\lambda^4$. Le bleu ($\sim \SI{450}{nm}$) est diffusé $\sim (700/450)^4 \approx 5,7$ fois plus que le rouge ($\sim \SI{700}{nm}$).
\begin{itemize}
\item \textbf{Midi (Soleil zénithal)} : trajet court dans l'atmosphère $\to$ le bleu est diffusé vers l'observateur $\to$ ciel bleu.
\item \textbf{Coucher de soleil} : trajet long (40 fois plus épais) $\to$ le bleu est entièrement diffusé avant d'arriver à l'œil $\to$ lumière directe rouge-orangée.
\end{itemize}
\end{experience}

\begin{exemplebox}[Pourquoi le ciel est-il bleu à midi et rouge-orangé au coucher du soleil à Kribi ?]
La lumière blanche du Soleil contient toutes les couleurs. En traversant l'atmosphère, elle est diffusée par les molécules d'air (azote, dioxygène), bien plus petites que les longueurs d'onde visibles : c'est la diffusion Rayleigh, d'autant plus efficace que $\lambda$ est courte.

\textbf{À midi} : le Soleil est haut, la traversée de l'atmosphère est courte. Le bleu est diffusé dans toutes les directions bien plus que le rouge : quelle que soit la direction où l'on regarde (hors du Soleil), on reçoit du bleu diffusé — le ciel est bleu.

\textbf{Au coucher du soleil sur la plage de Kribi} : la lumière traverse une épaisseur d'atmosphère beaucoup plus grande (jusqu'à 40 fois plus). Le bleu et le vert sont presque totalement diffusés sur le trajet, avant d'arriver jusqu'à toi. Il ne reste dans la lumière directe que le rouge et l'orange : le Soleil et les nuages proches de l'horizon flamboient.

Même physique, deux spectacles différents : tout dépend de l'épaisseur d'air traversée.
\end{exemplebox}

\courssub{Bilan énergétique : un principe unique, la conservation de l'énergie}

Dans chacune des trois interactions, faisons le \emph{bilan d'énergie} comme on ferait un bilan de forces en mécanique : on identifie le système, on liste les énergies entrantes et sortantes, et on écrit la conservation.

\begin{propriete}[Conservation de l'énergie (admise)]
Lors de \textbf{toute} interaction entre la lumière et la matière, l'énergie totale se conserve :
\[ E_{\text{lumière incidente}} + E_{\text{matière avant}} = E_{\text{lumière après}} + E_{\text{matière après}} \]
\begin{itemize}
\item \textbf{Émission} : la matière perd de l'énergie, la lumière en emporte exactement autant : $E_{\text{lumière émise}} = -\Delta E_{\text{matière}}$ ;
\item \textbf{Absorption} : la matière gagne exactement l'énergie perdue par la lumière : $\Delta E_{\text{matière}} = E_{\text{lumière absorbée}}$ ;
\item \textbf{Diffusion} : la fréquence ne change pas, donc (nous le verrons avec $E = h\nu$) la lumière diffusée emporte la même énergie que la lumière incidente ; seule la direction change.
\end{itemize}
\end{propriete}

\begin{methode}[Réaliser un bilan d'énergie lumière-matière]
\begin{enumerate}
\item \textbf{Choisir le système} : la matière (atome, surface, volume) + le champ lumineux incident.
\item \textbf{Identifier l'interaction} : émission, absorption ou diffusion (ou plusieurs simultanément).
\item \textbf{Lister les énergies} : $E_{\text{lum, in}}$, $E_{\text{mat, in}}$, $E_{\text{lum, out}}$, $E_{\text{mat, out}}$.
\item \textbf{Écrire la conservation} : $E_{\text{lum, in}} + E_{\text{mat, in}} = E_{\text{lum, out}} + E_{\text{mat, out}}$.
\item \textbf{Utiliser la quantification} (section 2 et 3) : $\Delta E_{\text{mat}} = \pm h\nu$ pour les transitions discrètes.
\item \textbf{Conclure} sur la variation d'énergie interne (température, état d'excitation) et les caractéristiques de la lumière sortante (fréquence, direction).
\end{enumerate}
\end{methode}

\begin{exoresolu}[Le forgeron du marché Mokolo]
\textbf{Énoncé.} Le forgeron porte un morceau de fer de masse $m = \SI{200}{g}$ à la température $\theta = \SI{900}{\celsius}$. À cette température, le fer émet une lumière orangée et rayonne une puissance lumineuse et thermique totale $P = \SI{150}{W}$. On suppose, pour simplifier, que le fer ne reçoit plus aucune énergie et que toute la puissance rayonnée est perdue. La capacité thermique massique du fer est $c_{\text{fer}} = \SI{450}{J.kg^{-1}.K^{-1}}$.
\begin{enumerate}
\item Quelle énergie le fer perd-il par rayonnement pendant $\Delta t = \SI{10}{s}$ ?
\item En déduire la baisse de température $\Delta\theta$ du fer pendant cette durée.
\item Pourquoi la couleur du fer change-t-elle au cours de ce refroidissement ?
\end{enumerate}
\tcblower
\textbf{Solution détaillée.}
\begin{enumerate}
\item L'énergie rayonnée est $E = P \times \Delta t = 150 \times 10 = \SI{1500}{J} = \SI{1,5}{kJ}$.

Par conservation de l'énergie, cette énergie emportée par le rayonnement (émission) est exactement l'énergie interne perdue par le fer.
\item La baisse d'énergie interne du fer s'écrit $E = m\, c_{\text{fer}}\, \Delta\theta$, d'où :
\[ \Delta\theta = \frac{E}{m\, c_{\text{fer}}} = \frac{1500}{0,200 \times 450} = \frac{1500}{90} \approx \SI{17}{\celsius} \]
Le fer passe d'environ $\SI{900}{\celsius}$ à $\SI{883}{\celsius}$ en dix secondes.
\item La couleur émise par un corps chaud dépend de sa température : plus le corps est chaud, plus la lumière émise est riche en courtes longueurs d'onde. En refroidissant, le fer passe du blanc-orangé au rouge sombre, puis cesse de briller dans le visible (il rayonne alors dans l'infrarouge, invisible). C'est précisément cette loi qui permettra, dans la section 5, de déterminer la température d'un corps chaud par analyse spectrale — y compris celle du Soleil.
\end{enumerate}
\end{exoresolu}

\courssub{Bilan : les trois interactions en un coup d'œil}

\begin{center}
\begin{tabularx}{\linewidth}{|l|X|X|X|}
\hline
\rowcolor{popblueL} \textbf{Interaction} & \textbf{Ce qui se passe} & \textbf{Bilan d'énergie} & \textbf{Exemples du quotidien} \\
\hline
\textbf{Émission} & La matière excitée retombe à un état de moindre énergie en libérant de la lumière. & La matière perd l'énergie emportée par la lumière. & Lampadaires au sodium de Yaoundé, LED des boutiques, braises du feu de bois, fer chauffé du forgeron. \\
\hline
\textbf{Absorption} & La matière capture l'énergie lumineuse et la convertit (chaleur, chimie, excitation). & La matière gagne l'énergie perdue par la lumière. & Tôle noire brûlante au soleil de Garoua, chlorophylle des cacaoyers, filtre coloré, crème solaire. \\
\hline
\textbf{Diffusion} & La matière renvoie la lumière dans toutes les directions, sans changer sa fréquence. & Énergie lumineuse conservée ; seule la direction change. & Ciel bleu, nuages blancs au-dessus de Kribi, faisceau du vidéoprojecteur visible dans la poussière, pages blanches du cahier. \\
\hline
\end{tabularx}
\end{center}

\begin{savaistu}[Blanc à Maroua : la physique dans la garde-robe]
À Maroua, où le thermomètre dépasse régulièrement $\SI{40}{\celsius}$ à l'ombre, les vêtements traditionnels sont souvent blancs ou clairs. Ce n'est pas un hasard : un tissu blanc \emph{diffuse} la quasi-totalité des radiations visibles et renvoie une bonne partie du rayonnement solaire, tandis qu'un tissu noir les \emph{absorbe} et convertit leur énergie en chaleur. À puissance solaire égale (environ $\SI{1000}{W.m^{-2}}$ au sol en plein soleil), un vêtement noir peut s'échauffer de plus de $\SI{15}{\celsius}$ de plus qu'un vêtement blanc. Les nomades du désert ont découvert empiriquement, il y a des millénaires, ce que la physique des interactions lumière-matière explique aujourd'hui : mieux vaut diffuser qu'absorber !
\end{savaistu}

\begin{exercice}[À toi de jouer]
Une feuille de papier blanc et une feuille de papier noir sont posées côte à côte sous le soleil de midi. Après dix minutes, la feuille noire est nettement plus chaude que la blanche, alors que les deux reçoivent la même puissance lumineuse.
\begin{enumerate}
\item Nommer l'interaction lumière-matière dominante pour chaque feuille.
\item Écrire le bilan d'énergie correspondant dans chaque cas.
\item Pourquoi les nuages, formés de gouttelettes d'eau pourtant transparentes, apparaissent-ils blancs ?
\end{enumerate}
\end{exercice}

\begin{corrige}[Exercice : feuilles au soleil]
\begin{enumerate}
\item Feuille noire : \textbf{absorption} dominante (elle absorbe presque toutes les radiations visibles). Feuille blanche : \textbf{diffusion} dominante (elle renvoie la lumière dans toutes les directions).
\item Feuille noire : l'énergie lumineuse absorbée est convertie en énergie interne : $\Delta U = E_{\text{absorbée}} > 0$, la température monte. Feuille blanche : la lumière reçue est renvoyée avec la même énergie (même fréquence) ; presque rien ne reste dans le papier, l'échauffement est faible.
\item Les gouttelettes d'eau sont plus grandes que les longueurs d'onde visibles : elles diffusent toutes les couleurs de la même façon. La superposition de toutes les couleurs diffusées redonne du blanc : le nuage est blanc, même si l'eau pure est transparente.
\end{enumerate}
\end{corrige}

\begin{aretenir}[L'essentiel de la section]
\begin{itemize}
\item La lumière se propage comme une \textbf{onde} ($\lambda = c/\nu$) mais échange son énergie avec la matière \textbf{par grains} : les photons.
\item Trois interactions fondamentales : \textbf{émission} (la matière excitée libère de la lumière), \textbf{absorption} (la matière capte l'énergie lumineuse), \textbf{diffusion} (la matière renvoie la lumière dans toutes les directions, sans changer sa fréquence).
\item Dans les trois cas, \textbf{l'énergie totale se conserve}.
\item La diffusion est d'autant plus efficace que la longueur d'onde est courte (petites particules) : c'est pourquoi le ciel est bleu à midi et rouge au couchant.
\end{itemize}
Dans la prochaine section, nous verrons \emph{pourquoi} la matière n'émet et n'absorbe que certaines couleurs bien précises : c'est la quantification des niveaux d'énergie.
\end{aretenir}

\coursec{Quantification des niveaux d'énergie de la matière}

\courssub{Un constat expérimental troublant : les spectres de raies}

Dans la section précédente, nous avons vu que la lumière et la matière échangent de l'énergie par trois mécanismes : l'\cle{émission}, l'\cle{absorption} et la \cle{diffusion}. Nous allons maintenant regarder de très près \emph{quelle} lumière un atome émet ou absorbe. La réponse, obtenue expérimentalement dès le XIX\textsuperscript{e} siècle, va bouleverser notre vision de la matière.

\begin{experience}[Spectres d'émission : lampe à incandescence contre lampe à vapeur atomique]
\textbf{Matériel} : une lampe à incandescence (filament de tungstène), une lampe à vapeur de sodium (lampadaires orange que l'on voit encore sur certains axes routiers), une lampe à vapeur de mercure, une fente fine, un prisme ou un réseau de diffraction, un écran.

\textbf{Protocole} : on éclaire la fente avec chaque source, on décompose la lumière transmise à l'aide du prisme, et on observe le spectre sur l'écran.

\textbf{Observations} :
\begin{itemize}
\item La lampe à incandescence donne un \cle{spectre continu} : toutes les couleurs de l'arc-en-ciel se succèdent sans interruption, du rouge au violet. Le spectre du Soleil (hors raies d'absorption) présente le même aspect continu.
\item La lampe à vapeur de sodium donne un \cle{spectre de raies} : on n'observe que quelques raies lumineuses colorées, isolées sur fond noir — notamment un doublet jaune-orange intense vers $\lambda = \SI{589}{nm}$. Entre les raies : rien, le noir complet.
\item La lampe à vapeur de mercure donne aussi un spectre de raies, mais avec des raies \emph{différentes} (raies violette, bleue, verte, jaune...).
\end{itemize}

\textbf{Interprétation} : un solide chauffé (filament) émet \emph{toutes} les longueurs d'onde, alors qu'un \emph{atome isolé} (dans une vapeur diluée) n'émet que certaines longueurs d'onde bien précises, qui forment sa « signature ». Chaque élément chimique possède son propre spectre de raies, aussi caractéristique qu'une empreinte digitale.
\end{experience}

\begin{popfigure}
\begin{center}
\begin{tikzpicture}
% Spectre continu
\shade[left color=poppink, right color=poppurple] (0,2.2) rectangle (8,3.2);
\node[left, font=\small\bfseries] at (0,2.7) {Incandescence};
\node[font=\small, align=center] at (4,1.8) {Spectre continu : toutes les longueurs d'onde sont présentes};
% Spectre de raies
\fill[black] (0,-0.6) rectangle (8,0.4);
\draw[line width=2pt, poporange] (4.55,-0.6) -- (4.55,0.4);
\draw[line width=2pt, popgold] (4.75,-0.6) -- (4.75,0.4);
\node[left, font=\small\bfseries] at (0,-0.1) {Vapeur de sodium};
\node[font=\small, align=center] at (4,-1.1) {Spectre de raies : seules quelques longueurs d'onde sont émises};
\draw[->, thick, popmuted] (0,-1.6) -- (8,-1.6) node[right, font=\small] {$\lambda$ croissant};
\node[font=\small, popmuted] at (0.4,-1.9) {rouge};
\node[font=\small, popmuted] at (7.6,-1.9) {violet};
\end{tikzpicture}
\end{center}
\legende{Comparaison entre le spectre continu d'une lampe à incandescence et le spectre de raies d'une lampe à vapeur de sodium (doublet jaune vers $\SI{589}{nm}$).}
\end{popfigure}

Ce constat pose une question fondamentale : \emph{pourquoi} un atome n'émet-il que certaines couleurs et pas d'autres ? Si l'énergie lumineuse émise ne prend que certaines valeurs, c'est que l'énergie de l'atome lui-même ne peut prendre que certaines valeurs. C'est l'idée révolutionnaire introduite par Niels Bohr en 1913.

\courssub{Les niveaux d'énergie de l'atome sont quantifiés}

\begin{definition}[Quantification, niveau d'énergie, état fondamental, état excité]
On dit qu'une grandeur est \cle{quantifiée} lorsqu'elle ne peut prendre que certaines valeurs discrètes (bien déterminées et séparées), et non pas n'importe quelle valeur continue.

L'énergie d'un atome est quantifiée : elle ne peut prendre que des valeurs particulières $E_1, E_2, E_3, \ldots$ appelées \cle{niveaux d'énergie} de l'atome.

\begin{itemize}
\item Le niveau d'énergie le plus bas, noté $E_1$, est le \cle{niveau fondamental} : c'est l'état stable de l'atome, dans lequel il se trouve naturellement lorsqu'il n'est pas perturbé.
\item Les niveaux $E_2, E_3, \ldots$, d'énergies supérieures, sont les \cle{niveaux excités} : l'atome ne peut y séjourner que de façon transitoire (typiquement $10^{-8}$ s) avant de redescendre vers un niveau inférieur.
\end{itemize}
\end{definition}

\begin{savaistu}[L'escalier énergétique]
Imagine un escalier : tu peux poser le pied sur la première marche, la deuxième, la troisième... mais jamais \emph{entre} deux marches. L'énergie d'un atome fonctionne pareil : elle est « en marches d'escalier », pas « en rampe continue ». C'est Niels Bohr (prix Nobel 1922) qui a proposé ce modèle pour l'hydrogène en 1913, en s'inspirant des travaux de Planck et d'Einstein. La mécanique quantique, développée dans les années 1920, a ensuite justifié cette quantification de façon rigoureuse.
\end{savaistu}

\begin{propriete}[Quantification de l'énergie des atomes — admise]
L'énergie d'un atome ne peut prendre que des valeurs discrètes $E_n$ (avec $n = 1, 2, 3, \ldots$), caractéristiques de l'élément chimique. Toute variation de l'énergie de l'atome correspond à une \cle{transition} d'un niveau $E_n$ vers un niveau $E_p$, et s'accompagne d'un échange d'énergie avec l'extérieur exactement égal à
\[ \Delta E = |E_n - E_p| \]
Cette propriété est admise en Première ; elle sera démontrée dans le cadre de la mécanique quantique en études supérieures.
\end{propriete}

\courssub{Le diagramme de niveaux d'énergie}

Pour visualiser les niveaux d'énergie d'un atome, on utilise un \cle{diagramme de niveaux d'énergie} : chaque niveau est représenté par un trait horizontal, placé à une hauteur proportionnelle à son énergie, l'axe vertical étant gradué en énergie.

\textbf{Conventions essentielles} (à connaître absolument pour le Probatoire) :
\begin{itemize}
\item Par convention, on choisit $E = 0$ pour l'atome \cle{ionisé}, c'est-à-dire l'atome ayant perdu un électron (électron arraché, infiniment éloigné et au repos).
\item Les états où l'électron est \emph{lié} au noyau ont donc des énergies \textbf{négatives} : il faut \emph{fournir} de l'énergie à l'atome pour les amener à zéro (arracher l'électron).
\item Le niveau fondamental est le niveau d'énergie \emph{la plus négative} (le plus bas sur le diagramme).
\item Plus un niveau est haut sur le diagramme (énergie moins négative), plus l'atome est excité, et moins l'électron est lié au noyau.
\end{itemize}

\begin{definition}[L'électronvolt]
L'\cle{électronvolt} (symbole eV) est une unité d'énergie adaptée à l'échelle atomique. C'est l'énergie acquise par un électron accéléré sous une tension de $\SI{1}{V}$ :
\[ 1~\text{eV} = e \times U = \num{1,60e-19}~\text{C} \times 1~\text{V} = \num{1,60e-19}~\text{J} \]
Les énergies des atomes étant de l'ordre de $10^{-19}$ à $10^{-18}$ joule, le joule est malcommode à cette échelle ; l'électronvolt donne des valeurs simples (quelques eV), d'où son intérêt pratique.
\end{definition}

\begin{exemplebox}[Diagramme simplifié de l'atome d'hydrogène]
L'atome d'hydrogène, le plus simple de tous (un proton, un électron), possède les premiers niveaux suivants :
\begin{center}
\begin{tabular}{|c|c|c|}
\hline
\rowcolor{popblueL} Niveau & Énergie (eV) & Énergie (J) \\
\hline
$E_1$ (fondamental) & $-13,6$ & $-\num{2,18e-18}$ \\
\hline
$E_2$ & $-3,4$ & $-\num{5,44e-19}$ \\
\hline
$E_3$ & $-1,5$ & $-\num{2,40e-19}$ \\
\hline
Atome ionisé & $0$ & $0$ \\
\hline
\end{tabular}
\end{center}
Vérification de la conversion pour $E_1$ : $E_1 = -13,6 \times \num{1,60e-19} = -\num{2,18e-18}~\text{J}$. Pour $E_3$ : $-1,5 \times \num{1,60e-19} = -\num{2,40e-19}~\text{J}$. L'énergie $E_\infty = 0$ correspond à l'ionisation : il faut fournir $13,6$ eV à l'atome dans son état fondamental pour arracher son électron. C'est l'\emph{énergie d'ionisation} de l'hydrogène.
\end{exemplebox}

\begin{popfigure}
\begin{center}
\begin{tikzpicture}[scale=1]
% Axe énergie
\draw[->, thick] (0,-5.2) -- (0,1.2) node[above, font=\small] {$E$ (eV)};
% Graduations
\node[left, font=\small] at (1,0) {$0$};
\node[left, font=\small] at (1,-1.5) {$-1{,}5$};
\node[left, font=\small] at (1,-3.4) {$-3{,}4$};
\node[left, font=\small] at (1,-4.6) {$-13{,}6$};
% Niveaux
\draw[very thick, popdark] (1,0) -- (6,0) node[right, font=\small] {$E_\infty = 0$ (ionisé)};
\draw[very thick, popblue] (1,-1.5) -- (6,-1.5) node[right, font=\small] {$E_3 = -1{,}5$ eV};
\draw[very thick, popblue] (1,-3.4) -- (6,-3.4) node[right, font=\small] {$E_2 = -3{,}4$ eV};
\draw[very thick, popblue] (1,-4.6) -- (6,-4.6) node[right, font=\small] {$E_1 = -13{,}6$ eV (fondamental)};
% Transition émission E3 -> E2
\draw[->, very thick, poporange] (3.2,-1.5) -- (3.2,-3.4);
\node[right, font=\small, poporange] at (3.3,-2.45) {émission : $\Delta E = 1{,}9$ eV};
% Transition absorption E1 -> E3
\draw[->, very thick, popgreen] (4.6,-4.6) -- (4.6,-1.5);
\node[right, font=\small, popgreen] at (4.7,-3.1) {absorption : $\Delta E = 12{,}1$ eV};
\end{tikzpicture}
\end{center}
\legende{Diagramme simplifié des niveaux d'énergie de l'atome d'hydrogène. Flèche vers le bas : émission (l'atome cède de l'énergie). Flèche vers le haut : absorption (l'atome reçoit de l'énergie). L'échelle verticale n'est pas linéaire pour des raisons de lisibilité.}
\end{popfigure}

\courssub{Transitions entre niveaux : le bilan énergétique}

Un atome ne change de niveau d'énergie que s'il échange avec l'extérieur \emph{exactement} la différence d'énergie entre les deux niveaux. Deux cas se présentent :

\begin{itemize}
\item \textbf{Absorption} : l'atome passe d'un niveau $E_n$ à un niveau \emph{supérieur} $E_p$ ($E_p > E_n$). Il doit \emph{recevoir} l'énergie $\Delta E = E_p - E_n > 0$, par exemple en absorbant un grain de lumière (nous verrons à la section suivante qu'il s'agit d'un photon).
\item \textbf{Émission} : l'atome, préalablement excité, redescend d'un niveau $E_n$ vers un niveau \emph{inférieur} $E_p$ ($E_p < E_n$). Il \emph{cède} l'énergie $\Delta E = E_n - E_p > 0$, sous forme de lumière émise.
\end{itemize}

Dans les deux cas, l'énergie échangée s'écrit de manière générale :
\[ \Delta E = |E_n - E_p| \]

C'est précisément ce qui explique les spectres de raies : comme les niveaux d'énergie sont quantifiés, les différences $\Delta E$ ne peuvent prendre que certaines valeurs, donc la lumière émise ou absorbée ne peut avoir que certaines énergies — c'est-à-dire certaines longueurs d'onde bien précises. Le spectre de raies est la \emph{photographie} des transitions permises entre les niveaux de l'atome.

\begin{methode}[Lire et exploiter un diagramme de niveaux d'énergie]
\begin{enumerate}
\item \textbf{Repérer le niveau fondamental} : c'est le niveau le plus bas (énergie la plus négative). L'atome non perturbé s'y trouve.
\item \textbf{Identifier le sens de la transition} : si le niveau d'arrivée est au-dessus du niveau de départ, c'est une \emph{absorption} (l'atome gagne de l'énergie) ; s'il est en dessous, c'est une \emph{émission} (l'atome perd de l'énergie).
\item \textbf{Calculer l'énergie échangée} : $\Delta E = |E_{\text{arrivée}} - E_{\text{départ}}|$, toujours positive, en soustrayant \emph{algébriquement} les énergies (attention aux signes négatifs !).
\item \textbf{Convertir si nécessaire} : multiplier par $\num{1,60e-19}$ pour passer des eV aux joules.
\item \textbf{Conclure} : préciser si l'atome a absorbé ou émis cette énergie.
\end{enumerate}
\end{methode}

\begin{exoresolu}[Transition $E_3 \to E_2$ de l'atome d'hydrogène]
\textbf{Énoncé.} Un atome d'hydrogène, initialement au niveau $E_3 = -1,5$ eV, revient au niveau $E_2 = -3,4$ eV.
\begin{enumerate}
\item S'agit-il d'une émission ou d'une absorption ? Justifier.
\item Calculer l'énergie échangée $\Delta E$ en eV puis en joules.
\item Quelle énergie faut-il fournir à un atome d'hydrogène dans son état fondamental pour l'ioniser ?
\end{enumerate}
\tcblower
\textbf{Solution détaillée.}
\begin{enumerate}
\item Le niveau d'arrivée $E_2 = -3,4$ eV est \emph{en dessous} du niveau de départ $E_3 = -1,5$ eV : l'atome redescend vers un état moins énergétique. Il s'agit donc d'une \textbf{émission} : l'atome cède de l'énergie sous forme de lumière.
\item L'énergie échangée est :
\begin{align*}
\Delta E &= |E_3 - E_2| = |-1,5 - (-3,4)| \\
&= |-1,5 + 3,4| = |1,9| = 1,9~\text{eV}
\end{align*}
Conversion en joules :
\[ \Delta E = 1,9 \times \num{1,60e-19} = \num{3,04e-19}~\text{J} \]
\item L'état fondamental est $E_1 = -13,6$ eV et l'atome ionisé correspond à $E_\infty = 0$. L'énergie d'ionisation vaut :
\[ \Delta E_{\text{ion}} = |E_\infty - E_1| = |0 - (-13,6)| = 13,6~\text{eV} = \num{2,18e-18}~\text{J} \]
Il faut donc fournir $13,6$ eV à l'atome pour arracher son électron.
\end{enumerate}
\end{exoresolu}

\begin{attention}[Les trois pièges classiques du diagramme énergétique]
\textbf{Piège 1 — Oublier le signe négatif des énergies.} Les énergies des états liés sont \emph{négatives}. Écrire $E_2 = 3,4$ eV au lieu de $-3,4$ eV fausse tout le calcul : on trouverait $|-1,5 - 3,4| = 4,9$ eV au lieu de $1,9$ eV. Recopie toujours les énergies avec leur signe.

\textbf{Piège 2 — Se tromper de sens pour la différence.} $\Delta E$ est une énergie \emph{échangée} : elle se calcule \textbf{toujours en valeur absolue}, $\Delta E = |E_n - E_p|$, et c'est l'analyse du sens de la transition (montée ou descente sur le diagramme) qui permet de dire si l'atome a absorbé ou émis. Une énergie échangée négative n'a pas de sens physique.

\textbf{Piège 3 — Confondre eV et joules.} Si l'énoncé demande le résultat en joules, la conversion par $\num{1,60e-19}$ est obligatoire ; un résultat laissé en eV est incomplet.
\end{attention}

\begin{savaistu}[Pourquoi le ciel des villes est-il orange ?]
Les lampadaires à vapeur de sodium, longtemps utilisés pour éclairer les rues de Douala et de Yaoundé, exploitent directement la quantification : les atomes de sodium excités par décharge électrique redescendent vers leur état fondamental en émettant presque exclusivement la raie jaune caractéristique à $\SI{589}{nm}$. C'est aussi grâce aux spectres de raies que les astrophysiciens déterminent la composition des étoiles : chaque raie d'absorption du spectre solaire trahit la présence d'un élément chimique précis dans l'atmosphère du Soleil. L'hélium a d'ailleurs été découvert \emph{dans le Soleil} (1868) avant d'être identifié sur Terre !
\end{savaistu}

\begin{exercice}[Pour s'entraîner]
On donne pour l'atome d'hydrogène : $E_1 = -13,6$ eV ; $E_2 = -3,4$ eV ; $E_3 = -1,5$ eV.
\begin{enumerate}
\item Un atome d'hydrogène passe du niveau fondamental au niveau $E_3$. Émission ou absorption ? Calculer l'énergie échangée en eV puis en joules.
\item Un atome excité au niveau $E_3$ peut-il redescendre directement au niveau fondamental ? Si oui, quelle énergie cède-t-il ?
\end{enumerate}
\end{exercice}

\begin{corrige}[Exercice]
\begin{enumerate}
\item Le niveau d'arrivée $E_3$ est au-dessus de $E_1$ : c'est une \textbf{absorption}. L'énergie à fournir est $\Delta E = |E_3 - E_1| = |-1,5 - (-13,6)| = 12,1$ eV, soit $\Delta E = 12,1 \times \num{1,60e-19} = \num{1,94e-18}~\text{J}$.
\item Oui, la transition directe $E_3 \to E_1$ est possible : c'est une émission d'énergie $\Delta E = |E_3 - E_1| = 12,1$ eV $= \num{1,94e-18}~\text{J}$. (L'atome peut aussi redescendre en deux étapes : $E_3 \to E_2$ puis $E_2 \to E_1$, en émettant $1,9$ eV puis $10,2$ eV — la somme redonne bien $12,1$ eV, l'énergie se conserve.)
\end{enumerate}
\end{corrige}

\begin{aretenir}[L'essentiel de la section]
\begin{itemize}
\item Les atomes émettent et absorbent des \textbf{spectres de raies} : leurs échanges d'énergie ne prennent que des valeurs discrètes.
\item L'énergie d'un atome est \textbf{quantifiée} : niveaux $E_1$ (fondamental, stable), $E_2, E_3, \ldots$ (excités), représentés sur un diagramme gradué en eV.
\item Convention : $E = 0$ pour l'atome ionisé ; les états liés ont des énergies \textbf{négatives}.
\item Toute transition entre deux niveaux s'accompagne d'un échange d'énergie $\Delta E = |E_n - E_p|$ : absorption si l'atome monte, émission s'il descend.
\item Unité adaptée : l'électronvolt, $1~\text{eV} = \num{1,60e-19}~\text{J}$.
\end{itemize}
\end{aretenir}

Nous savons désormais \emph{combien} d'énergie un atome échange lors d'une transition. Reste à comprendre \emph{comment} cette énergie est transportée par la lumière : c'est là qu'entre en scène le \cle{photon}, le grain de lumière d'Einstein, objet de la section suivante.

\coursec{Le photon et son énergie : $E = h\nu$}

Dans la section précédente, nous avons découvert un fait bouleversant : l'énergie d'un atome ne peut prendre que certaines valeurs bien déterminées, les \cle{niveaux d'énergie}. Un atome ne peut donc absorber ou émettre de la lumière que par \emph{quantités discrètes} d'énergie, exactement égales à l'écart entre deux de ses niveaux. Mais alors, si la matière n'échange l'énergie lumineuse que par « paquets », n'est-ce pas parce que la lumière elle-même est faite de paquets ? C'est précisément la révolution conceptuelle que nous allons étudier ici : la lumière, qui se propage comme une onde, se comporte comme un \emph{grain} lorsqu'elle échange de l'énergie avec la matière. Ce grain s'appelle le \cle{photon}.

\courssub{Du modèle ondulatoire au modèle corpusculaire : un siècle de débat}

Depuis les travaux de Huygens (XVII\textsuperscript{e} siècle), puis les expériences d'interférences de Young (1801) et la théorie électromagnétique de Maxwell (1865), la lumière était décrite comme une \emph{onde} : une perturbation qui se propage, caractérisée par une fréquence $\nu$ et une longueur d'onde $\lambda$. Ce modèle explique parfaitement la diffraction et les interférences... mais il échoue devant certains faits expérimentaux.

À la fin du XIX\textsuperscript{e} siècle, deux énigmes résistent au modèle ondulatoire :
\begin{itemize}
\item le \emph{rayonnement du corps chaud} : la théorie classique prévoit qu'un four chauffé devrait émettre une énergie infinie dans l'ultraviolet (« catastrophe ultraviolette »), ce qui est absurde ;
\item l'\emph{effet photoélectrique} : une plaque de zinc éclairée par une lumière ultraviolette éjecte des électrons, mais une lumière rouge, même très intense, n'y parvient jamais. Or, selon le modèle ondulatoire, une onde rouge suffisamment intense devrait apporter assez d'énergie !
\end{itemize}

En \textbf{1900}, Max \textbf{Planck} résout la première énigme en faisant une hypothèse audacieuse : les échanges d'énergie entre la matière et le rayonnement ne se font pas de façon continue, mais par \emph{quanta} (grains) d'énergie proportionnels à la fréquence. En \textbf{1905}, Albert \textbf{Einstein} va plus loin pour expliquer l'effet photoélectrique : ce ne sont pas seulement les échanges qui sont quantifiés, c'est la lumière \emph{elle-même} qui est constituée de particules d'énergie, que l'on nommera plus tard les \cle{photons}. Cette audace lui vaudra le prix Nobel de physique en 1921.

\begin{savaistu}[Pourquoi la lumière rouge n'arrache-t-elle aucun électron ?]
Imagine des enfants qui lancent des balles sur une bouteille posée sur un mur. Si chaque balle est trop légère, même en lançant mille balles, la bouteille ne tombera pas. Il faut qu'\emph{une seule} balle ait assez d'énergie. De même, un électron ne peut absorber qu'\emph{un seul photon à la fois} : si ce photon est trop « léger » (lumière rouge), aucun électron ne part, quelle que soit l'intensité lumineuse. Un photon ultraviolet, lui, est assez énergétique à lui tout seul. C'est exactement ce que prédit le modèle du photon.
\end{savaistu}

\courssub{Le photon : définition et caractéristiques}

\begin{definition}[Le photon]
Le \cle{photon} est le grain d'énergie élémentaire constituant la lumière (et plus généralement toute onde électromagnétique). C'est une particule :
\begin{itemize}
\item de \textbf{masse nulle} : $m = 0$ ;
\item de \textbf{charge électrique nulle} : $q = 0$ ;
\item se déplaçant dans le vide à la célérité de la lumière $c = \num{3,00e8}~\text{m}\cdot\text{s}^{-1}$ ;
\item transportant une énergie $E$ qui dépend uniquement de la fréquence $\nu$ de la radiation associée.
\end{itemize}
\end{definition}

\begin{attention}[Deux idées fausses à éliminer tout de suite]
\begin{itemize}
\item Le photon n'est \emph{pas} une petite bille de matière qui ralentirait dans le verre ou l'eau : sa masse est rigoureusement nulle et, dans le vide, sa vitesse est toujours $c$, quelle que soit sa couleur.
\item Le photon ne possède \emph{pas} de couleur en soi : c'est son énergie (donc sa fréquence) qui détermine la sensation de couleur lorsqu'il frappe notre rétine. Un photon de $\lambda = 589~\text{nm}$ est perçu comme jaune.
\end{itemize}
\end{attention}

\courssub{La constante de Planck et la relation fondamentale $E = h\nu$}

L'hypothèse de Planck, reprise par Einstein, s'énonce ainsi : l'énergie d'un photon est \emph{proportionnelle} à la fréquence de la radiation. Le coefficient de proportionnalité est une constante universelle de la nature, la \cle{constante de Planck}, notée $h$.

\begin{propriete}[Énergie d'un photon — relation fondamentale (admise)]
L'énergie d'un photon associé à une radiation de fréquence $\nu$ est :
\[ E = h\,\nu \]
avec :
\begin{itemize}
\item $E$ : énergie du photon, en joules (J) ;
\item $\nu$ : fréquence de la radiation, en hertz (Hz), c'est-à-dire en $\text{s}^{-1}$ ;
\item $h = \num{6,63e-34}~\text{J}\cdot\text{s}$ : constante de Planck.
\end{itemize}
Cette relation est \emph{admise} au niveau Première ; sa connaissance et son exploitation sont exigibles au Probatoire.
\end{propriete}

\textbf{Vérification par analyse dimensionnelle.} Le produit $h\nu$ s'exprime en $\text{J}\cdot\text{s} \times \text{s}^{-1} = \text{J}$ : on obtient bien une énergie. La relation est homogène.

\textbf{Lecture physique.} La petitesse extraordinaire de $h$ (de l'ordre de $10^{-34}$) explique pourquoi le caractère « granulaire » de la lumière nous échappe dans la vie quotidienne : chaque photon transporte une énergie minuscule, et les sources ordinaires en émettent des milliards de milliards chaque seconde, ce qui donne l'illusion d'un flux continu.

\courssub{De la fréquence à la longueur d'onde : $E = \dfrac{hc}{\lambda}$}

Dans le vide, une onde électromagnétique de fréquence $\nu$ se propage à la célérité $c$. Par définition de la longueur d'onde (distance parcourue pendant une période $T = 1/\nu$) :

\begin{propriete}[Relation entre longueur d'onde et fréquence dans le vide]
\[ \lambda = \frac{c}{\nu} \qquad \text{soit de façon équivalente} \qquad \nu = \frac{c}{\lambda} \]
avec $\lambda$ en mètres (m), $\nu$ en hertz (Hz) et $c = \num{3,00e8}~\text{m}\cdot\text{s}^{-1}$.
\end{propriete}

En combinant les deux relations $E = h\nu$ et $\nu = c/\lambda$, on obtient l'expression de l'énergie d'un photon en fonction de sa longueur d'onde :

\begin{aretenir}[Les deux formes à connaître par cœur]
\[ \boxed{E = h\nu} \qquad \text{et} \qquad \boxed{E = \frac{hc}{\lambda}} \]
Conséquence immédiate : l'énergie d'un photon est \textbf{inversement proportionnelle} à la longueur d'onde. Plus $\lambda$ est \emph{petite}, plus le photon est \emph{énergétique}.
\end{aretenir}

Cette conséquence a des effets concrets sur notre santé : les photons ultraviolets ($\lambda < 400~\text{nm}$) sont plus énergétiques que les photons visibles, et bien plus que les infrarouges. C'est pourquoi ce sont les UV du Soleil qui provoquent les coups de soleil et peuvent endommager l'ADN de nos cellules, alors que la lumière visible et l'infrarouge (qui ne fait que chauffer) sont sans danger à intensité modérée. Sous le soleil de Yaoundé ou de Maroua, la crème solaire ne protège pas de la chaleur : elle bloque les photons UV, les plus énergétiques !

\begin{popfigure}
\begin{center}
\begin{tikzpicture}
\begin{axis}[
  width=0.85\linewidth, height=6cm,
  xlabel={$1/\lambda$ \quad ($10^{6}~\text{m}^{-1}$)},
  ylabel={$E$ \quad (eV)},
  xmin=0, xmax=6.5, ymin=0, ymax=8,
  xtick={0,1,2,3,4,5,6},
  ytick={0,2,4,6,8},
  grid=both, grid style={popline!40},
  axis lines=left,
  label style={font=\small},
  tick label style={font=\small},
]
\addplot[popcurve,domain=0:6.2,samples=50]{1.24*x};
\addplot[only marks,mark=*,mark size=1.8pt,poporange] coordinates {(1.70,2.11)};
\node[poporange,font=\small,anchor=south west] at (axis cs:1.85,2.11) {photon jaune du sodium};
\node[popblue,font=\small,anchor=north west] at (axis cs:3.3,5.6) {$E = hc/\lambda$ : droite de pente $hc$};
\end{axis}
\end{tikzpicture}
\end{center}
\legende{Énergie d'un photon (en eV) en fonction de l'inverse de la longueur d'onde : la relation $E = hc/\lambda$ est une proportionnalité. En unités pratiques, $E(\text{eV}) = \dfrac{1{,}24}{\lambda(\mu\text{m})}$.}
\end{popfigure}

\courssub{Méthode de calcul et exemple chiffré}

\begin{methode}[Calculer l'énergie d'un photon à partir de sa longueur d'onde]
\begin{enumerate}
\item \textbf{Convertir} la longueur d'onde $\lambda$ en \emph{mètres} : $1~\text{nm} = 10^{-9}~\text{m}$, $1~\mu\text{m} = 10^{-6}~\text{m}$.
\item \textbf{Calculer la fréquence} : $\nu = \dfrac{c}{\lambda}$ avec $c = \num{3,00e8}~\text{m}\cdot\text{s}^{-1}$.
\item \textbf{Calculer l'énergie en joules} : $E = h\nu$ (ou directement $E = \dfrac{hc}{\lambda}$).
\item \textbf{Convertir éventuellement en électronvolts} : $1~\text{eV} = \num{1,60e-19}~\text{J}$, donc $E(\text{eV}) = \dfrac{E(\text{J})}{\num{1,60e-19}}$.
\end{enumerate}
\end{methode}

\begin{attention}[Le piège classique du Probatoire]
L'erreur la plus fréquente consiste à \textbf{laisser $\lambda$ en nanomètres} dans le calcul. Avec $\lambda = 589$ au lieu de $\num{589e-9}~\text{m}$, le résultat est faux d'un facteur $10^{9}$ ! Écris systématiquement la conversion sur ta copie : $\lambda = 589~\text{nm} = \num{589e-9}~\text{m}$. Autre piège : oublier de convertir les joules en électronvolts quand l'énoncé demande le résultat en eV.
\end{attention}

\begin{exemplebox}[Le photon jaune de la lampe à vapeur de sodium]
Les lampes à vapeur de sodium, qui éclairent en orange certaines rues et certains stades, émettent une lumière quasi monochromatique de longueur d'onde $\lambda = 589~\text{nm}$. Calculons l'énergie d'un photon émis.

\textbf{Étape 1 — Conversion :} $\lambda = 589~\text{nm} = \num{589e-9}~\text{m} = \num{5,89e-7}~\text{m}$.

\textbf{Étape 2 — Fréquence :}
\[ \nu = \frac{c}{\lambda} = \frac{\num{3,00e8}}{\num{5,89e-7}} = \num{5,09e14}~\text{Hz} \]

\textbf{Étape 3 — Énergie en joules :}
\[ E = h\nu = \num{6,63e-34} \times \num{5,09e14} = \num{3,37e-19}~\text{J} \]

\textbf{Étape 4 — Conversion en électronvolts :}
\[ E = \frac{\num{3,37e-19}}{\num{1,60e-19}} \approx 2{,}11~\text{eV} \]
Un photon jaune transporte donc une énergie d'environ $2~\text{eV}$ : c'est l'ordre de grandeur à retenir pour tout photon visible.
\end{exemplebox}

\courssub{Ordres de grandeur à travers le spectre électromagnétique}

Appliquons la relation $E = hc/\lambda$ à différents domaines du spectre, des ondes radio aux rayons gamma. Le tableau suivant rassemble les ordres de grandeur essentiels.

\begin{center}
\begin{tabularx}{\linewidth}{|l|X|X|X|}
\hline
\rowcolor{popblueL} \textbf{Domaine} & \textbf{Longueur d'onde $\lambda$} & \textbf{Fréquence $\nu$} & \textbf{Énergie du photon} \\
\hline
Ondes radio & $> 1~\text{m}$ & $< \num{3e8}~\text{Hz}$ & $< \num{1e-6}~\text{eV}$ \\
\hline
Micro-ondes & $1~\text{mm}$ à $1~\text{m}$ & $\num{3e8}$ à $\num{3e11}~\text{Hz}$ & $\num{1e-6}$ à $\num{1e-3}~\text{eV}$ \\
\hline
Infrarouge (IR) & $800~\text{nm}$ à $1~\text{mm}$ & $\num{3e11}$ à $\num{4e14}~\text{Hz}$ & $\num{1e-3}$ à $1{,}6~\text{eV}$ \\
\hline
Visible & $400$ à $800~\text{nm}$ & $\num{4e14}$ à $\num{8e14}~\text{Hz}$ & $1{,}6$ à $3{,}1~\text{eV}$ \\
\hline
Ultraviolet (UV) & $10$ à $400~\text{nm}$ & $\num{8e14}$ à $\num{3e16}~\text{Hz}$ & $3{,}1$ à $124~\text{eV}$ \\
\hline
Rayons X & $\num{1e-11}$ à $10~\text{nm}$ & $\num{3e16}$ à $\num{3e19}~\text{Hz}$ & $124~\text{eV}$ à $124~\text{keV}$ \\
\hline
Rayons $\gamma$ & $< \num{1e-11}~\text{m}$ & $> \num{3e19}~\text{Hz}$ & $> 124~\text{keV}$ \\
\hline
\end{tabularx}
\end{center}

\begin{popfigure}
\begin{center}
\begin{tikzpicture}[x=1cm,y=1cm]
% Bande du spectre
\foreach \i/\col in {0/popgreen,1/popgreen!70!popblue,2/popblue,3/poppurple,4/poppink,5/poporange}{
  \fill[\col!60] (\i*1.9,0) rectangle (\i*1.9+1.9,0.9);
}
\draw[thick,popdark] (0,0) rectangle (11.4,0.9);
% Etiquettes domaines
\node[font=\scriptsize\bfseries,popdark] at (0.95,0.45) {Radio};
\node[font=\scriptsize\bfseries,popdark] at (2.85,0.45) {Micro-ondes};
\node[font=\scriptsize\bfseries,popdark] at (4.75,0.45) {Infrarouge};
\node[font=\scriptsize\bfseries,popdark] at (6.65,0.45) {Visible};
\node[font=\scriptsize\bfseries,popdark] at (8.55,0.45) {UV};
\node[font=\scriptsize\bfseries,popdark] at (10.45,0.45) {X et $\gamma$};
% Fleche lambda
\draw[-{Stealth},thick,popdark] (0.3,-0.5) -- (11.1,-0.5);
\node[font=\small,popdark,anchor=west] at (11.2,-0.5) {$\lambda$ croissant};
\draw[-{Stealth},thick,popdark] (11.1,-1.2) -- (0.3,-1.2);
\node[font=\small,popdark,anchor=east] at (0.2,-1.2) {$E = h\nu$ croissant};
% Graduations lambda
\foreach \x/\lab in {0.3/$10^{3}$ m,1.9/1 m,3.8/$10^{-3}$ m,5.7/$10^{-6}$ m,7.6/$10^{-9}$ m,9.5/$10^{-12}$ m,11.1/$10^{-15}$ m}{
  \draw[popdark] (\x,-0.35) -- (\x,-0.65);
  \node[font=\scriptsize,popdark,anchor=north] at (\x,-0.7) {\lab};
}
% Energies
\node[font=\scriptsize,popdark,anchor=south] at (0.95,1.0) {$E < 10^{-6}$ eV};
\node[font=\scriptsize,popdark,anchor=south] at (6.65,1.0) {$E \approx 2$ eV};
\node[font=\scriptsize,popdark,anchor=south] at (10.45,1.0) {$E > 10^{5}$ eV};
\end{tikzpicture}
\end{center}
\legende{Le spectre électromagnétique : lorsque la longueur d'onde diminue (vers la droite), l'énergie des photons augmente. Les photons radio sont les moins énergétiques, les photons gamma les plus énergétiques.}
\end{popfigure}

Retenons trois valeurs repères :
\begin{itemize}
\item un \textbf{photon radio} (par exemple celui d'une station FM à $\nu = 100~\text{MHz}$) : $E = h\nu \approx \num{6,6e-26}~\text{J} \approx \num{4e-7}~\text{eV}$ — totalement inoffensif ;
\item un \textbf{photon visible} : $E \approx 2~\text{eV} \approx \num{3e-19}~\text{J}$ ;
\item un \textbf{photon UV} ($\lambda = 300~\text{nm}$) : $E \approx 4{,}1~\text{eV}$, soit le double d'un photon visible — assez pour briser des liaisons chimiques dans la peau.
\end{itemize}

\begin{savaistu}[Combien de photons émet une ampoule LED ?]
Une ampoule LED de $5~\text{W}$ éclaire une chambre à Douala. En supposant que toute la puissance électrique est convertie en lumière jaune-verte ($E_{\text{photon}} \approx \num{4e-19}~\text{J}$), le nombre de photons émis chaque seconde est :
\[ N = \frac{P}{E_{\text{photon}}} = \frac{5}{\num{4e-19}} \approx 10^{19}~\text{photons par seconde} \]
Dix milliards de milliards de photons chaque seconde ! Voilà pourquoi la lumière nous paraît continue : les grains sont si nombreux et si petits que notre œil ne perçoit qu'un flux uniforme, comme le sable d'une plage paraît lisse alors qu'il est fait de grains.
\end{savaistu}

\courssub{La dualité onde-corpuscule}

Faut-il alors abandonner le modèle ondulatoire ? Non ! Les expériences d'interférences et de diffraction restent inexplicables sans ondes. La physique moderne adopte une vision double, appelée \cle{dualité onde-corpuscule} :

\begin{aretenir}[Dualité onde-corpuscule]
La lumière présente un double aspect :
\begin{itemize}
\item elle \textbf{se propage} comme une \emph{onde} électromagnétique de fréquence $\nu$ et de longueur d'onde $\lambda$ (ce qui explique les interférences et la diffraction) ;
\item elle \textbf{échange son énergie} avec la matière par \emph{grains}, les photons, d'énergie $E = h\nu$ (ce qui explique l'effet photoélectrique et les spectres de raies).
\end{itemize}
Les deux aspects sont indissociables : c'est la fréquence $\nu$ de l'onde qui fixe l'énergie du corpuscule.
\end{aretenir}

\begin{exoresolu}[Photon rouge ou photon bleu : qui est le plus énergétique ?]
\textbf{Énoncé.} Un laser rouge émet une radiation de longueur d'onde $\lambda_R = 650~\text{nm}$ ; un pointeur violet émet $\lambda_V = 405~\text{nm}$.
\begin{enumerate}
\item Calculer l'énergie de chaque photon, en joules puis en électronvolts.
\item Quel photon est le plus énergétique ? Justifier.
\item Un détecteur nécessite au minimum $2{,}8~\text{eV}$ pour réagir. Quel(s) laser(s) peut(vent) le déclencher ?
\end{enumerate}
On donne : $h = \num{6,63e-34}~\text{J}\cdot\text{s}$ ; $c = \num{3,00e8}~\text{m}\cdot\text{s}^{-1}$ ; $1~\text{eV} = \num{1,60e-19}~\text{J}$.
\tcblower
\textbf{Solution détaillée.}

\textbf{1. Énergies des photons.} On utilise $E = \dfrac{hc}{\lambda}$ après conversion de $\lambda$ en mètres.

Pour le laser rouge, $\lambda_R = \num{650e-9}~\text{m}$ :
\begin{align*}
E_R &= \frac{hc}{\lambda_R} = \frac{\num{6,63e-34} \times \num{3,00e8}}{\num{650e-9}} \\
&= \num{3,06e-19}~\text{J} = \frac{\num{3,06e-19}}{\num{1,60e-19}} \approx 1{,}91~\text{eV}
\end{align*}

Pour le pointeur violet, $\lambda_V = \num{405e-9}~\text{m}$ :
\begin{align*}
E_V &= \frac{hc}{\lambda_V} = \frac{\num{6,63e-34} \times \num{3,00e8}}{\num{405e-9}} \\
&= \num{4,91e-19}~\text{J} \approx 3{,}07~\text{eV}
\end{align*}

\textbf{2. Comparaison.} $E_V > E_R$ : le photon violet est le plus énergétique. C'est cohérent avec $E = hc/\lambda$ : à longueur d'onde plus petite correspond une énergie plus grande.

\textbf{3. Déclenchement du détecteur.} Le seuil est de $2{,}8~\text{eV}$. Or $E_R = 1{,}91~\text{eV} < 2{,}8~\text{eV}$ et $E_V = 3{,}07~\text{eV} > 2{,}8~\text{eV}$. Seul le pointeur \emph{violet} déclenche le détecteur, même si le laser rouge est plus puissant : c'est l'énergie de \emph{chaque} photon qui compte, pas le nombre de photons.
\end{exoresolu}

\begin{exercice}[À toi de jouer]
Une station de radio de Yaoundé émet à la fréquence $\nu = 94{,}5~\text{MHz}$ ; une lampe germicide utilisée dans certains hôpitaux émet des UV de longueur d'onde $\lambda = 254~\text{nm}$.
\begin{enumerate}
\item Calculer l'énergie du photon radio, en joules et en eV.
\item Calculer l'énergie du photon UV, en joules et en eV.
\item Comparer ces énergies et expliquer pourquoi les UV germicides détruisent les microbes alors que les ondes radio nous traversent sans effet.
\end{enumerate}
\end{exercice}

\begin{corrige}[Exercice : photon radio contre photon UV]
\textbf{1. Photon radio.} $E = h\nu = \num{6,63e-34} \times \num{94,5e6} = \num{6,27e-26}~\text{J} \approx \num{3,9e-7}~\text{eV}$.

\textbf{2. Photon UV.} $E = \dfrac{hc}{\lambda} = \dfrac{\num{6,63e-34} \times \num{3,00e8}}{\num{254e-9}} = \num{7,83e-19}~\text{J} \approx 4{,}89~\text{eV}$.

\textbf{3. Comparaison.} Le photon UV est environ $10^{7}$ fois plus énergétique que le photon radio. À près de $5~\text{eV}$, il peut rompre des liaisons chimiques et détruire l'ADN des bactéries et virus ; le photon radio, avec moins d'un millionième d'électronvolt, ne peut provoquer aucune modification chimique : il traverse la matière vivante sans la perturber.
\end{corrige}

\textbf{En résumé de cette section.} La lumière est constituée de photons, grains d'énergie sans masse ni charge, se déplaçant à la célérité $c$ dans le vide. L'énergie d'un photon est donnée par la relation fondamentale $E = h\nu$, ou de façon équivalente $E = \dfrac{hc}{\lambda}$, avec $h = \num{6,63e-34}~\text{J}\cdot\text{s}$. Plus la longueur d'onde est petite, plus le photon est énergétique : c'est pourquoi les UV sont dangereux et les ondes radio inoffensives. Reste maintenant à comprendre \emph{comment} ces photons sont échangés avec les niveaux d'énergie quantifiés de la matière : ce sera l'objet de la section suivante, consacrée à la relation $\Delta E = h\nu$.

\coursec{Échanges d'énergie lumière-matière : $\Delta E = h\nu$}

Dans la section précédente, nous avons découvert le \cle{photon}, ce « grain » de lumière sans masse qui transporte une énergie $E = h\nu$, d'autant plus grande que la fréquence de la radiation est élevée. Nous avons aussi vu, en section 2, que l'énergie d'un atome est \cle{quantifiée} : elle ne peut prendre que certaines valeurs bien définies $E_1, E_2, E_3, \ldots$ Reste à réunir ces deux idées : \emph{comment} un photon et un atome échangent-ils de l'énergie ? La réponse est d'une élégance remarquable — et c'est elle qui explique pourquoi chaque élément chimique possède sa propre « carte d'identité » lumineuse.

\courssub{Le principe fondamental : un échange « sur mesure »}

\textbf{Une analogie pour comprendre.}\par

Imagine un escalier dans une maison de Yaoundé. Tu peux poser le pied sur la première marche, la deuxième, la troisième\ldots{} mais jamais \emph{entre} deux marches. Pour monter de la marche 1 à la marche 3, tu dois fournir \emph{exactement} l'effort correspondant à la différence de hauteur : un effort trop faible ne te fait pas monter, un effort trop grand est inutilisable tel quel. L'atome fonctionne de la même façon : ses niveaux d'énergie sont les « marches », et le photon est le « paquet d'énergie » qui permet de changer de marche. Mais attention : le paquet doit contenir \emph{exactement} la bonne quantité d'énergie, ni plus, ni moins.

\begin{propriete}[Condition d'échange d'énergie entre un photon et un atome (admise)]
Un atome ne peut absorber ou émettre un photon que si l'énergie de ce photon correspond \textbf{exactement} à la différence d'énergie entre deux de ses niveaux. Si $E_n$ et $E_p$ sont les énergies de deux niveaux (avec $E_p > E_n$), alors :
\[
\Delta E = E_p - E_n = h\nu = \dfrac{hc}{\lambda}
\]
où $h = \num{6,63e-34}\ \mathrm{J \cdot s}$ est la constante de Planck, $\nu$ la fréquence et $\lambda$ la longueur d'onde du photon absorbé ou émis.
\end{propriete}

Cette relation est la conséquence directe de deux principes que tu connais déjà :
\begin{itemize}
\item la \textbf{conservation de l'énergie} : l'énergie gagnée (ou perdue) par l'atome est fournie (ou emportée) par le photon ;
\item l'\textbf{énergie du photon} : $E = h\nu$ (section précédente).
\end{itemize}

\begin{attention}[La condition est stricte : « exactement »]
Un photon dont l'énergie $h\nu$ ne correspond à \textbf{aucune} transition possible de l'atome n'est \textbf{pas absorbé} : il traverse la matière ou est simplement diffusé. Il n'existe pas d'absorption « partielle » : un photon est absorbé entièrement ou pas du tout. C'est cette sélectivité qui rend la spectroscopie si précise pour identifier les éléments chimiques.
\end{attention}

\courssub{L'absorption : l'atome monte d'un niveau}

\begin{definition}[Absorption d'un photon]
L'\cle{absorption} est le phénomène par lequel un atome, initialement sur un niveau d'énergie $E_n$, \textbf{capture un photon} d'énergie $h\nu = E_p - E_n$ et passe sur le niveau supérieur $E_p$ (avec $E_p > E_n$). Le photon disparaît ; son énergie est entièrement transférée à l'atome, qui se trouve alors dans un \textbf{état excité}.
\end{definition}

Le bilan énergétique s'écrit simplement :
\[
\underbrace{E_n}_{\text{atome avant}} + \underbrace{h\nu}_{\text{photon absorbé}} = \underbrace{E_p}_{\text{atome après}}
\]

\begin{exemplebox}[Le sodium absorbe le jaune]
L'atome de sodium possède deux niveaux séparés de $\Delta E = \num{2,11}\ \mathrm{eV}$. Un photon jaune de longueur d'onde $\lambda = \num{589}\ \mathrm{nm}$ transporte une énergie :
\[
E = \dfrac{hc}{\lambda} = \dfrac{\num{6,63e-34} \times \num{3,00e8}}{\num{589e-9}} = \num{3,38e-19}\ \mathrm{J} = \dfrac{\num{3,38e-19}}{\num{1,60e-19}} \approx \num{2,11}\ \mathrm{eV}
\]
L'énergie du photon correspond exactement à la transition : le photon est absorbé. C'est pourquoi la vapeur de sodium éclairée en lumière blanche « pompe » précisément la composante jaune du spectre.
\end{exemplebox}

\courssub{L'émission : l'atome redescend en émettant un photon}

Un atome excité ne reste pas longtemps sur son niveau supérieur : l'état excité est \emph{instable} (durée de vie typique de l'ordre de $10^{-8}\ \mathrm{s}$). Il « désexcite » spontanément vers un niveau d'énergie inférieure.

\begin{definition}[Émission d'un photon]
L'\cle{émission} est le phénomène par lequel un atome excité, sur un niveau d'énergie $E_n$, redescend sur un niveau inférieur $E_p$ ($E_p < E_n$) en \textbf{émettant un photon} qui emporte l'énergie excédentaire :
\[
h\nu = E_n - E_p
\]
\end{definition}

Le bilan énergétique est l'inverse de celui de l'absorption :
\[
\underbrace{E_n}_{\text{atome excité}} = \underbrace{E_p}_{\text{atome après}} + \underbrace{h\nu}_{\text{photon émis}}
\]

C'est exactement ce qui se passe dans les lampes à vapeur de sodium qui éclairent en jaune certains axes routiers, ou dans les enseignes lumineuses au néon des quartiers de Douala : les atomes, excités par des décharges électriques, se désexcitent en émettant des photons de longueurs d'onde bien précises.

\begin{popfigure}
\begin{center}
\begin{tikzpicture}[scale=1.0, >=Stealth]
% Niveaux d'énergie
\draw[very thick, popblue] (0,0) -- (4,0) node[right, popdark] {$E_n$ (niveau inférieur)};
\draw[very thick, poporange] (0,3) -- (4,3) node[right, popdark] {$E_p$ (niveau supérieur)};
\draw[dashed, popmuted] (0,1.5) -- (4,1.5);
\node[popmuted, right] at (4,1.5) {\small niveaux intermédiaires};
% Absorption (à gauche)
\draw[->, very thick, popgreen] (0.8,0.08) -- (0.8,2.92);
\node[popgreen, left, align=center] at (0.7,1.5) {\small \textbf{Absorption}};
% Photon absorbé (flèche ondulée entrante)
\draw[->, thick, poppurple, decorate, decoration={snake, amplitude=0.8mm, segment length=4mm}] (-2.2,0.6) -- (0.7,0.35);
\node[poppurple, left] at (-2.2,0.6) {\small photon $h\nu$};
% Émission (à droite)
\draw[->, very thick, poppink] (3.2,2.92) -- (3.2,0.08);
\node[poppink, right, align=center] at (3.3,1.5) {\small \textbf{Émission}};
% Photon émis (flèche ondulée sortante)
\draw[->, thick, poppurple, decorate, decoration={snake, amplitude=0.8mm, segment length=4mm}] (3.3,0.35) -- (6.2,0.6);
\node[poppurple, right] at (6.2,0.6) {\small photon $h\nu$};
% Delta E
\draw[<->, thick, popdark] (2,0) -- (2,3) node[midway, right, popdark] {\small $\Delta E = E_p - E_n = h\nu$};
\end{tikzpicture}
\end{center}
\legende{Absorption : l'atome capture un photon d'énergie $h\nu = E_p - E_n$ et monte du niveau $E_n$ au niveau $E_p$. Émission : l'atome excité redescend de $E_p$ vers $E_n$ en émettant un photon de même énergie. Les deux phénomènes mettent en jeu la \emph{même} longueur d'onde.}
\end{popfigure}

\courssub{Conséquence spectrale : la signature de chaque élément}

Puisque les niveaux d'énergie d'un atome sont \textbf{quantifiés} et \textbf{discrets}, les différences $E_p - E_n$ ne peuvent prendre qu'un nombre limité de valeurs. Par conséquent, les photons absorbés ou émis n'ont eux aussi que certaines longueurs d'onde bien déterminées :
\[
\lambda = \dfrac{hc}{E_p - E_n}
\]
C'est l'explication de la \textbf{discontinuité des spectres de raies} : au lieu d'un spectre continu (comme celui d'un corps chaud), un gaz d'atomes émet ou absorbe uniquement des \emph{raies} isolées.

\begin{propriete}[Signature spectrale d'un élément]
Les raies d'\textbf{émission} et les raies d'\textbf{absorption} d'un même élément correspondent aux \textbf{mêmes transitions} entre niveaux : elles ont donc les \textbf{mêmes longueurs d'onde}. Chaque élément chimique possède ainsi un spectre de raies qui lui est propre : c'est sa \cle{signature spectrale}, aussi unique qu'une empreinte digitale.
\end{propriete}

\begin{popfigure}
\begin{center}
\begin{tikzpicture}[scale=1.0]
% Spectre d'émission : fond noir, raies brillantes
\fill[black] (0,2.2) rectangle (10,3.4);
\foreach \x/\col in {1.6/poppurple, 2.3/popblue, 3.1/popgreen, 5.2/popgold, 6.4/poporange, 7.8/poppink}{
  \fill[\col] (\x,2.2) rectangle (\x+0.12,3.4);
}
\node[popdark, left, align=right] at (-0.2,2.8) {\small \textbf{Émission}\\ \small (raies brillantes)};
% Spectre d'absorption : fond continu (dégradé simulé), raies noires
\shade[left color=poppurple, right color=poppink] (0,0) rectangle (10,1.2);
\foreach \x in {1.6, 2.3, 3.1, 5.2, 6.4, 7.8}{
  \fill[black] (\x,0) rectangle (\x+0.12,1.2);
}
\node[popdark, left, align=right] at (-0.2,0.6) {\small \textbf{Absorption}\\ \small (raies noires)};
% Axe longueur d'onde
\draw[->, thick, popdark] (0,-0.3) -- (10.3,-0.3);
\node[popdark, right] at (10.3,-0.3) {\small $\lambda$};
\node[popdark, below] at (1,-0.3) {\small 400 nm};
\node[popdark, below] at (9,-0.3) {\small 750 nm};
% Correspondances
\foreach \x in {1.6, 2.3, 3.1, 5.2, 6.4, 7.8}{
  \draw[dashed, popmuted, thin] (\x+0.06,2.2) -- (\x+0.06,1.2);
}
\end{tikzpicture}
\end{center}
\legende{Spectres du même élément. En haut : spectre d'émission — raies brillantes colorées sur fond noir. En bas : spectre d'absorption — raies noires aux \emph{mêmes positions} sur le fond continu de la lumière blanche. Les pointillés soulignent la correspondance exacte des longueurs d'onde.}
\end{popfigure}

\begin{methode}[Exploiter un diagramme de niveaux d'énergie]
Pour déterminer la radiation absorbée ou émise lors d'une transition :
\begin{enumerate}
\item \textbf{Identifier} les deux niveaux de la transition sur le diagramme ($E_n$ et $E_p$).
\item \textbf{Calculer} la différence d'énergie : $\Delta E = |E_n - E_p|$ (toujours positive ; convertir en joules si les niveaux sont en eV : $1\ \mathrm{eV} = \num{1,60e-19}\ \mathrm{J}$).
\item \textbf{En déduire} la fréquence $\nu = \dfrac{\Delta E}{h}$ puis la longueur d'onde $\lambda = \dfrac{c}{\nu} = \dfrac{hc}{\Delta E}$.
\item \textbf{Vérifier} le domaine spectral : UV si $\lambda < 400\ \mathrm{nm}$, visible si $400\ \mathrm{nm} < \lambda < 750\ \mathrm{nm}$ (environ), IR si $\lambda > 750\ \mathrm{nm}$.
\end{enumerate}
Astuce de calcul : avec $\Delta E$ en eV et $\lambda$ en nm, on peut utiliser directement $\lambda\ (\mathrm{nm}) \approx \dfrac{1240}{\Delta E\ (\mathrm{eV})}$, car $hc \approx \num{1,24e-6}\ \mathrm{eV \cdot m}$ (valeur usuelle au Probatoire).
\end{methode}

\begin{attention}[Pièges fréquents au Probatoire]
\begin{itemize}
\item \textbf{Oublier la conversion eV $\rightarrow$ J} avant d'appliquer $\lambda = hc/\Delta E$ avec $h$ en $\mathrm{J \cdot s}$. C'est l'erreur numéro un !
\item \textbf{Signe de $\Delta E$} : l'énergie du photon est toujours \emph{positive}. Pour l'émission $E_n \to E_p$ ($E_n > E_p$), on écrit $h\nu = E_n - E_p$ ; pour l'absorption, $h\nu = E_p - E_n$. En cas de doute, prends la valeur absolue.
\item \textbf{Confondre fréquence et longueur d'onde} : $\lambda = c/\nu$, donc grande fréquence $\Leftrightarrow$ petite longueur d'onde $\Leftrightarrow$ photon très énergétique.
\end{itemize}
\end{attention}

\courssub{Application : l'identification chimique par spectroscopie}

La spectroscopie est l'une des techniques d'analyse les plus puissantes de la chimie moderne : en dispersant la lumière émise par une lampe à vapeur (sodium, mercure, hydrogène\ldots) ou la lumière traversant un gaz, on observe les raies et on en déduit \emph{sans ambiguïté} les éléments présents. C'est ainsi que l'on analyse la composition des étoiles, que l'on contrôle la pureté des matériaux, ou que l'on détecte des traces de substances dans les laboratoires hospitaliers.

\begin{savaistu}[L'hélium : un élément découvert dans le Soleil avant la Terre]
En 1868, lors d'une éclipse totale de Soleil, l'astronome français Jules Janssen observe une raie jaune brillante dans le spectre de la couronne solaire, à $\lambda = \num{587,6}\ \mathrm{nm}$. Or cette raie ne correspond à \emph{aucun} élément connu sur Terre ! L'astronome Lockyer en conclut qu'il s'agit d'un élément nouveau, qu'il baptise \textbf{hélium}, du grec \emph{hélios} (Soleil). Il faudra attendre 1895 pour que le chimiste Ramsay l'identifie sur Terre, dans un minerai d'uranium. C'est la preuve éclatante que la spectroscopie permet d'analyser la matière à des millions de kilomètres de distance.
\end{savaistu}

\courssub{Exercice guidé : la raie rouge H$\alpha$ de l'hydrogène}

\begin{exoresolu}[La raie H$\alpha$ de l'hydrogène]
\textbf{Énoncé.} Les premiers niveaux d'énergie de l'atome d'hydrogène sont donnés par $E_n = -\dfrac{\num{13,6}}{n^2}\ \mathrm{eV}$. On donne : $h = \num{6,63e-34}\ \mathrm{J \cdot s}$ ; $c = \num{3,00e8}\ \mathrm{m \cdot s^{-1}}$ ; $1\ \mathrm{eV} = \num{1,60e-19}\ \mathrm{J}$.
\begin{enumerate}
\item Calculer les énergies des niveaux $n = 2$ et $n = 3$.
\item Un atome d'hydrogène excité sur le niveau $n = 3$ se désexcite vers le niveau $n = 2$. Calculer l'énergie du photon émis, en eV puis en joules.
\item En déduire la longueur d'onde $\lambda$ du photon émis. Ce photon de $\lambda \approx \num{658}\ \mathrm{nm}$ peut-il effectivement être émis par l'hydrogène ? À quel domaine spectral appartient-il ?
\end{enumerate}
\tcblower
\textbf{Solution détaillée.}
\begin{enumerate}
\item Niveaux d'énergie :
\[
E_2 = -\dfrac{\num{13,6}}{2^2} = -\dfrac{\num{13,6}}{4} = \num{-3,40}\ \mathrm{eV} \qquad ; \qquad E_3 = -\dfrac{\num{13,6}}{3^2} = -\dfrac{\num{13,6}}{9} \approx \num{-1,51}\ \mathrm{eV}
\]
\item Énergie du photon émis lors de la transition $E_3 \to E_2$ :
\[
\Delta E = E_3 - E_2 = \num{-1,51} - (\num{-3,40}) = \num{1,89}\ \mathrm{eV}
\]
Conversion en joules :
\[
\Delta E = \num{1,89} \times \num{1,60e-19} = \num{3,02e-19}\ \mathrm{J}
\]
\item Longueur d'onde du photon :
\[
\lambda = \dfrac{hc}{\Delta E} = \dfrac{\num{6,63e-34} \times \num{3,00e8}}{\num{3,02e-19}} \approx \num{6,59e-7}\ \mathrm{m} = \num{659}\ \mathrm{nm}
\]
Vérification avec la formule rapide : $\lambda \approx \dfrac{1240}{\num{1,89}} \approx \num{656}\ \mathrm{nm}$ (valeur de référence). La légère différence provient des valeurs arrondies des constantes imposées. \textbf{Conclusion} : oui, un photon de $\lambda \approx \num{658}\ \mathrm{nm}$ (valeur de référence $\num{656}\ \mathrm{nm}$) peut être émis par l'hydrogène — c'est exactement la transition $E_3 \to E_2$. Comme $400\ \mathrm{nm} < 658\ \mathrm{nm} < 750\ \mathrm{nm}$, cette radiation appartient au \textbf{domaine visible}, dans le \textbf{rouge} : c'est la célèbre raie \textbf{H$\alpha$}, la raie la plus intense du spectre visible de l'hydrogène, celle qui donne sa teinte rosée aux nébuleuses photographiées par les télescopes.
\end{enumerate}
\end{exoresolu}

\begin{aretenir}[L'essentiel de la section]
\begin{itemize}
\item Un atome n'absorbe ou n'émet un photon que si son énergie correspond \textbf{exactement} à une transition : $\Delta E = h\nu = \dfrac{hc}{\lambda}$.
\item \textbf{Absorption} : $h\nu = E_p - E_n$ (l'atome monte, le photon disparaît). \textbf{Émission} : $h\nu = E_n - E_p$ (l'atome redescend, le photon est créé).
\item La quantification des niveaux explique la \textbf{discontinuité} des spectres de raies.
\item Raies d'émission et d'absorption d'un même élément : \textbf{mêmes longueurs d'onde} — chaque élément a sa signature spectrale, base de l'identification chimique par spectroscopie.
\item Raie H$\alpha$ de l'hydrogène : $\Delta E = \num{1,89}\ \mathrm{eV}$, $\lambda \approx \num{656}\ \mathrm{nm}$ (rouge).
\end{itemize}
\end{aretenir}

\coursec{Profil spectral d'une source lumineuse}

Dans la section précédente, nous avons établi que les échanges d'énergie entre la lumière et la matière sont \emph{quantifiés} : un atome absorbe ou émet une énergie $\Delta E = h\nu$ sous la forme d'un unique photon. Si nous observons maintenant une source macroscopique (une lampe, le Soleil, un métal chauffé), elle contient un nombre immense d'atomes qui émettent collectivement. La répartition de l'énergie lumineuse émise en fonction de la longueur d'onde (ou de la fréquence) constitue ce que l'on appelle le \textbf{profil spectral}. C'est l'empreinte digitale de la source, qui nous renseigne sur sa nature, sa température et les mécanismes physiques à l'œuvre.

\courssub{Définition et interprétation du profil spectral}

\begin{definition}[Profil spectral d'une source lumineuse]
Le \textbf{profil spectral} (ou distribution spectrale de l'éclairement, de la luminance ou de la puissance rayonnée) est la courbe qui donne l'intensité lumineuse $I$ (ou la puissance spectrale) émise par une source en fonction de la longueur d'onde $\lambda$ (ou de la fréquence $\nu$) : $I = f(\lambda)$.
\end{definition}

Pour un \textbf{corps chaud} (filament d'une lampe à incandescence, surface du Soleil, métal porté au rouge par un forgeron), l'émission est due à l'agitation thermique des charges (électrons, ions) dans la matière. Cette agitation étant continue et aléatoire, le rayonnement émis couvre un domaine continu de longueurs d'onde : on parle de \textbf{spectre continu}. La forme de ce spectre est une \textbf{courbe en cloche} (asymétrique) qui présente un \textbf{maximum d'émission} à une longueur d'onde caractéristique notée $\lambda_{\text{max}}$.

\begin{center}
\begin{popfigure}
\begin{tikzpicture}
\begin{axis}[
    width=\linewidth, height=0.55\linewidth,
    axis lines=middle,
    xlabel={$\lambda$ (nm)}, ylabel={$I(\lambda)$ (u.a.)},
    xmin=100, xmax=2500,
    ymin=0, ymax=1.1,
    xtick={400,500,600,700,1000,1500,2000},
    ytick=\empty,
    xlabel style={anchor=north west}, ylabel style={anchor=south east},
    domain=100:2500, samples=400,
    clip=false,
    enlargelimits=0.02,
    legend style={at={(0.98,0.98)}, anchor=north east, font=\small, fill=white, draw=popmuted}
]
% Universal Planck function normalized by T^5: y = 1/(x*T)^5 * 1/(exp(1.4388e7/(x*T))-1)
% x in nm, T in K. Constant hc/k = 1.4388e7 nm K.
% Scaling factor 3e-11 makes peak ~1.
% T = 3000 K (poporange) -> peak at 966 nm
\addplot[poporange, thick, domain=150:2500, samples=300] 
    ({x}, {3e-11 * (1/(x*3000)^5) * (1/(exp(1.4388e7/(x*3000))-1))});
\addlegendentry{$T = 3000\,\text{K}$ (Lampe incandescente)}

% T = 4500 K (popgold) -> peak at 644 nm
\addplot[popgold, thick, domain=150:2500, samples=300] 
    ({x}, {3e-11 * (1/(x*4500)^5) * (1/(exp(1.4388e7/(x*4500))-1))});
\addlegendentry{$T = 4500\,\text{K}$}

% T = 6000 K (popblue) -> peak at 483 nm
\addplot[popblue, thick, domain=150:2500, samples=300] 
    ({x}, {3e-11 * (1/(x*6000)^5) * (1/(exp(1.4388e7/(x*6000))-1))});
\addlegendentry{$T = 6000\,\text{K}$ (Soleil)}

% Vertical lines for lambda_max (Wien: lambda_max = 2.898e6 / T nm)
% 3000K -> 966 nm
% 4500K -> 644 nm
% 6000K -> 483 nm
\draw[poporange, dashed, thick] (966,0) -- (966,0.95) node[above, poporange, font=\small] {$\lambda_{\text{max}}$};
\draw[popgold, dashed, thick] (644,0) -- (644,0.98) node[above, popgold, font=\small] {$\lambda_{\text{max}}$};
\draw[popblue, dashed, thick] (483,0) -- (483,1.0) node[above, popblue, font=\small] {$\lambda_{\text{max}}$};

% Visible domain shading
\fill[popgreen!15] (380,0) rectangle (780,1.1);
\node[popgreen!70!black, font=\small, rotate=90] at (370,0.5) {Domaine visible};
\end{axis}
\end{tikzpicture}
\legende{Profils spectraux de corps noirs (approximation corps chaud) à trois températures. Le maximum d'émission $\lambda_{\text{max}}$ se décale vers les courtes longueurs d'onde (le bleu) quand la température augmente (loi de Wien). La zone verte hachurée représente le domaine visible (380--780 nm). Les courbes sont normalisées pour avoir la même hauteur maximale, ce qui met en évidence le déplacement du pic.}
\end{popfigure}
\end{center}

\courssub{Influence de la température : loi de Wien}

L'observation des courbes de la figure ci-dessus révèle deux faits majeurs, découverts expérimentalement à la fin du XIX\textsuperscript{e} siècle (Lummer, Pringsheim, Wien) avant d'être expliqués par Planck :

\begin{enumerate}
    \item \textbf{Déplacement du maximum :} Lorsque la température $T$ du corps chaud augmente, la longueur d'onde $\lambda_{\text{max}}$ correspondant au maximum d'émission \textbf{diminue}. Le pic du spectre se déplace vers les courtes longueurs d'onde (vers le bleu, l'ultraviolet).
    \item \textbf{Augmentation de l'intensité totale :} La surface sous la courbe (énergie totale émise par unité de surface et par seconde, c'est-à-dire la luminance énergétique) augmente très rapidement avec la température (loi de Stefan-Boltzmann : $P \propto T^4$).
\end{enumerate}

Le lien quantitatif entre $\lambda_{\text{max}}$ et $T$ est donné par la \textbf{loi de Wien} (admise au programme).

\begin{propriete}[Loi de Wien — Déplacement du maximum spectral (admise)]
Pour un corps noir (ou un corps chaud se comportant comme un corps noir), le produit de la longueur d'onde $\lambda_{\text{max}}$ (en mètres) au maximum d'émission par la température thermodynamique $T$ (en kelvins) est une constante :
\[
\boxed{\lambda_{\text{max}} \times T = 2,90 \times 10^{-3} \ \text{m}\cdot\text{K}}
\]
\textbf{Analyse dimensionnelle :} $[\lambda_{\text{max}}] = \text{m}$, $[T] = \text{K}$, donc $[\lambda_{\text{max}} T] = \text{m}\cdot\text{K}$. La constante $b = 2,90 \times 10^{-3} \ \text{m}\cdot\text{K}$ est homogène.
\end{propriete}

\textbf{Conséquence visible : la couleur de la température.}
Un forgeron à Douala ou à Garoua n'a pas besoin de thermomètre pour connaître la température de son fer : il \emph{regarde} sa couleur.
\begin{itemize}
    \item \textbf{Fer rouge sombre} ($\approx 800\,\text{K}$ à $1000\,\text{K}$) : $\lambda_{\text{max}} \approx 2900\,\text{nm}$ à $3600\,\text{nm}$ (infrarouge). Le pic est dans l'infrarouge, mais la queue de la courbe dans le rouge visible donne cette teinte.
    \item \textbf{Fer orange/jaune} ($\approx 1500\,\text{K}$ à $2000\,\text{K}$) : $\lambda_{\text{max}}$ rentre dans le visible (rouge/orange).
    \item \textbf{Fer blanc} ($> 2500\,\text{K}$) : $\lambda_{\text{max}}$ est au cœur du visible (vert/bleu), toutes les couleurs visibles sont émises fortement $\rightarrow$ lumière blanche.
\end{itemize}
C'est la même physique qui dicte la couleur des étoiles : \textbf{les étoiles bleues sont plus chaudes que les étoiles rouges}.

\begin{methode}[Déterminer la température d'un corps chaud par analyse spectrale]
\begin{enumerate}
    \item \textbf{Relever} la longueur d'onde $\lambda_{\text{max}}$ (en nm le plus souvent) sur le profil spectral expérimental ou donné.
    \item \textbf{Convertir} $\lambda_{\text{max}}$ en \textbf{mètres} ($1\,\text{nm} = 10^{-9}\,\text{m}$).
    \item \textbf{Appliquer} la loi de Wien : $T = \dfrac{2,90 \times 10^{-3}}{\lambda_{\text{max}}}$ (résultat en kelvins).
    \item \textbf{Convertir} éventuellement en degrés Celsius : $\theta(^\circ\text{C}) = T(\text{K}) - 273$.
\end{enumerate}
\end{methode}

\begin{exoresolu}[Température de surface du Soleil]
Le profil spectral du Soleil présente un maximum d'émission pour $\lambda_{\text{max}} \approx 500\,\text{nm}$ (vert/bleu). Déterminer la température de surface du Soleil (photosphère).
\tcblower
\textbf{Solution détaillée.}
\begin{enumerate}
    \item $\lambda_{\text{max}} = 500\,\text{nm} = 500 \times 10^{-9}\,\text{m} = 5,00 \times 10^{-7}\,\text{m}$.
    \item Loi de Wien : $\lambda_{\text{max}} T = 2,90 \times 10^{-3}\,\text{m}\cdot\text{K}$.
    \item $T = \dfrac{2,90 \times 10^{-3}}{5,00 \times 10^{-7}} = \dfrac{2,90}{5,00} \times 10^{4} = 0,58 \times 10^{4}\,\text{K} = 5\,800\,\text{K}$.
    \item En Celsius : $\theta = 5\,800 - 273 = 5\,527\,^\circ\text{C}$.
\end{enumerate}
\textbf{Réponse :} La température de surface du Soleil est d'environ \textbf{$5\,800\,\text{K}$} ($5\,527\,^\circ\text{C}$).
\end{exoresolu}

\begin{attention}[Piège classique : Celsius ou Kelvin ?]
La loi de Wien \textbf{exige impérativement} la température en \textbf{kelvins (K)}. Ne jamais insérer une température en degrés Celsius ($^\circ\text{C}$) dans la formule $\lambda_{\text{max}} T = \text{cste}$.
\begin{itemize}
    \item \textbf{Faux :} $T = 500\,^\circ\text{C} \Rightarrow \lambda_{\text{max}} = \frac{2,90 \times 10^{-3}}{500} = 5,8 \times 10^{-6}\,\text{m}$ (Erreur !)
    \item \textbf{Vrai :} $T = 500 + 273 = 773\,\text{K} \Rightarrow \lambda_{\text{max}} = \frac{2,90 \times 10^{-3}}{773} \approx 3,75 \times 10^{-6}\,\text{m}$.
\end{itemize}
Rappel : $0\,\text{K} = -273,15\,^\circ\text{C}$. La conversion est $\theta(^\circ\text{C}) = T(\text{K}) - 273$ (valeur approchée suffisante au Probatoire).
\end{attention}

\courssub{Profils spectraux des sources modernes et efficacité énergétique}

Tous les corps chauds (bougie, lampe à incandescence, Soleil) ont un profil spectral \emph{continu} en cloche. Les sources modernes (LED, tubes fluorescents, lampes à décharge) fonctionnent sur des principes différents (électroluminescence, fluorescence) et leurs profils spectraux sont radicalement différents : ils présentent des \textbf{pics étroits} (raies d'émission) ou des bandes étroites, superposés parfois à un fond continu.

\begin{center}
\begin{popfigure}
\begin{tikzpicture}
\begin{axis}[
    width=\linewidth, height=0.55\linewidth,
    axis lines=middle,
    xlabel={$\lambda$ (nm)}, ylabel={$I(\lambda)$ (u.a.)},
    xmin=300, xmax=800,
    ymin=0, ymax=1.1,
    xtick={350,400,450,500,550,600,650,700,750},
    ytick=\empty,
    xlabel style={anchor=north west}, ylabel style={anchor=south east},
    domain=300:800, samples=500,
    clip=false,
    enlargelimits=0.02,
    legend style={at={(0.98,0.98)}, anchor=north east, font=\small, fill=white, draw=popmuted}
]
% Incandescent ~ 2700K (Blackbody) - Raw Planck (not normalized by T^5) to show low visible intensity relative to IR peak
% B_lambda ~ 1/lambda^5 * 1/(exp(C2/lambda T)-1), C2=1.4388e7 nm K
% Scaling 2.5e-28 for visibility on 0-1 scale (les valeurs brutes de la loi de Planck sont énormes en unités SI).
\addplot[poporange, thick, domain=300:800] 
    ({x}, {2.5e-28 * (1/(x*1e-9)^5) * (1/(exp(1.4388e-2/(x*1e-9*2700))-1))});
\addlegendentry{Lampe à incandescence ($T\approx 2700\,\text{K}$)}

% LED blanche typique : Pic bleu ~450nm (puce InGaN) + Pic large jaune ~560nm (phosphore YAG:Ce)
% Gaussian peaks
\def\gauss#1#2#3{exp(-(x-#1)^2/(2*#2^2))*#3}
% Blue peak: center 450, sigma 15, height 1.0
\addplot[popblue, thick, domain=300:800] 
    ({x}, {\gauss{450}{15}{1.0}});
% Yellow peak: center 560, sigma 60, height 0.8
\addplot[popblue, thick, domain=300:800] 
    ({x}, {\gauss{560}{60}{0.8}});
\addlegendentry{LED blanche (puce bleue + phosphore)}

% Visible domain shading
\fill[popgreen!10] (380,0) rectangle (780,1.1);
\end{axis}
\end{tikzpicture}
\legende{Comparaison des profils spectraux : lampe à incandescence (courbe continue orange, corps noir à 2700 K, maximum dans l'infrarouge) et LED blanche (deux pics étroits : bleu $\approx 450\,\text{nm}$ et jaune large $\approx 560\,\text{nm}$). La LED concentre son énergie dans le visible.}
\end{popfigure}
\end{center}

\begin{center}
\begin{tabularx}{\linewidth}{|l|X|X|X|}
\hline
\rowcolor{popblueL}
\textbf{Source} & \textbf{Type de spectre} & \textbf{Aspect du profil $I(\lambda)$} & \textbf{Efficacité lumineuse (lm/W)} \\
\hline
Bougie & Continu (corps chaud $\approx 1800\,\text{K}$) & Cloche très décalée vers IR (rouge/orange) & $\approx 0,3$ (très faible) \\
\hline
Lampe à incandescence & Continu (corps chaud $\approx 2700\,\text{K}$) & Cloche décalée vers IR (blanc chaud) & $10$ à $17$ (faible : 90\% en chaleur) \\
\hline
Soleil (surface) & Continu + raies d'absorption & Cloche centrée sur visible ($\lambda_{\text{max}}\approx 500\,\text{nm}$) & $\approx 93$ (référence) \\
\hline
Tube fluorescent & Raies d'émission (Hg) + Continu (phosphores) & Pic UV fort + pics visibles (bleu, vert, rouge) & $50$ à $100$ (moyenne) \\
\hline
LED blanche & Pics étroits (Bleu + Jaune/Vert/Rouge) & Deux bandes principales étroites & $100$ à $200+$ (élevée, peu de chaleur) \\
\hline
\end{tabularx}
\end{center}

\paragraph{Interprétation physique et enjeu énergétique.}
\begin{itemize}
    \item \textbf{Incandescence :} L'électricité chauffe un filament (effet Joule). La majeure partie de l'énergie est rayonnée dans l'infrarouge (chaleur) invisible pour l'œil, d'où le faible rendement lumineux. C'est le principe du \emph{corps noir}.
    \item \textbf{LED (Light Emitting Diode) :} L'émission provient de la \textbf{recombinaison radiative} d'électrons et de trous dans une jonction semi-conductrice (GaN, InGaN). L'énergie du photon émis est proche de la bande interdite $E_g$ : $h\nu \approx E_g$. Cela donne un pic \textbf{très étroit} (quasi monochromatique). Pour faire du blanc, on utilise une puce bleue ($E_g \approx 2,7\,\text{eV}$, $\lambda \approx 450\,\text{nm}$) qui excite un phosphore jaune (ou vert+rouge). Le profil spectral montre donc un pic bleu résiduel et une bande jaune large.
    \item \textbf{Efficacité :} Comme la LED émet essentiellement dans le domaine visible (là où l'œil est sensible), son \textbf{efficacité lumineuse} (flux lumineux en lumens par watt électrique consommé) est 10 à 15 fois supérieure à l'incandescence. C'est un enjeu majeur pour l'électrification rurale au Cameroun (kits solaires + LED) et la réduction de la facture ENEO.
\end{itemize}

\begin{savaistu}[Code couleur des robinets vs Température des étoiles]
Dans la vie courante (robinets, brûleurs de gaz), le \textbf{bleu} code pour le \textbf{froid} et le \textbf{rouge} pour le \textbf{chaud}. En astrophysique et en physique du corps noir, c'est \textbf{l'inverse} !
\begin{itemize}
    \item Étoile \textbf{rouge} (ex: Bételgeuse) : $T \approx 3\,500\,\text{K}$ (relativement \textbf{froide}).
    \item Étoile \textbf{jaune} (ex: Soleil) : $T \approx 5\,800\,\text{K}$.
    \item Étoile \textbf{bleue} (ex: Rigel, Sirius A) : $T > 10\,000\,\text{K}$ (très \textbf{chaude}).
\end{itemize}
La loi de Wien ($\lambda_{\text{max}} T = \text{cste}$) l'explique : plus $T$ est grand, plus $\lambda_{\text{max}}$ est petit (vers le bleu/UV). Le code couleur domestique est une convention psychologique (feu = danger = rouge ; eau/glace = bleu), pas une loi physique.
\end{savaistu}

\courssub{Vers le spectre solaire : raies d'absorption}

Le profil spectral du Soleil n'est pas une courbe de corps noir \emph{parfaite}. Si l'on observe avec une haute résolution spectrale (spectrographe), on découvre une multitude de \textbf{raies sombres} (défauts d'intensité) superposées au continuum lumineux : ce sont les \textbf{raies de Fraunhofer}. Elles témoignent de l'absorption de photons de fréquences précises par les atomes de l'atmosphère solaire (photosphère et chromosphère) plus froide que les couches profondes. C'est l'application directe de la relation $\Delta E = h\nu$ vue précédemment : les atomes d'hydrogène, hélium, sodium, fer, calcium... absorbent sélectivement les photons dont l'énergie correspond à leurs sauts quantiques. Nous analyserons ce \textbf{spectre d'absorption} en détail dans la prochaine section.

\begin{exercice}[Application : Température d'une étoile et type spectral]
On observe le profil spectral de trois étoiles A, B, C. Leurs longueurs d'onde d'émission maximale sont :
\begin{itemize}
    \item Étoile A : $\lambda_{\text{max}} = 290\,\text{nm}$ (ultraviolet)
    \item Étoile B : $\lambda_{\text{max}} = 500\,\text{nm}$ (vert)
    \item Étoile C : $\lambda_{\text{max}} = 1\,200\,\text{nm}$ (infrarouge)
\end{itemize}
\begin{enumerate}
    \item Calculer la température de surface de chaque étoile en kelvins.
    \item Classer ces étoiles du plus chaud au plus froid.
    \item Donner la couleur apparente probable de chaque étoile (Rouge, Blanc, Bleu).
\end{enumerate}
\end{exercice}

\begin{corrige}[Exercice : Température d'une étoile et type spectral]
\begin{enumerate}
    \item \textbf{Loi de Wien :} $T = \frac{2,90 \times 10^{-3}}{\lambda_{\text{max}}}$ (avec $\lambda_{\text{max}}$ en m).
    \begin{itemize}
        \item \textbf{Étoile A :} $\lambda_{\text{max}} = 290\,\text{nm} = 2,90 \times 10^{-7}\,\text{m}$.
              $T_A = \frac{2,90 \times 10^{-3}}{2,90 \times 10^{-7}} = 10^4\,\text{K} = \textbf{10\,000 K}$.
        \item \textbf{Étoile B :} $\lambda_{\text{max}} = 500\,\text{nm} = 5,00 \times 10^{-7}\,\text{m}$.
              $T_B = \frac{2,90 \times 10^{-3}}{5,00 \times 10^{-7}} = 5,8 \times 10^3\,\text{K} = \textbf{5\,800 K}$.
        \item \textbf{Étoile C :} $\lambda_{\text{max}} = 1\,200\,\text{nm} = 1,20 \times 10^{-6}\,\text{m}$.
              $T_C = \frac{2,90 \times 10^{-3}}{1,20 \times 10^{-6}} \approx 2,42 \times 10^3\,\text{K} = \textbf{2\,420 K}$.
    \end{itemize}
    \item \textbf{Classement (chaud $\rightarrow$ froid) :} A ($10\,000\,\text{K}$) $>$ B ($5\,800\,\text{K}$) $>$ C ($2\,420\,\text{K}$).
    \item \textbf{Couleur apparente :}
    \begin{itemize}
        \item A ($10\,000\,\text{K}$, pic dans UV) : Rayonnement fort dans le bleu/violet visible $\rightarrow$ \textbf{Bleue}.
        \item B ($5\,800\,\text{K}$, pic dans vert) : Rayonnement équilibré sur tout le visible $\rightarrow$ \textbf{Blanche/Jaune} (type solaire).
        \item C ($2\,420\,\text{K}$, pic dans IR) : Queue de courbe dans le rouge/orange visible $\rightarrow$ \textbf{Rouge/Orange}.
    \end{itemize}
\end{enumerate}
\end{corrige}

\coursec{Le spectre solaire : forme et raies}

Dans la section précédente, nous avons appris à décrire le \cle{profil spectral} d'une source lumineuse : la répartition de l'intensité lumineuse en fonction de la longueur d'onde. Nous avons vu qu'un corps chauffé (filament d'une lampe, métal porté au rouge) émet un spectre \emph{continu} dont le maximum d'émission se déplace vers les courtes longueurs d'onde quand la température augmente (loi de Wien), tandis qu'un gaz excité émet un spectre de \emph{raies} caractéristiques de ses atomes. Que se passe-t-il alors lorsqu'on observe le Soleil, cette immense boule de gaz incandescent qui nous éclaire chaque jour à Douala comme à Maroua ? La réponse est fascinante : son spectre est \emph{les deux à la fois}, et c'est précisément cette double nature qui permet aux astronomes de connaître la température et la composition chimique des étoiles sans jamais y toucher.

\courssub{Observation : un fond continu traversé de raies noires}

\begin{experience}[Observation du spectre solaire]
\textbf{Matériel :} un prisme (ou un réseau de diffraction), une fente fine, un écran blanc.\\
\textbf{Protocole :} on envoie un fin pinceau de lumière solaire (jamais l'image directe du Soleil dans l'œil !) à travers la fente, puis à travers le prisme, et on recueille le spectre sur l'écran.\\
\textbf{Observations :}
\begin{itemize}
\item on obtient un étalement \emph{continu} de couleurs, du violet au rouge, qui se prolonge au-delà du visible : dans l'\cle{ultraviolet} (UV, $\lambda < 400\ \mathrm{nm}$) côté violet, et dans l'\cle{infrarouge} (IR, $\lambda > 800\ \mathrm{nm}$) côté rouge ;
\item en examinant attentivement avec un instrument de bonne résolution, on découvre que ce fond coloré est traversé de \emph{centaines de fines raies noires}, à des longueurs d'onde très précises.
\end{itemize}
\textbf{Interprétation préliminaire :} le fond continu est la signature d'un \emph{corps chaud dense} ; les raies noires sont des \emph{absorptions} : des photons de certaines énergies ont disparu entre le Soleil et nous.
\end{experience}

\begin{definition}[Raies de Fraunhofer]
Les \cle{raies de Fraunhofer} sont les raies d'\emph{absorption} (raies sombres) observées dans le spectre de la lumière solaire. Elles doivent leur nom à l'opticien allemand Joseph von Fraunhofer qui en catalogua plusieurs centaines dès 1814 et repéra les principales par des lettres : A, B, C, D, E, F, G, H, K. Chaque raie correspond à l'absorption de photons dont l'énergie $h\nu$ coïncide exactement avec une transition d'un atome présent dans l'atmosphère solaire.
\end{definition}

\begin{popfigure}
\begin{center}
\begin{tikzpicture}[x=1.6cm,y=0.055cm]
% fond continu (arc-en-ciel schématique)
\fill[poppurple!60] (0,0) rectangle (0.8,60);
\fill[popblue!60] (0.8,0) rectangle (1.6,60);
\fill[popgreen!60] (1.6,0) rectangle (2.4,60);
\fill[popgold!70] (2.4,0) rectangle (3.2,60);
\fill[poporange!70] (3.2,0) rectangle (4.0,60);
\fill[popink!70] (4.0,0) rectangle (4.8,60);
% raies noires
\fill[black] (1.05,0) rectangle (1.12,60); % H beta ~486
\fill[black] (2.95,0) rectangle (3.02,60); % D sodium ~589
\fill[black] (3.85,0) rectangle (3.92,60); % H alpha ~656
\fill[black] (1.55,0) rectangle (1.60,60); % Mg ~517
\fill[black] (2.35,0) rectangle (2.40,60); % Fe ~527
% annotations
\draw[<-,thick] (1.085,60) -- (1.085,78) node[above,font=\small] {H$\beta$ (486 nm)};
\draw[<-,thick] (1.575,60) -- (1.575,70) node[above,font=\small] {Mg};
\draw[<-,thick] (2.985,60) -- (2.985,78) node[above,font=\small] {D du sodium (589 nm)};
\draw[<-,thick] (3.885,60) -- (3.885,70) node[above,font=\small] {H$\alpha$ (656 nm)};
\draw[->,thick] (0,-8) -- (5.0,-8) node[right,font=\small] {$\lambda$ croissant};
\node[below,font=\small] at (0.4,-8) {UV};
\node[below,font=\small] at (4.5,-8) {IR};
\end{tikzpicture}
\end{center}
\legende{Reproduction schématique du spectre solaire : fond continu de l'UV à l'IR, traversé de raies noires de Fraunhofer (les principales raies visibles sont repérées).}
\end{popfigure}

\courssub{Origine du fond continu : la photosphère, un corps chaud à 5800 K}

Le Soleil n'a pas de surface solide : ce que nous appelons sa « surface » est la \cle{photosphère}, une couche de gaz dense d'environ $500\ \mathrm{km}$ d'épaisseur, portée à une température moyenne $T \approx 5800\ \mathrm{K}$. Un gaz dense et chaud se comporte, pour le rayonnement, comme un \emph{corps chaud} : les agitations thermiques incessantes de ses particules produisent des photons de \emph{toutes} les énergies, d'où un spectre continu.

\begin{propriete}[Fond continu et loi de Wien]
Le fond continu du spectre solaire suit approximativement le profil d'un corps chaud à la température $T$ de la photosphère. La longueur d'onde du maximum d'émission est donnée par la \cle{loi de Wien} :
\[ \lambda_{\max} \cdot T = \num{2,90e-3}\ \mathrm{m \cdot K} \]
où $T$ s'exprime en kelvins et $\lambda_{\max}$ en mètres. Réciproquement, la mesure de $\lambda_{\max}$ permet de déterminer la température de surface d'une étoile :
\[ T = \frac{\num{2,90e-3}}{\lambda_{\max}}. \]
\end{propriete}

\begin{exemplebox}[Température de surface du Soleil]
Le maximum d'émission du spectre solaire est mesuré à $\lambda_{\max} \approx 500\ \mathrm{nm}$ (lumière verte-bleue — c'est d'ailleurs à cette couleur que notre œil est le plus sensible, fruit de l'évolution !). Appliquons la loi de Wien :
\[ T = \frac{\num{2,90e-3}\ \mathrm{m \cdot K}}{500\times 10^{-9}\ \mathrm{m}} = \num{5,8e3}\ \mathrm{K}. \]
On retrouve bien $T \approx 5800\ \mathrm{K}$, soit environ $5500\ ^\circ\mathrm{C}$. À titre de comparaison, le filament d'une lampe à incandescence n'atteint que $2700\ \mathrm{K}$ : son maximum est dans l'infrarouge, c'est pourquoi elle chauffe plus qu'elle n'éclaire.
\end{exemplebox}

\courssub{Origine des raies noires : l'absorption sélective par la chromosphère}

Au-dessus de la photosphère s'étend la \cle{chromosphère}, une atmosphère solaire de gaz beaucoup plus \emph{dilué} et plus \emph{froid} (environ $4500\ \mathrm{K}$ dans ses couches basses), contenant des atomes d'hydrogène, d'hélium, de sodium, de calcium, de fer, etc.

Appliquons maintenant le modèle corpusculaire de la lumière, cœur de toute cette leçon. Les photons émis par la photosphère traversent la chromosphère. Un photon d'énergie $E = h\nu$ ne peut être absorbé par un atome que si son énergie correspond \emph{exactement} à l'écart entre deux niveaux d'énergie de cet atome :
\[ \Delta E = E_{\text{final}} - E_{\text{initial}} = h\nu = \frac{hc}{\lambda}. \]
\begin{itemize}
\item Les photons dont l'énergie ne correspond à \emph{aucune} transition des atomes de la chromosphère la traversent sans interaction : ils forment le fond continu.
\item Les photons dont l'énergie correspond exactement à une transition sont \emph{absorbés} : ils manquent à l'arrivée, d'où une \emph{raie noire} à la longueur d'onde $\lambda = hc/\Delta E$.
\end{itemize}
Les atomes excités par absorption se désexcitent ensuite en réémettant des photons de même énergie, mais \emph{dans toutes les directions} : la direction de l'observateur se trouve donc appauvrie en ces photons précis, et la raie reste sombre.

\begin{propriete}[Structure du spectre solaire]
Le spectre solaire est la superposition de deux contributions :
\[ \text{Spectre solaire} = \underbrace{\text{fond continu de corps chaud}}_{\text{photosphère, } T \approx 5800\ \mathrm{K}} \;+\; \underbrace{\text{raies d'absorption}}_{\text{chromosphère (H, He, Na, Ca, Fe\dots)}} \]
La position de $\lambda_{\max}$ renseigne sur la \emph{température} de la photosphère ; la position des raies noires renseigne sur la \emph{composition chimique} de la chromosphère.
\end{propriete}

\begin{popfigure}
\begin{center}
\begin{tikzpicture}[scale=1.0,>=Stealth]
% photosphère
\fill[poporange!80] (0,0) circle (1.5);
\fill[popgold!60] (0,0) circle (1.1);
\node[font=\small\bfseries] at (0,0.2) {Photosphère};
\node[font=\scriptsize] at (0,-0.35) {$T \approx 5800$ K};
\node[font=\scriptsize] at (0,-0.75) {émission continue};
% chromosphère
\draw[thick,popink,dashed] (0,0) circle (1.85);
\node[font=\scriptsize,popink] at (0,-2.15) {Chromosphère (gaz froid : H, Na, Ca, Fe\dots)};
% photons sortants
\draw[->,very thick,popgreen] (1.85,0.3) -- (3.4,0.3) node[midway,above,font=\scriptsize] {photons quelconques};
\draw[->,very thick,popgreen] (1.7,0.9) -- (3.4,1.3);
\draw[->,very thick,popgreen] (1.7,-0.9) -- (3.4,-1.3);
% photon absorbé
\draw[->,very thick,popink] (1.85,-0.3) -- (2.6,-0.3);
\draw[thick,popink] (2.6,-0.45) -- (2.75,-0.15);
\draw[thick,popink] (2.6,-0.15) -- (2.75,-0.45);
\node[font=\scriptsize,popink,anchor=west] at (2.85,-0.55) {photon absorbé : $h\nu = \Delta E$};
% spectre final
\begin{scope}[xshift=4.2cm,yshift=-0.9cm]
\fill[poppurple!50] (0,0) rectangle (0.5,1.8);
\fill[popblue!50] (0.5,0) rectangle (1.0,1.8);
\fill[popgreen!50] (1.0,0) rectangle (1.5,1.8);
\fill[popgold!60] (1.5,0) rectangle (2.0,1.8);
\fill[poporange!60] (2.0,0) rectangle (2.5,1.8);
\fill[popink!60] (2.5,0) rectangle (3.0,1.8);
\fill[black] (0.68,0) rectangle (0.72,1.8);
\fill[black] (1.85,0) rectangle (1.89,1.8);
\fill[black] (2.42,0) rectangle (2.46,1.8);
\node[font=\scriptsize,below] at (1.5,0) {spectre observé : fond continu + raies noires};
\end{scope}
\end{tikzpicture}
\end{center}
\legende{Coupe schématique du Soleil : la photosphère émet un rayonnement continu ; la chromosphère absorbe sélectivement les photons dont l'énergie $h\nu$ égale un écart $\Delta E$ entre niveaux de ses atomes ; le spectre observé présente des raies noires.}
\end{popfigure}

\courssub{Applications : composition chimique et température des étoiles}

\begin{methode}[Analyser le spectre d'une étoile]
\begin{enumerate}
\item \textbf{Repérer $\lambda_{\max}$}, la longueur d'onde du maximum du fond continu, puis calculer la température de surface par la loi de Wien : $T = \num{2,90e-3}/\lambda_{\max}$. Une étoile bleutée ($\lambda_{\max}$ petit) est plus chaude qu'une étoile rougeoyante.
\item \textbf{Relever les longueurs d'onde des raies noires} et les comparer aux spectres d'émission des éléments mesurés en laboratoire. Chaque coïncidence prouve la présence de l'élément correspondant dans l'atmosphère de l'étoile.
\item \textbf{Conclure} : température de surface + composition chimique, obtenues sans quitter la Terre !
\end{enumerate}
\end{methode}

\begin{exoresolu}[La raie D du sodium dans le spectre solaire]
\textbf{Énoncé.} On observe dans le spectre solaire une raie noire intense à $\lambda = 589\ \mathrm{nm}$ (raie D de Fraunhofer). En laboratoire, la vapeur de sodium émet une raie jaune à cette même longueur d'onde, correspondant à une transition d'énergie $\Delta E$.\\
1. Calculer l'énergie du photon absorbé, en joules puis en électronvolts.\\
2. Que peut-on conclure sur la composition du Soleil ?\\
On donne : $h = \num{6,63e-34}\ \mathrm{J \cdot s}$, $c = \num{3,00e8}\ \mathrm{m/s}$, $1\ \mathrm{eV} = \num{1,60e-19}\ \mathrm{J}$.
\tcblower
\textbf{Solution.}\\
1. D'après la relation de Planck-Einstein et $\nu = c/\lambda$ :
\[ E = h\nu = \frac{hc}{\lambda} = \frac{\num{6,63e-34}\times \num{3,00e8}}{589\times 10^{-9}} = \num{3,38e-19}\ \mathrm{J}. \]
Conversion en électronvolts :
\[ E = \frac{\num{3,38e-19}}{\num{1,60e-19}} \approx 2,11\ \mathrm{eV}. \]
Vérification dimensionnelle : $\dfrac{[h][c]}{[\lambda]} = \dfrac{\mathrm{J \cdot s}\times \mathrm{m/s}}{\mathrm{m}} = \mathrm{J}$. Correct.\\
2. La raie noire à $589\ \mathrm{nm}$ signifie que des photons de cette énergie exacte ont été absorbés par des atomes de la chromosphère. Or seul le sodium possède une transition à cette énergie. \textbf{Conclusion : il y a du sodium dans l'atmosphère du Soleil.} C'est ainsi qu'on a établi que le Soleil contient environ 74 \% d'hydrogène, 24 \% d'hélium et 2 \% d'éléments plus lourds (en masse).
\end{exoresolu}

\begin{attention}[Une raie sombre ne signifie pas « élément absent » !]
C'est le piège classique du Probatoire. Une raie \emph{noire} à la longueur d'onde d'un élément prouve au contraire sa \textbf{présence} dans l'atmosphère de l'étoile : ce sont ses atomes qui ont absorbé les photons. Rédige ainsi : « la raie d'absorption à $589\ \mathrm{nm}$ atteste la présence de sodium dans la chromosphère », et jamais « le Soleil ne contient pas de sodium ». Autre piège : ne pas confondre spectre d'\emph{émission} (raies brillantes sur fond noir, gaz chaud dilué) et spectre d'\emph{absorption} (raies noires sur fond continu, gaz froid devant une source chaude).
\end{attention}

\begin{savaistu}[L'analyse spectrale au service de la société]
La spectroscopie ne sert pas qu'aux astronomes ! Elle permet de contrôler la \emph{qualité des matériaux} (pureté des métaux dans l'industrie), de \emph{détecter les polluants} dans l'air des grandes villes comme Douala ou Yaoundé (dioxyde de soufre, oxydes d'azote, plomb), d'analyser la composition des eaux, et même d'identifier des substances en médecine légale. L'hélium doit d'ailleurs son nom au Soleil (\emph{helios} en grec) : il a été découvert en 1868 par sa raie jaune dans le spectre solaire, \emph{avant} d'être identifié sur Terre !
\end{savaistu}

\courssub{Bilan de la leçon}

\begin{aretenir}[L'essentiel de la leçon « Interaction lumière-matière : le photon »]
\begin{itemize}
\item La lumière est constituée de \cle{photons}, grains d'énergie $E = h\nu = \dfrac{hc}{\lambda}$, avec $h = \num{6,63e-34}\ \mathrm{J \cdot s}$ et $\lambda = c/\nu$.
\item Les niveaux d'énergie de la matière sont \cle{quantifiés} : un atome ne peut échanger que des énergies bien précises $\Delta E = h\nu$, d'où l'\emph{émission} (photon émis, l'atome descend d'un niveau), l'\emph{absorption} (photon absorbé, l'atome monte d'un niveau) et la \emph{diffusion} (photon renvoyé dans une autre direction).
\item Le \cle{profil spectral} d'un corps chaud est continu, avec un maximum à $\lambda_{\max} = \num{2,90e-3}/T$ (loi de Wien) ; celui d'un gaz dilué est un spectre de raies, carte d'identité de l'élément.
\item Le spectre solaire combine les deux : fond continu de la photosphère ($T \approx 5800\ \mathrm{K}$) et raies de Fraunhofer absorbées par la chromosphère, révélant sa composition chimique.
\end{itemize}
\end{aretenir}

\begin{popfigure}
\begin{center}
\begin{center}\tcbox[colback=poporange!60,colframe=poporange!60,coltext=white,arc=9pt,boxsep=4pt,left=6pt,right=6pt]{\bfseries\small \textbf{Le photon} $E = h\nu = hc/\lambda$}\end{center}
\vspace{-2pt}
\begin{tcbraster}[raster columns=2,raster equal height=rows,raster column skip=6pt,raster row skip=6pt,enhanced,blankest]
\begin{tcolorbox}[enhanced,colback=white,colframe=popgreen!40,boxrule=0.9pt,arc=3pt,title={Niveaux d'énergie quantifiés},fonttitle=\bfseries\scriptsize,coltitle=white,colbacktitle=popgreen!40,left=4pt,right=4pt,top=2pt,bottom=2pt,boxsep=1pt,toptitle=1.5pt,bottomtitle=1.5pt,halign=left]
\begin{itemize}[nosep,leftmargin=*,itemsep=1pt,font=\scriptsize]
\item Absorption : $\Delta E = h\nu$
\item Émission : $\Delta E = h\nu$
\end{itemize}
\end{tcolorbox}
\begin{tcolorbox}[enhanced,colback=white,colframe=popblue!40,boxrule=0.9pt,arc=3pt,title={Spectres},fonttitle=\bfseries\scriptsize,coltitle=white,colbacktitle=popblue!40,left=4pt,right=4pt,top=2pt,bottom=2pt,boxsep=1pt,toptitle=1.5pt,bottomtitle=1.5pt,halign=left]
\begin{itemize}[nosep,leftmargin=*,itemsep=1pt,font=\scriptsize]
\item Raies : carte d'identité de l'atome
\item Continu : corps chaud
\end{itemize}
\end{tcolorbox}
\begin{tcolorbox}[enhanced,colback=white,colframe=popgold!50,boxrule=0.9pt,arc=3pt,title={Corps chaud},fonttitle=\bfseries\scriptsize,coltitle=white,colbacktitle=popgold!50,left=4pt,right=4pt,top=2pt,bottom=2pt,boxsep=1pt,toptitle=1.5pt,bottomtitle=1.5pt,halign=left]
\begin{itemize}[nosep,leftmargin=*,itemsep=1pt,font=\scriptsize]
\item Wien : $\lambda_{\max} T = \num{2,9e-3}$
\end{itemize}
\end{tcolorbox}
\begin{tcolorbox}[enhanced,colback=white,colframe=poppink!40,boxrule=0.9pt,arc=3pt,title={Spectre solaire},fonttitle=\bfseries\scriptsize,coltitle=white,colbacktitle=poppink!40,left=4pt,right=4pt,top=2pt,bottom=2pt,boxsep=1pt,toptitle=1.5pt,bottomtitle=1.5pt,halign=left]
\begin{itemize}[nosep,leftmargin=*,itemsep=1pt,font=\scriptsize]
\item Fraunhofer : composition
\item $\lambda_{\max}$ : température
\end{itemize}
\end{tcolorbox}
\end{tcbraster}

\end{center}
\legende{Carte mentale de synthèse : le photon relie les niveaux d'énergie, les spectres et la température des corps chauds.}
\end{popfigure}

\begin{exercice}[Pour s'entraîner]
Une étoile présente un maximum d'émission à $\lambda_{\max} = 725\ \mathrm{nm}$ et des raies noires à $410\ \mathrm{nm}$, $434\ \mathrm{nm}$, $486\ \mathrm{nm}$ et $656\ \mathrm{nm}$ (raies de Balmer de l'hydrogène).\\
1. Déterminer sa température de surface. Cette étoile est-elle plus chaude ou plus froide que le Soleil ? Quelle est sa couleur apparente ?\\
2. Quel élément est présent dans son atmosphère ? Justifier.
\end{exercice}

\begin{corrige}[Exercice]
1. Loi de Wien : $T = \dfrac{\num{2,90e-3}}{725\times 10^{-9}} = \num{4,0e3}\ \mathrm{K}$. C'est plus froid que le Soleil ($5800\ \mathrm{K}$) ; le maximum étant dans le rouge proche, l'étoile paraît rouge-orangée (comme Bételgeuse).\\
2. Les quatre raies noires coïncident exactement avec les raies d'émission de l'hydrogène mesurées en laboratoire : des atomes d'hydrogène de l'atmosphère de l'étoile ont absorbé les photons d'énergie $hc/\lambda$ correspondant à leurs transitions. L'atmosphère de l'étoile contient donc de l'\textbf{hydrogène}.
\end{corrige}

\newpage

\coursec{Activité d'intégration}

\courssub{Objectifs de l'activité}
Cette activité d'intégration vise à mobiliser l'ensemble des \textbf{savoir-faire} de la leçon « Interaction lumière-matière : le photon » dans une situation authentique d'analyse spectrale. À l'issue de cette mission, vous devez être capable de :
\begin{itemize}
    \item Appliquer la \cle{loi de Wien} pour déterminer la température de surface d'un corps chaud (le Soleil) à partir de son profil spectral continu.
    \item Utiliser les relations \cle{$\Delta E = h\nu = \dfrac{hc}{\lambda}$} et \cle{$\lambda = \dfrac{c}{\nu}$} pour calculer l'énergie d'un photon (en joules et en électronvolts) et identifier la transition atomique correspondante sur un diagramme de niveaux d'énergie.
    \item Exploiter un diagramme de niveaux d'énergie de l'\cle{hydrogène} et du \cle{sodium} pour justifier l'origine de raies d'absorption précises dans le spectre solaire.
    \item Rédiger une explication scientifique rigoureuse, accessible à un élève de Seconde, articulant les concepts de \cle{photon}, d'\cle{absorption}, d'\cle{émission} et de \cle{niveaux d'énergie quantifiés} pour rendre compte de la structure du spectre solaire (fond continu + raies noires).
\end{itemize}

\courssub{Situation : Mission Astronomie — Portrait du Soleil depuis la classe}

Dans le cadre du club « Sciences et Espace » de votre établissement (Lycée de Nkolbisson, Yaoundé), vous participez à un atelier d'astronomie amateur encadré par votre professeur de Physique et un ingénieur de l'\emph{Institut de Recherche Géologique et Minière (IRGM)}. L'objectif : caractériser le Soleil sans quitter la salle de classe, en exploitant des données spectrales réelles fournies par le satellite \emph{SDO (Solar Dynamics Observatory)} de la NASA.

Votre binôme reçoit un \textbf{dossier d'analyse} contenant trois documents techniques :

\medskip
\noindent\textbf{Document 1 : Profil spectral continu du Soleil}
La courbe de densité spectrale de puissance $P(\lambda)$ (en $\text{W}\cdot\text{m}^{-3}$) en fonction de la longueur d'onde $\lambda$ présente un maximum net pour :
\[
\lambda_{\text{max}} = 502\,\text{nm} = 5{,}02 \times 10^{-7}\,\text{m}.
\]
Ce profil est caractéristique d'un \cle{corps noir} à température constante.

\medskip
\noindent\textbf{Document 2 : Extrait du spectre solaire (région visible)}
Un spectromètre à haute résolution révèle un fond continu sur lequel se superposent de nombreuses \cle{raies d'absorption} (raies noires). Trois raies intenses sont relevées avec précision :
\begin{center}
\begin{tabular}{|c|c|c|}\hline
\textbf{Raie} & \textbf{Longueur d'onde $\lambda$ (nm)} & \textbf{Couleur observée} \\ \hline
Raie $\alpha$ & $656{,}3$ & Rouge (H$\alpha$) \\
Raie $\beta$ & $589{,}0$ & Jaune-orangé (Doublet Na D) \\
Raie $\gamma$ & $486{,}1$ & Bleu-vert (H$\beta$) \\ \hline
\end{tabular}
\end{center}

\medskip
\noindent\textbf{Document 3 : Diagrammes simplifiés des niveaux d'énergie (en électronvolts, eV)}
Les valeurs sont données par rapport à l'état fondamental ($E=0$ eV pour l'ionisé, valeurs négatives pour les états liés).

\begin{popfigure}
\centering
\begin{tikzpicture}[scale=0.8, every node/.style={font=\small}]
    % Hydrogène
    \draw[thick] (0,0) -- (3,0) node[right] {$n=1 \quad E_1 = -13{,}60\,\text{eV}$};
    \draw[thick] (0,-1.51) -- (3,-1.51) node[right] {$n=2 \quad E_2 = -3{,}40\,\text{eV}$};
    \draw[thick] (0,-3.02) -- (3,-3.02) node[right] {$n=3 \quad E_3 = -1{,}51\,\text{eV}$};
    \draw[thick] (0,-4.25) -- (3,-4.25) node[right] {$n=4 \quad E_4 = -0{,}85\,\text{eV}$};
    \draw[thick] (0,-5.44) -- (3,-5.44) node[right] {$n=\infty \quad E_\infty = 0{,}00\,\text{eV}$};
    \node[align=center, popdark] at (1.5, 0.7) {\textbf{Hydrogène (H)}};
    % Flèches transitions absorption (vers le haut)
    \draw[popblue, thick, ->] (3.5, -1.51) -- (3.5, -3.02) node[midway, right] {$3 \leftarrow 2$};
    \draw[popblue, thick, ->] (3.5, -3.02) -- (3.5, -4.25) node[midway, right] {$4 \leftarrow 2$};
    
    % Sodium (décalé à droite)
    \begin{scope}[xshift=6cm]
    \draw[thick] (0,0) -- (3,0) node[right] {$3s \quad E = -5{,}14\,\text{eV}$};
    \draw[thick] (0,-2.10) -- (3,-2.10) node[right] {$3p \quad E = -3{,}04\,\text{eV}$};
    \draw[thick] (0,-3.75) -- (3,-3.75) node[right] {$4s \quad E = -1{,}95\,\text{eV}$};
    \draw[thick] (0,-5.14) -- (3,-5.14) node[right] {$E_\infty = 0{,}00\,\text{eV}$};
    \node[align=center, popdark] at (1.5, 0.7) {\textbf{Sodium (Na)}};
    \draw[poporange, thick, ->] (3.5, 0) -- (3.5, -2.10) node[midway, right] {$3p \leftarrow 3s$};
    \end{scope}
\end{tikzpicture}
\legende{Diagrammes des niveaux d'énergie de l'hydrogène (série de Balmer) et du sodium (transition de résonance). Les flèches indiquent les absorptions photoniques (énergie croissante).}
\end{popfigure}

\courssub{Consignes}
À partir des documents fournis et de vos connaissances du cours, réaliser les tâches suivantes en montrant toute la démarche (formule, substitution numérique, résultat encadré avec unité) :

\begin{enumerate}
    \item \textbf{Température de surface :} En utilisant la \cle{loi de Wien}, calculer la température de surface $T_{\odot}$ du Soleil (en kelvin). Préciser la valeur de la constante de Wien utilisée.
    \item \textbf{Analyse de la raie H$\alpha$ (656,3 nm) :}
    \begin{enumerate}
        \item Calculer l'énergie $E_{\text{ph}}$ du photon correspondant à cette raie, en \textbf{joules (J)} puis en \textbf{électronvolts (eV)}.
        \item Sur le diagramme de l'hydrogène, identifier la transition d'absorption responsable de cette raie (préciser les numéros quantiques principaux $n_i$ et $n_f$). Justifier par le calcul de $\Delta E$.
    \end{enumerate}
    \item \textbf{Analyse de la raie Na D (589,0 nm) :}
    \begin{enumerate}
        \item Calculer l'énergie du photon à 589,0 nm en eV.
        \item Montrer que cette énergie correspond à la transition fondamentale (de résonance) du sodium $3p \leftarrow 3s$ en utilisant les valeurs du diagramme.
    \end{enumerate}
    \item \textbf{Rapport de synthèse (10 lignes maximum) :} Rédiger un court texte adressé à un élève de Seconde pour expliquer pourquoi le spectre solaire présente \textbf{un fond continu} et \textbf{des raies noires}. Votre explication doit impérativement utiliser les mots : \cle{photon}, \cle{absorption}, \cle{émission}, \cle{niveaux d'énergie}.
\end{enumerate}

\courssub{Ressources à mobiliser}
\begin{itemize}
    \item \textbf{Loi de Wien :} $\lambda_{\text{max}} T = b$ avec $b = 2{,}898 \times 10^{-3}\,\text{m}\cdot\text{K}$ (constante de Wien).
    \item \textbf{Relations onde-particule :} $c = \lambda \nu = 3{,}00 \times 10^8\,\text{m}\cdot\text{s}^{-1}$ ; $E = h\nu = \dfrac{hc}{\lambda}$.
    \item \textbf{Constantes :} $h = 6{,}63 \times 10^{-34}\,\text{J}\cdot\text{s}$ ; $1\,\text{eV} = 1{,}60 \times 10^{-19}\,\text{J}$ ; $hc \approx 1240\,\text{eV}\cdot\text{nm}$ (raccourci numérique très utile).
    \item \textbf{Modèle corpusculaire :} La lumière échange son énergie par paquets (photons) d'énergie $E = h\nu$.
    \item \textbf{Quantification :} L'atome n'accepte d'échanger de l'énergie que par sauts entre niveaux discrets : $\Delta E = E_f - E_i = h\nu_{\text{photon}}$.
    \item \textbf{Origine du spectre :} Fond continu $\leftrightarrow$ émission thermique du corps noir (photosphère) ; Raies noires $\leftrightarrow$ absorption sélective par les atomes de l'atmosphère solaire (chromosphère).
\end{itemize}

\courssub{Démarche et corrigé commenté}
\begin{exoresolu}[Corrigé détaillé de la Mission Astronomie]
\textbf{Énoncé :} Voir les consignes ci-dessus.

\tcblower

\textbf{Solution détaillée}

\medskip
\noindent\textbf{1. Température de surface du Soleil par la loi de Wien}
\medskip
La loi de Wien relie la longueur d'onde d'émission maximale $\lambda_{\text{max}}$ d'un corps noir à sa température thermodynamique $T$ :
\[
\lambda_{\text{max}} T = b \quad \Rightarrow \quad T = \frac{b}{\lambda_{\text{max}}}
\]
Données : $b = 2{,}898 \times 10^{-3}\,\text{m}\cdot\text{K}$ ; $\lambda_{\text{max}} = 502\,\text{nm} = 502 \times 10^{-9}\,\text{m} = 5{,}02 \times 10^{-7}\,\text{m}$.
\[
T_{\odot} = \frac{2{,}898 \times 10^{-3}}{5{,}02 \times 10^{-7}} = \frac{2{,}898}{5{,}02} \times 10^{4} \approx 0{,}5773 \times 10^{4} \approx 5773\,\text{K}
\]
\fbox{$T_{\odot} \approx 5\,770\,\text{K}$} (arrondi à 3 chiffres significatifs cohérents avec $\lambda_{\text{max}}$).
\emph{Analyse dimensionnelle :} $[b] = \text{m}\cdot\text{K}$, $[\lambda] = \text{m}$ $\Rightarrow$ $[T] = \text{K}$. Correct.

\medskip
\noindent\textbf{2. Analyse de la raie H$\alpha$ (656,3 nm)}
\medskip
\paragraph{a) Énergie du photon}
On utilise la relation $E = \dfrac{hc}{\lambda}$.
\emph{En électronvolts (eV) :} On utilise le raccourci $hc \approx 1240\,\text{eV}\cdot\text{nm}$.
\[
E_{\text{eV}} = \frac{1240\,\text{eV}\cdot\text{nm}}{656{,}3\,\text{nm}} \approx 1{,}889\,\text{eV}
\]
\emph{En joules (J) :} Conversion via $1\,\text{eV} = 1{,}60 \times 10^{-19}\,\text{J}$.
\[
E_{\text{J}} = 1{,}889 \times 1{,}60 \times 10^{-19} \approx 3{,}02 \times 10^{-19}\,\text{J}
\]
(Vérification directe : $E = \frac{6{,}63\times10^{-34} \times 3{,}00\times10^8}{656{,}3\times10^{-9}} = 3{,}03\times10^{-19}\,\text{J}$, cohérent).
\fbox{$E_{\text{ph}} \approx 1{,}89\,\text{eV} \quad (\text{soit } 3{,}02 \times 10^{-19}\,\text{J})$}

\paragraph{b) Identification de la transition dans l'hydrogène}
Le diagramme donne les niveaux : $E_1=-13{,}60$, $E_2=-3{,}40$, $E_3=-1{,}51$, $E_4=-0{,}85$ eV.
Une raie d'absorption correspond à un saut vers un niveau d'énergie \emph{supérieur} ($n_i < n_f$). L'énergie absorbée est $\Delta E = E_f - E_i > 0$.
On teste les transitions de la série de Balmer (arrivée sur $n=2$, domaine visible) :
\begin{align*}
\text{Transition } 3 \leftarrow 2 \quad (n_i=2 \to n_f=3) &: \Delta E = E_3 - E_2 = (-1{,}51) - (-3{,}40) = 1{,}89\,\text{eV}. \\
\text{Transition } 4 \leftarrow 2 \quad (n_i=2 \to n_f=4) &: \Delta E = E_4 - E_2 = (-0{,}85) - (-3{,}40) = 2{,}55\,\text{eV}.
\end{align*}
Seule la transition $3 \leftarrow 2$ donne $\Delta E = 1{,}89\,\text{eV}$, identique à l'énergie du photon calculée.
\fbox{La raie H$\alpha$ (656,3 nm) correspond à l'absorption $n_i=2 \to n_f=3$ (transition $3 \leftarrow 2$) de l'hydrogène.}

\medskip
\noindent\textbf{3. Analyse de la raie Na D (589,0 nm)}
\medskip
\paragraph{a) Énergie du photon}
\[
E_{\text{eV}} = \frac{1240}{589{,}0} \approx 2{,}105\,\text{eV}
\]
\fbox{$E_{\text{ph}} \approx 2{,}11\,\text{eV}$}

\paragraph{b) Correspondance avec la transition fondamentale du sodium}
Le diagramme du sodium donne : $E(3s) = -5{,}14\,\text{eV}$ (état fondamental) ; $E(3p) = -3{,}04\,\text{eV}$ (premier état excité).
La transition de résonance (fondamentale) est $3p \leftarrow 3s$ ($n_i=3s \to n_f=3p$).
\[
\Delta E_{\text{Na}} = E(3p) - E(3s) = (-3{,}04) - (-5{,}14) = 2{,}10\,\text{eV}.
\]
On trouve $\Delta E_{\text{Na}} = 2{,}10\,\text{eV}$ pour une énergie photonique de $2{,}11\,\text{eV}$. L'accord est excellent (différence $< 0{,}5\,\%$ due aux arrondis des niveaux fournis).
\fbox{La raie à 589,0 nm correspond bien à la transition de résonance $3s \to 3p$ du sodium (raies D).}

\medskip
\noindent\textbf{4. Rapport de synthèse (exemple de rédaction attendue)}
\medskip
\begin{quote}
\textit{« Le Soleil émet un \textbf{fond continu} car sa surface (photosphère), à environ 5800 K, se comporte comme un \textbf{corps noir} : les atomes y sont si denses qu'ils émettent des \textbf{photons} de toutes les énergies par \textbf{émission} thermique. En traversant l'atmosphère solaire (chromosphère), plus froide et peu dense, cette lumière rencontre des atomes d'hydrogène et de sodium. Ces atomes ne peuvent absorber que des \textbf{photons} dont l'énergie correspond \textbf{exactement} à l'écart entre deux de leurs \textbf{niveaux d'énergie quantifiés} ($\Delta E = h\nu$). C'est cette \textbf{absorption} sélective qui crée les \textbf{raies noires} : les photons de ces énergies précises manquent dans le spectre reçu. »} 
\end{quote}
\emph{(9 lignes, vocabulaire imposé respecté, causalité physique claire).}
\end{exoresolu}

\courssub{Bilan de l'activité}
Cette mission « Portrait du Soleil » a permis de valider les \textbf{compétences clés} du module Optique :
\begin{itemize}
    \item \textbf{Modélisation corps noir / Loi de Wien :} Passer d'une observation spectrale ($\lambda_{\text{max}}$) à une grandeur physique macroscopique ($T \approx 5770\,\text{K}$), confirmant que le Soleil rayonne comme un corps noir.
    \item \textbf{Quantification de l'énergie / Modèle corpusculaire :} Utiliser $E = hc/\lambda$ pour passer de l'onde (longueur d'onde mesurée) au corpuscule (énergie du photon en J et eV), pierre angulaire de la physique quantique.
    \item \textbf{Exploitation de diagrammes de niveaux :} Lire un diagramme d'énergie, calculer des écarts $\Delta E$, et faire le lien univoque $\Delta E = h\nu$ pour identifier l'espèce chimique (H, Na) et la transition quantique ($n=2\to3$, $3s\to3p$). C'est la base de la \cle{spectroscopie d'absorption}.
    \item \textbf{Communication scientifique :} Structurer un discours explicatif rigoureux (fond continu = émission thermique corps noir ; raies = absorption sélective par niveaux quantifiés) en utilisant le vocabulaire précis du programme (\cle{photon}, \cle{absorption}, \cle{émission}, \cle{niveaux d'énergie}).
\end{itemize}
\medskip
\begin{attention}[Pièges évités grâce à cette activité]
\begin{itemize}
    \item \textbf{Unités :} Ne pas confondre nm et m dans la loi de Wien ($b$ en m.K). Toujours convertir $\lambda$ en mètres.
    \item \textbf{Sens de la transition :} Absorption $\Rightarrow$ atome monte en énergie ($E_f > E_i$, $\Delta E > 0$). Émission $\Rightarrow$ atome descend ($E_f < E_i$). Sur le diagramme, les flèches d'absorption pointent vers le HAUT.
    \item \textbf{Origine du fond continu :} Ce n'est pas l'addition de raies d'émission, mais bien le rayonnement thermique d'un milieu dense (photosphère) approximé par un corps noir.
    \item \textbf{Précision du vocabulaire :} « Niveaux d'énergie » (quantifiés, discrets) et non « couches » ou « orbites » (termes du modèle de Bohr dépassé, bien que les numéros quantiques $n$ persistent).
\end{itemize}
\end{attention}

\begin{savaistu}[Le Soleil, laboratoire de physique nucléaire]
Les raies noires du spectre solaire (raies de \textbf{Fraunhofer}, 1814) ont permis de découvrir l'\textbf{hélium} (He) en 1868 : une raie jaune à 587,6 nm, proche du doublet Na D, ne correspondait à aucun élément terrestre connu. On l'a attribuée à un nouvel élément : \emph{hélios} (Soleil en grec). Aujourd'hui, la spectroscopie d'absorption (et d'émission) est l'outil principal pour connaître la composition chimique, la température, la pression, le champ magnétique (effet Zeeman) et la vitesse radiale (effet Doppler-Fizeau) des étoiles. Au Cameroun, l'IRGM et l'Université de Yaoundé I utilisent ces techniques pour l'analyse de roches (LIBS - Laser Induced Breakdown Spectroscopy) et la surveillance de la qualité de l'air.
\end{savaistu}

\newpage

\coursec{Méthodes et exercices résolus}

\courssub{Interaction lumière-matière : émission, absorption, diffusion}

\begin{definition}[Interaction lumière-matière]
Lorsqu'un rayonnement lumineux rencontre la matière, trois processus fondamentaux peuvent se produire :
\begin{itemize}
\item \textbf{L'émission} : un atome, une molécule ou un corps chaud cède de l'énergie sous forme de photons (lumière). C'est le cas d'une lampe, du Soleil, d'un filament.
\item \textbf{L'absorption} : un atome ou une molécule capte l'énergie d'un photon incident et passe à un état d'énergie supérieure. Le photon disparaît.
\item \textbf{La diffusion} : le photon change de direction (et parfois d'énergie) sans être détruit ni créé par une transition quantique interne. On distingue la diffusion élastique (Rayleigh, Mie : ciel bleu, nuages blancs) et la diffusion inélastique (Raman, Compton -- cette dernière en Terminale).
\end{itemize}
\end{definition}

\begin{experience}[Observation de spectres : signature de la matière]
\textbf{Matériel :} Spectroscope (ou réseau de diffraction), lampes à décharge (hydrogène, hélium, mercure, sodium), lampe à incandescence, lumière du Soleil (précautions !).
\textbf{Protocole :} Observer à travers le spectroscope la lumière émise par chaque source.
\textbf{Observations :}
\begin{itemize}
\item Lampe à incandescence / Soleil (noyau) : \textbf{spectre continu} (arc-en-ciel complet), signature d'un corps chaud dense.
\item Lampes à décharge (gaz raréfiés) : \textbf{spectres de raies} (raies colorées sur fond noir), chaque gaz a son motif unique.
\item Lumière du Soleil (bord du disque) : spectre continu traversé de \textbf{raies sombres} (raies de Fraunhofer), signature d'absorption par l'atmosphère solaire.
\end{itemize}
\textbf{Interprétation :} La matière n'échange l'énergie qu'avec des photons d'énergies bien précises, caractéristiques de sa structure quantique. C'est la base de la spectroscopie (identification chimique, température des astres, vélocimétrie Doppler).
\end{experience}

\begin{propriete}[Quantification de l'énergie de la matière]
L'énergie interne d'un système quantique (atome, molécule, noyau) ne peut prendre que des valeurs discrètes, appelées \textbf{niveaux d'énergie} $E_1, E_2, \dots, E_n$ (souvent négatifs par rapport à l'état libre $E_\infty = 0$). L'état de plus basse énergie $E_1$ est l'\textbf{état fondamental} ; les autres sont des \textbf{états excités}.
\end{propriete}

\begin{aretenir}[Modèle corpusculaire de la lumière : le photon]
La lumière se comporte comme un flux de particules d'énergie (quanta) appelées \textbf{photons}.
\begin{itemize}
\item Énergie d'un photon : $E = h\nu$ avec $h = \num{6,63e-34}~\text{J.s}$ (constante de Planck).
\item Relation onde-corpuscule : $\lambda = \dfrac{c}{\nu}$ avec $c = \num{3,00e8}~\text{m.s}^{-1}$ (vitesse de la lumière dans le vide).
\item Énergie en fonction de la longueur d'onde : $E = \dfrac{hc}{\lambda}$.
\item Conversion pratique : $hc \approx \num{1240}~\text{eV.nm}$ (utile pour calculs rapides en eV et nm).
\end{itemize}
\end{aretenir}

\courssub{Méthode 1 : Calculer l'énergie d'un photon à partir de sa fréquence ou de sa longueur d'onde}

\begin{methode}[Exploiter les relations $E = h\nu$ et $\lambda = \dfrac{c}{\nu}$]
\begin{enumerate}
\item \textbf{Identifier la donnée} : fréquence $\nu$ (en hertz) ou longueur d'onde $\lambda$ (en mètres). Convertir systématiquement en unités SI : $1~\text{nm} = 10^{-9}~\text{m}$, $1~\text{THz} = 10^{12}~\text{Hz}$.
\item \textbf{Si la longueur d'onde est donnée}, calculer d'abord la fréquence : $\nu = \dfrac{c}{\lambda}$ avec $c = \num{3,00e8}~\text{m.s}^{-1}$.
\item \textbf{Calculer l'énergie} : $E = h\nu$ avec $h = \num{6,63e-34}~\text{J.s}$. On peut combiner : $E = \dfrac{hc}{\lambda}$.
\item \textbf{Convertir en électronvolts si demandé} : $1~\text{eV} = \num{1,60e-19}~\text{J}$, donc $E(\text{eV}) = \dfrac{E(\text{J})}{\num{1,60e-19}}$.
\item \textbf{Vérifier l'ordre de grandeur} : un photon du domaine visible a une énergie de l'ordre de $10^{-19}$~J, soit quelques eV.
\end{enumerate}
\end{methode}

\begin{exoresolu}[Énergie d'un photon rouge et d'un photon ultraviolet]
\textbf{Énoncé.} On donne $h = \num{6,63e-34}~\text{J.s}$, $c = \num{3,00e8}~\text{m.s}^{-1}$, $1~\text{eV} = \num{1,60e-19}~\text{J}$.
\begin{enumerate}
\item Calculer l'énergie d'un photon de lumière rouge de longueur d'onde $\lambda_1 = 650~\text{nm}$, en joules puis en électronvolts.
\item Un photon ultraviolet a une fréquence $\nu_2 = \num{1,20e15}~\text{Hz}$. Calculer son énergie en eV et sa longueur d'onde. Comparer avec le photon rouge.
\end{enumerate}
\tcblower
\textbf{Solution.}
\begin{enumerate}
\item \textbf{Photon rouge.} Conversion : $\lambda_1 = 650~\text{nm} = \num{6,50e-7}~\text{m}$.\\
Fréquence : $\nu_1 = \dfrac{c}{\lambda_1} = \dfrac{\num{3,00e8}}{\num{6,50e-7}} = \num{4,62e14}~\text{Hz}$.\\
Énergie : $E_1 = h\nu_1 = \num{6,63e-34} \times \num{4,62e14} = \num{3,06e-19}~\text{J}$.\\
En électronvolts : $E_1 = \dfrac{\num{3,06e-19}}{\num{1,60e-19}} = 1,91~\text{eV}$.
\item \textbf{Photon UV.} $E_2 = h\nu_2 = \num{6,63e-34} \times \num{1,20e15} = \num{7,96e-19}~\text{J} = \dfrac{\num{7,96e-19}}{\num{1,60e-19}} = 4,97~\text{eV}$.\\
Longueur d'onde : $\lambda_2 = \dfrac{c}{\nu_2} = \dfrac{\num{3,00e8}}{\num{1,20e15}} = \num{2,50e-7}~\text{m} = 250~\text{nm}$.
\end{enumerate}
\textbf{Conclusion :} le photon ultraviolet ($4,97~\text{eV}$) est environ $2,6$ fois plus énergétique que le photon rouge ($1,91~\text{eV}$) : plus la longueur d'onde est petite, plus le photon est énergétique. C'est pourquoi les UV sont dangereux pour la peau (risque de mutation ADN), contrairement à la lumière rouge.
\end{exoresolu}

\courssub{Méthode 2 : Exploiter un diagramme de niveaux d'énergie}

\begin{methode}[Utiliser la relation $\Delta E = h\nu$ sur un diagramme d'énergie]
\begin{enumerate}
\item \textbf{Repérer les deux niveaux} concernés par la transition : niveau initial $E_i$ et niveau final $E_f$ (souvent négatifs, en eV, avec $E_\infty = 0$ pour l'atome ionisé).
\item \textbf{Calculer la variation d'énergie} : $\Delta E = |E_f - E_i|$, toujours positive.
\item \textbf{Identifier le phénomène} :
\begin{itemize}
\item $E_f < E_i$ : l'atome \textbf{émet} un photon d'énergie $h\nu = E_i - E_f$ (déssexcitation) ;
\item $E_f > E_i$ : l'atome \textbf{absorbe} un photon d'énergie exactement égale à $h\nu = E_f - E_i$ (excitation).
\end{itemize}
\item \textbf{En déduire la fréquence puis la longueur d'onde} : $\nu = \dfrac{\Delta E}{h}$ puis $\lambda = \dfrac{c}{\nu} = \dfrac{hc}{\Delta E}$.
\item \textbf{Convertir les eV en joules avant tout calcul} avec $h$ en J.s : $\Delta E(\text{J}) = \Delta E(\text{eV}) \times \num{1,60e-19}$.
\item \textbf{Conclure sur le domaine spectral} (UV, visible, IR) en comparant $\lambda$ aux bornes du visible ($400~\text{nm}$ à $800~\text{nm}$).
\end{enumerate}
\end{methode}

\begin{exoresolu}[Raie rouge de l'hydrogène (série de Balmer)]
\textbf{Énoncé.} Les niveaux d'énergie de l'atome d'hydrogène sont donnés par $E_n = -\dfrac{13,6}{n^2}$ (en eV). On donne $h = \num{6,63e-34}~\text{J.s}$, $c = \num{3,00e8}~\text{m.s}^{-1}$, $1~\text{eV} = \num{1,60e-19}~\text{J}$.
\begin{enumerate}
\item Calculer les énergies des niveaux $n = 2$ et $n = 3$.
\item L'atome passe du niveau $n = 3$ au niveau $n = 2$. S'agit-il d'une émission ou d'une absorption ? Calculer l'énergie du photon échangé.
\item Déterminer la longueur d'onde de la radiation correspondante. À quel domaine appartient-elle ?
\end{enumerate}
\tcblower
\textbf{Solution.}
\begin{enumerate}
\item $E_2 = -\dfrac{13,6}{2^2} = -3,40~\text{eV}$ ; \quad $E_3 = -\dfrac{13,6}{3^2} = -\dfrac{13,6}{9} = -1,51~\text{eV}$.
\item L'atome descend de $E_3 = -1,51~\text{eV}$ vers $E_2 = -3,40~\text{eV}$ : son énergie diminue, il \textbf{émet} un photon.\\
$\Delta E = E_3 - E_2 = -1,51 - (-3,40) = 1,89~\text{eV}$.\\
Conversion : $\Delta E = 1,89 \times \num{1,60e-19} = \num{3,02e-19}~\text{J}$.
\item Longueur d'onde : 
\begin{align*}
\lambda &= \dfrac{hc}{\Delta E} = \dfrac{\num{6,63e-34} \times \num{3,00e8}}{\num{3,02e-19}} = \num{6,59e-7}~\text{m} = 659~\text{nm}.
\end{align*}
\end{enumerate}
\textbf{Conclusion :} la radiation émise a pour longueur d'onde $\lambda \approx 659~\text{nm}$ : elle appartient au domaine visible (rouge). C'est la célèbre raie rouge $H_\alpha$ de l'hydrogène, observable dans les tubes spectraux en TP et dans les spectres stellaires. (La valeur de référence est $656,3~\text{nm}$ ; l'écart vient des arrondis sur $h$, $c$ et $E_n$).
\end{exoresolu}

\courssub{Méthode 3 : Déterminer si un photon peut être absorbé (quantification)}

\begin{methode}[Condition d'absorption d'un photon par un atome]
\begin{enumerate}
\item \textbf{Calculer l'énergie du photon incident} : $E = h\nu = \dfrac{hc}{\lambda}$ (convertir en eV si le diagramme est en eV).
\item \textbf{Comparer aux écarts du diagramme} : le photon n'est absorbé que si son énergie correspond \emph{exactement} à un écart entre deux niveaux : $E = E_f - E_i$.
\item \textbf{Cas particulier de l'ionisation} : si $E \geq |E_1|$ (énergie du niveau fondamental en valeur absolue), l'électron est arraché ; l'excès d'énergie devient énergie cinétique de l'électron : $E_c = E - |E_1|$.
\item \textbf{Conclure clairement} : photon absorbé (avec la transition correspondante) ou non absorbé (photon transmis ou diffusé élastiquement).
\end{enumerate}
\end{methode}

\begin{attention}[Réflexe de rédaction : quantification stricte]
Écrire explicitement la condition de quantification : « l'énergie du photon doit être \textbf{exactement égale} à l'écart entre deux niveaux ». Un photon d'énergie « presque égale » n'est \textbf{pas} absorbé (sauf élargissement naturel/pression/Doppler, négligeable au Probatoire). C'est ce qui distingue le modèle quantique du modèle classique.
\end{attention}

\begin{exoresolu}[Absorption sélective par la vapeur de mercure]
\textbf{Énoncé.} Un atome de mercure possède les niveaux d'énergie : $E_1 = -10,4~\text{eV}$ (fondamental), $E_2 = -5,5~\text{eV}$, $E_3 = -3,7~\text{eV}$. On l'éclaire avec des photons de longueur d'onde $\lambda = 254~\text{nm}$. On donne $h = \num{6,63e-34}~\text{J.s}$, $c = \num{3,00e8}~\text{m.s}^{-1}$, $1~\text{eV} = \num{1,60e-19}~\text{J}$.
\begin{enumerate}
\item Calculer l'énergie des photons incidents en eV.
\item Ces photons peuvent-ils être absorbés par un atome dans son état fondamental ? Si oui, préciser la transition.
\item Un photon d'énergie $6,0~\text{eV}$ peut-il être absorbé ? Justifier.
\end{enumerate}
\tcblower
\textbf{Solution.}
\begin{enumerate}
\item $E = \dfrac{hc}{\lambda} = \dfrac{\num{6,63e-34} \times \num{3,00e8}}{\num{2,54e-7}} = \num{7,83e-19}~\text{J}$, soit $E = \dfrac{\num{7,83e-19}}{\num{1,60e-19}} = 4,89~\text{eV}$.
\item Écarts possibles depuis le fondamental $E_1$ :
\begin{itemize}
\item $E_2 - E_1 = -5,5 - (-10,4) = 4,9~\text{eV}$ ;
\item $E_3 - E_1 = -3,7 - (-10,4) = 6,7~\text{eV}$.
\end{itemize}
L'énergie du photon ($4,89~\text{eV} \approx 4,9~\text{eV}$) correspond exactement à l'écart $E_2 - E_1$ : le photon est \textbf{absorbé} et l'atome passe du niveau $E_1$ au niveau $E_2$.
\item Pour un photon de $6,0~\text{eV}$ : aucun écart depuis $E_1$ ne vaut $6,0~\text{eV}$ ($4,9$ et $6,7~\text{eV}$), et $6,0~\text{eV} < |E_1| = 10,4~\text{eV}$ donc pas d'ionisation. Le photon n'est \textbf{pas absorbé} (il est transmis ou diffusé).
\end{enumerate}
\textbf{Conclusion :} seuls les photons dont l'énergie coïncide exactement avec un écart de niveaux sont absorbés : c'est la \cle{quantification} des échanges d'énergie. C'est ce principe qui fait fonctionner les lampes fluorescentes et les lampadaires à vapeur de mercure (lumière blanche bleutée) utilisés pour l'éclairage public à Yaoundé et Douala.
\end{exoresolu}

\courssub{Méthode 4 : Interpréter un spectre de raies (émission ou absorption)}

\begin{methode}[Analyser un spectre pour identifier un élément ou un phénomène]
\begin{enumerate}
\item \textbf{Identifier le type de spectre} : raies brillantes sur fond noir $\rightarrow$ spectre d'\textbf{émission} (gaz chaud sous faible pression) ; raies noires sur fond coloré continu $\rightarrow$ spectre d'\textbf{absorption} (gaz froid devant une source de lumière blanche).
\item \textbf{Relever les longueurs d'onde des raies} et les convertir en mètres.
\item \textbf{Calculer les énergies correspondantes} $E = \dfrac{hc}{\lambda}$ et les comparer aux écarts du diagramme de niveaux de l'élément suspecté.
\item \textbf{Exploiter l'unicité} : chaque élément chimique possède un spectre de raies qui lui est propre, véritable « carte d'identité ». La coïncidence des raies permet l'identification formelle.
\item \textbf{Pour un spectre d'absorption} : les raies noires se trouvent aux \emph{mêmes} longueurs d'onde que les raies d'émission du même gaz (les photons absorbés sont ceux que le gaz sait émettre).
\end{enumerate}
\end{methode}

\begin{exoresolu}[Identification d'un gaz par son spectre]
\textbf{Énoncé.} Le spectre d'émission d'une lampe inconnue présente deux raies intenses : $\lambda_a = 589~\text{nm}$ et $\lambda_b = 589,6~\text{nm}$ (doublet jaune). On donne les écarts d'énergie caractéristiques de quelques transitions : sodium : $2,11~\text{eV}$ ; mercure : $2,27~\text{eV}$ ; hydrogène : $1,89~\text{eV}$. On donne $h = \num{6,63e-34}~\text{J.s}$, $c = \num{3,00e8}~\text{m.s}^{-1}$, $1~\text{eV} = \num{1,60e-19}~\text{J}$.
\begin{enumerate}
\item Calculer l'énergie des photons associés à la raie $\lambda_a$, en eV.
\item Identifier le gaz contenu dans la lampe. Justifier.
\item Citer une application courante de ce type de lampe au Cameroun.
\end{enumerate}
\tcblower
\textbf{Solution.}
\begin{enumerate}
\item $E_a = \dfrac{hc}{\lambda_a} = \dfrac{\num{6,63e-34} \times \num{3,00e8}}{\num{5,89e-7}} = \num{3,38e-19}~\text{J}$, soit $E_a = \dfrac{\num{3,38e-19}}{\num{1,60e-19}} = 2,11~\text{eV}$.
\item Comparaison : $E_a = 2,11~\text{eV}$ correspond exactement à l'écart caractéristique du \textbf{sodium} (doublet jaune à $589,0$ et $589,6~\text{nm}$). Ni le mercure ($2,27~\text{eV}$) ni l'hydrogène ($1,89~\text{eV}$) ne conviennent. La lampe contient donc de la vapeur de sodium.
\item Les lampes à vapeur de sodium (basse ou haute pression), à la lumière jaune-orangée quasi monochromatique, équipent de nombreux lampadaires et l'éclairage des axes routiers (par exemple l'autoroute Yaoundé--Nsimalen) pour leur excellent rendement lumineux et leur faible coût énergétique.
\end{enumerate}
\textbf{Conclusion :} l'énergie des photons émis ($2,11~\text{eV}$) coïncide avec le doublet du sodium : le spectre de raies permet d'identifier sans ambiguïté l'élément chimique présent dans la source.
\end{exoresolu}

\courssub{Méthode 5 : Expliquer les caractéristiques du spectre solaire}

\begin{methode}[Interpréter la forme et les raies du spectre du Soleil]
\begin{enumerate}
\item \textbf{Décrire le profil spectral} : le spectre solaire est un spectre \textbf{continu} dont l'intensité passe par un maximum ; il ressemble au profil du \textbf{corps noir} (loi de Planck).
\item \textbf{Expliquer l'origine du fond continu} : la photosphère (surface visible du Soleil, environ $\num{5800}~\text{K}$) est un corps chaud dense et opaque qui émet un rayonnement continu dont le profil spectral ne dépend que de la température (lois de Wien et Stefan-Boltzmann).
\item \textbf{Expliquer les raies noires (raies de Fraunhofer)} : la chromosphère, couche de gaz plus froid ($\approx \num{4000}~\text{K}$) et peu dense entourant la photosphère, \textbf{absorbe} certains photons dont l'énergie correspond exactement aux écarts entre niveaux de ses atomes (H, He, Fe, Ca, Na...). Il en résulte des raies sombres superposées au fond continu.
\item \textbf{Exploiter les raies} : leur position permet d'identifier les éléments chimiques présents dans l'atmosphère solaire ; leur étude a d'ailleurs permis de découvrir l'hélium (du grec \emph{hélios}, Soleil) en 1868, bien avant sa découverte sur Terre (1895).
\item \textbf{Déterminer la température de surface} : repérer la longueur d'onde $\lambda_{\max}$ du maximum d'émission et appliquer la loi de Wien : $\lambda_{\max} \cdot T = \num{2,90e-3}~\text{m.K}$, avec $\lambda_{\max}$ en mètres et $T$ en kelvins.
\end{enumerate}
\end{methode}

\begin{exoresolu}[Température de surface du Soleil et raies de Fraunhofer]
\textbf{Énoncé.} Le profil spectral du Soleil présente un maximum d'intensité pour $\lambda_{\max} = 500~\text{nm}$. Son spectre montre aussi une raie sombre intense à $\lambda = 656~\text{nm}$, attribuée à l'hydrogène (transition entre les niveaux $n = 2$ et $n = 3$). On donne la constante de Wien $b = \num{2,90e-3}~\text{m.K}$, $h = \num{6,63e-34}~\text{J.s}$, $c = \num{3,00e8}~\text{m.s}^{-1}$, $1~\text{eV} = \num{1,60e-19}~\text{J}$.
\begin{enumerate}
\item Expliquer l'origine du fond continu du spectre solaire et celle de la raie sombre à $656~\text{nm}$.
\item Calculer la température de surface du Soleil.
\item Vérifier que la raie à $656~\text{nm}$ correspond bien à la transition $n = 3 \rightarrow n = 2$ de l'hydrogène ($E_2 = -3,40~\text{eV}$, $E_3 = -1,51~\text{eV}$).
\end{enumerate}
\tcblower
\textbf{Solution.}
\begin{enumerate}
\item Le \textbf{fond continu} provient de la photosphère, surface dense et chaude du Soleil, qui émet un rayonnement de corps chaud (corps noir) dont le profil ne dépend que de sa température. La \textbf{raie sombre} à $656~\text{nm}$ est une raie d'absorption : les atomes d'hydrogène de la chromosphère (couche externe plus froide) absorbent les photons d'énergie $E_3 - E_2$, ce qui creuse une raie noire dans le spectre continu.
\item Loi de Wien : 
\begin{align*}
T &= \dfrac{b}{\lambda_{\max}} = \dfrac{\num{2,90e-3}}{\num{5,00e-7}} = \num{5,80e3}~\text{K}.
\end{align*}
La température de surface du Soleil vaut environ $\num{5800}~\text{K}$.
\item Écart d'énergie : $\Delta E = E_3 - E_2 = -1,51 - (-3,40) = 1,89~\text{eV} = 1,89 \times \num{1,60e-19} = \num{3,02e-19}~\text{J}$.\\
Longueur d'onde associée : $\lambda = \dfrac{hc}{\Delta E} = \dfrac{\num{6,63e-34} \times \num{3,00e8}}{\num{3,02e-19}} = \num{6,59e-7}~\text{m} = 659~\text{nm} \approx 656~\text{nm}$ (écart dû aux arrondis des données $h, c, E_n$).
\end{enumerate}
\textbf{Conclusion :} le maximum d'émission à $500~\text{nm}$ donne une température de surface d'environ $\num{5800}~\text{K}$, et la raie sombre à $656~\text{nm}$ s'interprète parfaitement comme l'absorption par l'hydrogène de la chromosphère : l'analyse spectrale permet à la fois de « peser » la température des astres et d'en dresser la carte d'identité chimique.
\end{exoresolu}

\begin{savaistu}[L'hélium, découvert dans le Soleil avant la Terre]
En 1868, lors d'une éclipse solaire observée en Inde, l'astronome français Jules Janssen et l'anglais Norman Lockyer repèrent une raie jaune brillante à $587,6~\text{nm}$ dans le spectre de la chromosphère, qui ne correspond à aucun élément connu sur Terre. Lockyer la baptise \emph{hélium} (du grec \emph{hélios}, Soleil). Il faudra attendre 1895 pour que William Ramsay l'isole sur Terre dans un minerai d'uranium (cléveite). Aujourd'hui, l'hélium est crucial pour les IRM des hôpitaux (supraconductivité), le soudage TIG, et les ballons scientifiques (CNES, météo).
\end{savaistu}

\courssub{Méthode 6 : Déterminer la température d'un corps chaud par analyse spectrale}

\begin{methode}[Appliquer la loi de Wien]
\begin{enumerate}
\item \textbf{Tracer ou exploiter le profil spectral} $I = f(\lambda)$ de la source : repérer l'abscisse $\lambda_{\max}$ du maximum d'intensité.
\item \textbf{Convertir $\lambda_{\max}$ en mètres} (nm $\rightarrow$ m : multiplier par $10^{-9}$ ; $\mu$m $\rightarrow$ m : multiplier par $10^{-6}$).
\item \textbf{Appliquer la loi de Wien} : $\lambda_{\max} \times T = \num{2,90e-3}~\text{m.K}$, d'où $T = \dfrac{\num{2,90e-3}}{\lambda_{\max}}$.
\item \textbf{Exprimer $T$ en kelvins}, puis éventuellement convertir en degrés Celsius : $\theta(°\text{C}) = T(\text{K}) - 273$.
\item \textbf{Interpréter} : plus le corps est chaud, plus $\lambda_{\max}$ est petit (déplacement vers le bleu/UV) ; un corps moins chaud émet surtout dans le rouge ou l'infrarouge.
\end{enumerate}
\end{methode}

\begin{attention}[Conditions d'application de la loi de Wien]
La loi $\lambda_{\max} T = \num{2,90e-3}~\text{m.K}$ ne s'applique qu'au profil continu d'un corps chaud dense (corps noir : filament, surface stellaire, métal incandescent, lave), \textbf{jamais} à une raie isolée d'un spectre de gaz dilué. Ne jamais écrire $\lambda_{\text{raie}} \times T = b$.
\end{attention}

\begin{exoresolu}[Température d'un filament et d'une étoile]
\textbf{Énoncé.} On donne $b = \num{2,90e-3}~\text{m.K}$.
\begin{enumerate}
\item Le filament d'une lampe halogène émet un rayonnement dont le maximum se situe à $\lambda_{\max} = \num{1,0e-6}~\text{m}$. Calculer sa température en kelvins puis en degrés Celsius. Dans quel domaine se situe ce maximum ?
\item Une étoile bleutée a un maximum d'émission à $\lambda'_{\max} = 290~\text{nm}$. Comparer sa température à celle du Soleil ($\num{5800}~\text{K}$).
\end{enumerate}
\tcblower
\textbf{Solution.}
\begin{enumerate}
\item Loi de Wien : $T = \dfrac{b}{\lambda_{\max}} = \dfrac{\num{2,90e-3}}{\num{1,0e-6}} = \num{2,90e3}~\text{K}$, soit $\theta = 2900 - 273 \approx \num{2630}~°\text{C}$.\\
Le maximum se situe à $\num{1,0}~\mu\text{m} = 1000~\text{nm} > 800~\text{nm}$ : c'est le domaine \textbf{infrarouge}. Une grande partie du rayonnement du filament est donc invisible (chaleur), ce qui explique le faible rendement lumineux des lampes à incandescence/halogène (majoritairement chauffantes).
\item $T' = \dfrac{\num{2,90e-3}}{\num{2,90e-7}} = \num{1,00e4}~\text{K}$.\\
Comparaison : $\dfrac{T'}{T_{\text{Soleil}}} = \dfrac{\num{10000}}{\num{5800}} \approx 1,7$ : l'étoile est environ $1,7$ fois plus chaude que le Soleil. Son pic dans l'UV proche explique sa couleur bleue (pic visible décalé vers le bleu).
\end{enumerate}
\textbf{Conclusion :} le filament atteint environ $\num{2900}~\text{K}$ avec un maximum dans l'infrarouge, tandis que l'étoile bleue, à $\num{10000}~\text{K}$, émet surtout dans l'ultraviolet proche : la position du maximum spectral est un véritable « thermomètre » pour les corps chauds, des lampes aux étoiles.
\end{exoresolu}

\begin{aretenir}[Bilan des relations clés du chapitre]
\begin{center}
\begin{tabularx}{\linewidth}{|l|X|X|}\hline
\textbf{Grandeur / Loi} & \textbf{Formule} & \textbf{Unités SI / Remarques} \\ \hline
Énergie photon & $E = h\nu = \dfrac{hc}{\lambda}$ & $E$ en J, $h=\num{6,63e-34}$ J.s, $\nu$ en Hz, $\lambda$ en m, $c=\num{3,00e8}$ m/s \\
Conversion eV & $1~\text{eV} = \num{1,60e-19}~\text{J}$ & $E(\text{eV}) = \dfrac{E(\text{J})}{\num{1,60e-19}}$ \\
Raccourci eV-nm & $hc \approx 1240~\text{eV.nm}$ & $E(\text{eV}) = \dfrac{1240}{\lambda(\text{nm})}$ \\
Échange lumière-matière & $\Delta E = h\nu$ & $\Delta E = |E_f - E_i|$ (quantification) \\
Loi de Wien (corps noir) & $\lambda_{\max} T = \num{2,90e-3}~\text{m.K}$ & $\lambda_{\max}$ en m, $T$ en K. Pic du spectre \textbf{continu} uniquement. \\ \hline
\end{tabularx}
\end{center}
\end{aretenir}

\begin{attention}[Les 5 erreurs qui coûtent des points au Probatoire]
\begin{enumerate}
\item \textbf{Oublier les conversions avant le calcul.} La longueur d'onde doit être en \textbf{mètres} ($\lambda = 650~\text{nm} = \num{6,50e-7}~\text{m}$) et l'énergie en \textbf{joules} quand on utilise $h = \num{6,63e-34}~\text{J.s}$. Calculer $E = h\nu$ avec $\Delta E$ resté en eV donne un résultat faux d'un facteur $\num{1,6e-19}$ : convertir d'abord, calculer ensuite.
\item \textbf{Confondre émission et absorption.} Émission : l'atome descend vers un niveau \emph{inférieur} et \emph{cède} un photon ($h\nu = E_i - E_f > 0$). Absorption : l'atome monte et \emph{reçoit} un photon ($h\nu = E_f - E_i > 0$). Le sens de la flèche sur le diagramme et le verbe (émet/absorbe) doivent être cohérents et explicités.
\item \textbf{Croire qu'un photon d'énergie « suffisante » est toujours absorbé.} Faux : hors ionisation ($E \geq |E_1|$), l'absorption n'a lieu que si l'énergie du photon est \emph{exactement égale} à un écart entre deux niveaux. C'est la quantification : un photon de $6,0~\text{eV}$ ne peut pas provoquer une transition de $4,9~\text{eV}$ en « laissant la monnaie ».
\item \textbf{Inverser la relation entre énergie et longueur d'onde.} $E = \dfrac{hc}{\lambda}$ : grande longueur d'onde $\Leftrightarrow$ photon \emph{peu} énergétique. Le rouge ($650~\text{nm}$, $\approx 1,9~\text{eV}$) est moins énergétique que le bleu ($450~\text{nm}$, $\approx 2,8~\text{eV}$). Vérifier systématiquement la cohérence du résultat avec cette tendance physique.
\item \textbf{Appliquer la loi de Wien à tort et négliger les unités de température.} La loi $\lambda_{\max} T = \num{2,90e-3}~\text{m.K}$ ne s'applique qu'au profil continu d'un corps chaud (corps noir), jamais à une raie isolée ; la température s'obtient en \textbf{kelvins}, et il faut écrire la conversion $\theta = T - 273$ si la réponse est demandée en degrés Celsius. Enfin, ne pas oublier la phrase de conclusion avec l'unité et l'interprétation physique (domaine spectral, couleur, température).
\end{enumerate}
\end{attention}

\newpage

\coursec{Exercices}

\courssub{Je vérifie mes connaissances}

\begin{exercice}[QCM : le photon]
Pour chaque affirmation, choisir la (ou les) bonne(s) réponse(s).
\begin{enumerate}
\item L'énergie d'un photon de fréquence $\nu$ vaut :
\begin{enumerate}
\item[a)] $E = \dfrac{h}{\nu}$ \quad b)] $E = h\nu$ \quad c)] $E = \dfrac{hc}{\nu}$
\end{enumerate}
\item Quand la fréquence d'une radiation double, l'énergie de ses photons :
\begin{enumerate}
\item[a)] est divisée par 2 \quad b)] reste inchangée \quad c)] double
\end{enumerate}
\item Un photon ultraviolet est, par rapport à un photon infrarouge :
\begin{enumerate}
\item[a)] moins énergétique \quad b)] plus énergétique \quad c)] de même énergie
\end{enumerate}
\item Le spectre d'une lampe à vapeur de mercure est :
\begin{enumerate}
\item[a)] continu \quad b)] un spectre de raies d'émission \quad c)] un spectre de bandes
\end{enumerate}
\end{enumerate}
\end{exercice}

\begin{exercice}[Vrai ou faux ?]
Répondre par vrai ou faux en justifiant chaque réponse par une phrase.
\begin{enumerate}
\item Les niveaux d'énergie d'un atome peuvent prendre n'importe quelle valeur.
\item Un atome ne peut absorber un photon que si l'énergie de celui-ci correspond exactement à la différence d'énergie entre deux de ses niveaux.
\item La lumière blanche du Soleil, traversant l'atmosphère solaire, donne un spectre présentant des raies sombres.
\item La constante de Planck $h$ s'exprime en joules par seconde ($\mathrm{J\,s^{-1}}$).
\end{enumerate}
\end{exercice}

\begin{exercice}[Définitions]
Définir brièvement les termes suivants : \cle{photon} ; \cle{quantification} des niveaux d'énergie ; \cle{profil spectral} d'une source lumineuse ; \cle{raie d'absorption}.
\end{exercice}

\begin{exercice}[Calculs directs]
On donne $h = \num{6,63e-34}\ \mathrm{J\,s}$, $c = \num{3,00e8}\ \mathrm{m\,s^{-1}}$, $1\ \mathrm{eV} = \num{1,60e-19}\ \mathrm{J}$.
\begin{enumerate}
\item Calculer la fréquence $\nu$ d'une radiation de longueur d'onde $\lambda = 500\ \mathrm{nm}$.
\item En déduire l'énergie du photon correspondant, en joules puis en électronvolts.
\item Quelle est la longueur d'onde d'un photon d'énergie $E = 3{,}0\ \mathrm{eV}$ ? À quel domaine du spectre cette radiation appartient-elle ?
\end{enumerate}
\end{exercice}

\courssub{Je m'entraîne}

\begin{exercice}[Lumière rouge et ultraviolets]
On donne $h = \num{6,63e-34}\ \mathrm{J\,s}$, $c = \num{3,00e8}\ \mathrm{m\,s^{-1}}$, $1\ \mathrm{eV} = \num{1,60e-19}\ \mathrm{J}$.
\begin{enumerate}
\item Calculer l'énergie (en J puis en eV) d'un photon de lumière rouge de longueur d'onde $\lambda_1 = 700\ \mathrm{nm}$.
\item Calculer l'énergie (en J puis en eV) d'un photon ultraviolet de longueur d'onde $\lambda_2 = 300\ \mathrm{nm}$.
\item Comparer ces deux énergies. Expliquer alors pourquoi les UV provoquent des coups de soleil (altération des molécules de la peau) alors que la lumière rouge, même intense, ne le fait pas.
\end{enumerate}
\end{exercice}

\begin{exercice}[Diagramme de niveaux de l'atome de mercure]
On donne quelques niveaux d'énergie de l'atome de mercure : $E_0 = -10{,}44\ \mathrm{eV}$ (fondamental), $E_1 = -5{,}55\ \mathrm{eV}$, $E_2 = -5{,}52\ \mathrm{eV}$, $E_3 = -3{,}74\ \mathrm{eV}$. On donne $h = \num{6,63e-34}\ \mathrm{J\,s}$, $c = \num{3,00e8}\ \mathrm{m\,s^{-1}}$, $1\ \mathrm{eV} = \num{1,60e-19}\ \mathrm{J}$.
\begin{enumerate}
\item Représenter schématiquement ce diagramme de niveaux d'énergie.
\item L'atome, porté sur un niveau excité, se désexcite directement vers le niveau fondamental. Calculer les longueurs d'onde des trois radiations susceptibles d'être émises.
\item Préciser le domaine spectral (UV, visible, IR) de chacune de ces radiations.
\end{enumerate}
\end{exercice}

\begin{exercice}[Absorption d'un photon par le sodium]
Le diagramme simplifié de l'atome de sodium comporte les niveaux : $E_0 = -5{,}14\ \mathrm{eV}$ (fondamental), $E_1 = -3{,}03\ \mathrm{eV}$, $E_2 = -1{,}94\ \mathrm{eV}$. Un photon d'énergie $E = 2{,}11\ \mathrm{eV}$ frappe un atome de sodium initialement au repos (niveau fondamental). On donne $h = \num{6,63e-34}\ \mathrm{J\,s}$, $c = \num{3,00e8}\ \mathrm{m\,s^{-1}}$.
\begin{enumerate}
\item Calculer les différences d'énergie $E_1 - E_0$ et $E_2 - E_0$.
\item Le photon peut-il être absorbé par l'atome ? Justifier. Si oui, préciser la transition correspondante.
\item Calculer la longueur d'onde associée à cette transition. À quelle couleur correspond-elle ?
\end{enumerate}
\end{exercice}

\begin{exercice}[Température d'un filament]
Le profil spectral d'une lampe à incandescence présente un maximum d'émission à $\lambda_{\max} = 966\ \mathrm{nm}$. On admet la loi de Wien : $\lambda_{\max} \cdot T = \num{2,90e-3}\ \mathrm{m\,K}$.
\begin{enumerate}
\item Calculer la température $T$ du filament en kelvins, puis en degrés Celsius.
\item Dans quel domaine spectral se situe $\lambda_{\max}$ ?
\item Expliquer pourquoi une grande partie de l'énergie consommée par une telle lampe est perdue sous forme de rayonnement infrarouge (chaleur) plutôt que de lumière visible.
\end{enumerate}
\end{exercice}

\courssub{Je me prépare au Probatoire}

\begin{exercice}[La lampe à vapeur de sodium de l'éclairage public]
Dans plusieurs villes du Cameroun, l'éclairage public utilise des lampes à vapeur de sodium qui émettent une lumière jaune caractéristique de longueur d'onde $\lambda = 589\ \mathrm{nm}$. Le diagramme simplifié de l'atome de sodium donne : $E_0 = -5{,}14\ \mathrm{eV}$ (fondamental), $E_1 = -3{,}03\ \mathrm{eV}$, $E_2 = -1{,}94\ \mathrm{eV}$, $E_3 = -1{,}52\ \mathrm{eV}$.\\
Données : $h = \num{6,63e-34}\ \mathrm{J\,s}$, $c = \num{3,00e8}\ \mathrm{m\,s^{-1}}$, $1\ \mathrm{eV} = \num{1,60e-19}\ \mathrm{J}$.
\begin{enumerate}
\item Rappeler la relation donnant l'énergie d'un photon en fonction de la fréquence, puis celle reliant longueur d'onde et fréquence.
\item Calculer, en joules puis en électronvolts, l'énergie d'un photon de la radiation jaune émise.
\item Montrer que cette radiation correspond à une transition entre deux niveaux du diagramme ; préciser lesquels et représenter cette transition par une flèche.
\item Un photon d'énergie $3{,}20\ \mathrm{eV}$ arrive sur un atome de sodium au niveau fondamental. Peut-il être absorbé ? Justifier.
\item Expliquer, à l'aide du modèle du photon et de la quantification des niveaux d'énergie, pourquoi le spectre de cette lampe est un spectre de raies discontinu, alors que celui d'une lampe à incandescence est continu.
\item (Hors barème classique) La puissance lumineuse de la lampe est $P = 10\ \mathrm{W}$. En supposant que toute cette puissance est rayonnée sous forme de photons de longueur d'onde $589\ \mathrm{nm}$, calculer le nombre de photons émis par seconde.
\end{enumerate}
\end{exercice}

\begin{exercice}[Étude spectrale d'une étoile]
Un astrophysicien amateur de l'observatoire de l'université de Yaoundé étudie la lumière d'une étoile. Le profil spectral de cette étoile présente un maximum d'émission à $\lambda_{\max} = 580\ \mathrm{nm}$, et son spectre montre des raies d'absorption situées exactement aux longueurs d'onde des raies d'émission de l'hydrogène mesurées en laboratoire.\\
Données : loi de Wien $\lambda_{\max} \cdot T = \num{2,90e-3}\ \mathrm{m\,K}$ ; température de surface du Soleil $T_S \approx \num{5800}\ \mathrm{K}$ ; $h = \num{6,63e-34}\ \mathrm{J\,s}$, $c = \num{3,00e8}\ \mathrm{m\,s^{-1}}$, $1\ \mathrm{eV} = \num{1,60e-19}\ \mathrm{J}$.
\begin{enumerate}
\item Décrire l'allure générale du spectre d'une étoile : fond continu et raies. Préciser l'origine de chacune de ces deux composantes.
\item Calculer la température de surface $T$ de l'étoile en kelvins.
\item Comparer $T$ à la température de surface du Soleil. Cette étoile est-elle plus chaude ou plus froide que le Soleil ? Sa teinte est-elle plutôt jaune, rouge ou bleutée ? Justifier.
\item Expliquer, à l'aide du modèle corpusculaire de la lumière et de la quantification des niveaux d'énergie, pourquoi la présence de raies d'absorption aux longueurs d'onde de l'hydrogène prouve que l'atmosphère de l'étoile contient de l'hydrogène.
\item L'une des raies d'absorption observées correspond à la transition de l'hydrogène d'énergie $\Delta E = 1{,}89\ \mathrm{eV}$. Calculer la longueur d'onde de cette raie. Appartient-elle au domaine visible ?
\end{enumerate}
\end{exercice}

\begin{exercice}[Le spectre du Soleil et l'énergie solaire au Cameroun]
Le Cameroun dispose d'un fort potentiel solaire, exploité par des centrales photovoltaïques en zone rurale. La lumière du Soleil, analysée à l'aide d'un spectroscope, donne un spectre continu traversé par de fines raies sombres (raies de Fraunhofer). Le maximum d'émission du profil spectral solaire se situe à $\lambda_{\max} \approx 500\ \mathrm{nm}$.\\
Données : $\lambda_{\max} \cdot T = \num{2,90e-3}\ \mathrm{m\,K}$ ; $h = \num{6,63e-34}\ \mathrm{J\,s}$, $c = \num{3,00e8}\ \mathrm{m\,s^{-1}}$, $1\ \mathrm{eV} = \num{1,60e-19}\ \mathrm{J}$.
\begin{enumerate}
\item Expliquer qualitativement la formation du spectre solaire : pourquoi observe-t-on un fond continu ? Pourquoi des raies sombres apparaissent-elles ?
\item Calculer la température de surface du Soleil déduite de la loi de Wien.
\item Calculer l'énergie, en joules puis en électronvolts, d'un photon de longueur d'onde $\lambda_{\max} = 500\ \mathrm{nm}$.
\item Les cellules photovoltaïques au silicium ne convertissent en électricité que les photons d'énergie supérieure à $E_g = 1{,}12\ \mathrm{eV}$. Montrer qu'un photon de $\lambda = 500\ \mathrm{nm}$ est utilisable, puis déterminer la longueur d'onde maximale $\lambda_g$ des photons encore exploitables par la cellule.
\item Les photons infrarouges de longueur d'onde supérieure à $\lambda_g$ ne produisent pas d'électricité : proposer une explication à l'aide du modèle du photon, et conclure sur une des causes de la limitation du rendement des panneaux solaires.
\end{enumerate}
\end{exercice}

\newpage

\coursec{Corrigés des exercices}

\courssub{Je vérifie mes connaissances}

\begin{corrige}[Exercice 1 — QCM : le photon]
\begin{enumerate}
\item \textbf{Réponse b)} : $E = h\nu$. L'énergie d'un photon est proportionnelle à la fréquence $\nu$ de la radiation, le coefficient de proportionnalité étant la constante de Planck $h$. Vérification dimensionnelle : $[h\nu] = \mathrm{J\,s}\times\mathrm{s^{-1}} = \mathrm{J}$, c'est bien une énergie. La réponse c) est fausse car $\dfrac{hc}{\nu}$ n'est pas homogène à une énergie : $\left[\dfrac{hc}{\nu}\right] = \dfrac{\mathrm{J\,s}\times\mathrm{m\,s^{-1}}}{\mathrm{s^{-1}}} = \mathrm{J\,m}$.
\item \textbf{Réponse c)} : comme $E = h\nu$, l'énergie est proportionnelle à la fréquence ; si $\nu$ double, $E$ double.
\item \textbf{Réponse b)} : les ultraviolets ont une fréquence plus grande que les infrarouges ($\nu_{\mathrm{UV}} > \nu_{\mathrm{IR}}$), donc $E_{\mathrm{UV}} > E_{\mathrm{IR}}$ : le photon UV est plus énergétique.
\item \textbf{Réponse b)} : une lampe à vapeur de mercure contient un gaz d'atomes sous faible pression ; les atomes excités se désexcitent en émettant des photons d'énergies bien précises : le spectre est un spectre de \cle{raies d'émission} (raies brillantes sur fond noir).
\end{enumerate}
\end{corrige}

\begin{corrige}[Exercice 2 — Vrai ou faux ?]
\begin{enumerate}
\item \textbf{Faux.} Les niveaux d'énergie d'un atome sont \cle{quantifiés} : ils ne peuvent prendre que certaines valeurs discrètes, caractéristiques de l'atome.
\item \textbf{Vrai.} L'absorption n'est possible que si $E_{\text{photon}} = h\nu = |E_f - E_i|$, c'est-à-dire si l'énergie du photon correspond \emph{exactement} à un écart entre deux niveaux ; sinon le photon n'interagit pas avec l'atome (il est transmis sans absorption résonnante).
\item \textbf{Vrai.} Le fond continu émis par la photosphère traverse la chromosphère (atmosphère solaire) dont les atomes absorbent les photons d'énergies égales à leurs écarts de niveaux : il apparaît des raies sombres (raies de Fraunhofer) sur le fond continu.
\item \textbf{Faux.} $h$ s'exprime en joules fois secondes : $[h] = \dfrac{[E]}{[\nu]} = \dfrac{\mathrm{J}}{\mathrm{s^{-1}}} = \mathrm{J\,s}$, et non en $\mathrm{J\,s^{-1}}$.
\end{enumerate}
\end{corrige}

\begin{corrige}[Exercice 3 — Définitions]
\begin{itemize}
\item \cle{Photon} : particule élémentaire, grain d'énergie associé à une radiation électromagnétique dans le modèle corpusculaire de la lumière ; il se déplace à la célérité $c$ dans le vide, a une masse nulle au repos et transporte l'énergie $E = h\nu$.
\item \cle{Quantification} des niveaux d'énergie : propriété de la matière (atomes, molécules) de ne pouvoir exister que dans des états d'énergie discrètes, bien déterminées, et non dans un continuum de valeurs.
\item \cle{Profil spectral} d'une source lumineuse : courbe donnant l'intensité lumineuse (ou la luminance spectrale) émise en fonction de la longueur d'onde ; il renseigne sur la nature et la température de la source.
\item \cle{Raie d'absorption} : raie sombre observée dans le spectre d'une lumière blanche ayant traversé un gaz, correspondant à la longueur d'onde des photons absorbés par les atomes du gaz (transition d'un niveau d'énergie vers un niveau supérieur).
\end{itemize}
\end{corrige}

\begin{corrige}[Exercice 4 — Calculs directs]
\begin{enumerate}
\item Fréquence de la radiation :
\[ \nu = \frac{c}{\lambda} = \frac{3,00\times10^8}{500\times10^{-9}} = 6,00\times10^{14}\ \mathrm{Hz}. \]
\item Énergie du photon :
\[ E = h\nu = 6,63\times10^{-34}\times6,00\times10^{14} = 3,98\times10^{-19}\ \mathrm{J}. \]
Conversion en électronvolts :
\[ E = \frac{3,98\times10^{-19}}{1,60\times10^{-19}} = 2,49\ \mathrm{eV}. \]
\item Pour $E = 3,0\ \mathrm{eV} = 3,0\times1,60\times10^{-19} = 4,80\times10^{-19}\ \mathrm{J}$ :
\[ \lambda = \frac{hc}{E} = \frac{6,63\times10^{-34}\times3,00\times10^8}{4,80\times10^{-19}} = 4,14\times10^{-7}\ \mathrm{m} = 414\ \mathrm{nm}. \]
Cette radiation appartient au domaine \textbf{visible} (limite du violet, proche de l'ultraviolet).
\end{enumerate}
\end{corrige}

\courssub{Je m'entraîne}

\begin{corrige}[Exercice 5 — Lumière rouge et ultraviolets]
\begin{enumerate}
\item Photon rouge ($\lambda_1 = 700\ \mathrm{nm}$) :
\[ E_1 = \frac{hc}{\lambda_1} = \frac{6,63\times10^{-34}\times3,00\times10^8}{700\times10^{-9}} = 2,84\times10^{-19}\ \mathrm{J} = \frac{2,84\times10^{-19}}{1,60\times10^{-19}} = 1,78\ \mathrm{eV}. \]
\item Photon ultraviolet ($\lambda_2 = 300\ \mathrm{nm}$) :
\[ E_2 = \frac{hc}{\lambda_2} = \frac{6,63\times10^{-34}\times3,00\times10^8}{300\times10^{-9}} = 6,63\times10^{-19}\ \mathrm{J} = 4,14\ \mathrm{eV}. \]
\item Comparaison : $\dfrac{E_2}{E_1} = \dfrac{\lambda_1}{\lambda_2} = \dfrac{700}{300} \approx 2,3$ : un photon UV est environ $2,3$ fois plus énergétique qu'un photon rouge. Les échanges d'énergie entre lumière et matière se font \emph{photon par photon} : un photon UV de $4,14\ \mathrm{eV}$ transporte assez d'énergie pour rompre des liaisons chimiques et altérer les molécules de la peau (coups de soleil). Un photon rouge de $1,78\ \mathrm{eV}$, même s'ils arrivent en grand nombre (lumière intense), ne peut individuellement fournir l'énergie nécessaire à ces altérations : c'est l'énergie de \emph{chaque} photon qui compte, pas seulement l'énergie totale du faisceau.
\end{enumerate}
\end{corrige}

\begin{corrige}[Exercice 6 — Diagramme de niveaux de l'atome de mercure]
\begin{enumerate}
\item Diagramme de niveaux (échelle non respectée) :
\begin{popfigure}
\begin{tikzpicture}[scale=1]
\draw[thick,popdark] (0,0) -- (4,0) node[right]{$E_0 = -10,44\ \mathrm{eV}$ (fondamental)};
\draw[thick,popblue] (0,1.8) -- (4,1.8) node[right]{$E_1 = -5,55\ \mathrm{eV}$};
\draw[thick,popblue] (0,2.0) -- (4,2.0) node[right]{$E_2 = -5,52\ \mathrm{eV}$};
\draw[thick,poporange] (0,3.4) -- (4,3.4) node[right]{$E_3 = -3,74\ \mathrm{eV}$};
\draw[-{Stealth},very thick,popgreen] (1,1.8) -- (1,0.05);
\draw[-{Stealth},very thick,poppurple] (2,2.0) -- (2,0.05);
\draw[-{Stealth},very thick,poporange] (3,3.4) -- (3,0.05);
\end{tikzpicture}
\legende{Transitions de désexcitation directe vers le niveau fondamental.}
\end{popfigure}
\item Lors d'une désexcitation de $E_n$ vers $E_0$, le photon émis a pour énergie $\Delta E = E_n - E_0$, et $\lambda = \dfrac{hc}{\Delta E}$ avec $hc = 1,989\times10^{-25}\ \mathrm{J\,m}$.
\begin{itemize}
\item Transition $E_1 \to E_0$ : $\Delta E_1 = -5,55 - (-10,44) = 4,89\ \mathrm{eV} = 7,82\times10^{-19}\ \mathrm{J}$ :
\[ \lambda_1 = \frac{1,989\times10^{-25}}{7,82\times10^{-19}} = 2,54\times10^{-7}\ \mathrm{m} = 254\ \mathrm{nm}. \]
\item Transition $E_2 \to E_0$ : $\Delta E_2 = -5,52 - (-10,44) = 4,92\ \mathrm{eV} = 7,87\times10^{-19}\ \mathrm{J}$ :
\[ \lambda_2 = \frac{1,989\times10^{-25}}{7,87\times10^{-19}} = 2,53\times10^{-7}\ \mathrm{m} = 253\ \mathrm{nm}. \]
\item Transition $E_3 \to E_0$ : $\Delta E_3 = -3,74 - (-10,44) = 6,70\ \mathrm{eV} = 1,07\times10^{-18}\ \mathrm{J}$ :
\[ \lambda_3 = \frac{1,989\times10^{-25}}{1,07\times10^{-18}} = 1,86\times10^{-7}\ \mathrm{m} = 186\ \mathrm{nm}. \]
\end{itemize}
\item Les trois longueurs d'onde sont inférieures à $400\ \mathrm{nm}$ : les trois radiations appartiennent au domaine \textbf{ultraviolet}. (C'est pourquoi les lampes à vapeur de mercure servent de source UV, par exemple pour la stérilisation.)
\end{enumerate}
\end{corrige}

\begin{corrige}[Exercice 7 — Absorption d'un photon par le sodium]
\begin{enumerate}
\item Différences d'énergie depuis le fondamental :
\[ E_1 - E_0 = -3,03 - (-5,14) = 2,11\ \mathrm{eV} ; \qquad E_2 - E_0 = -1,94 - (-5,14) = 3,20\ \mathrm{eV}. \]
\item Le photon incident a une énergie $E = 2,11\ \mathrm{eV}$, exactement égale à $E_1 - E_0$. La condition d'absorption $E_{\text{photon}} = |E_f - E_i|$ est donc satisfaite : \textbf{le photon est absorbé}, et l'atome effectue la transition du niveau fondamental $E_0$ vers le niveau excité $E_1$.
\item Longueur d'onde associée :
\[ \lambda = \frac{hc}{E_1 - E_0} = \frac{6,63\times10^{-34}\times3,00\times10^8}{2,11\times1,60\times10^{-19}} = \frac{1,989\times10^{-25}}{3,376\times10^{-19}} = 5,89\times10^{-7}\ \mathrm{m} = 589\ \mathrm{nm}. \]
Cette radiation correspond à la couleur \textbf{jaune} (c'est la célèbre raie jaune du sodium, utilisée dans l'éclairage public).
\end{enumerate}
\end{corrige}

\begin{corrige}[Exercice 8 — Température d'un filament]
\begin{enumerate}
\item D'après la loi de Wien :
\[ T = \frac{2,90\times10^{-3}}{\lambda_{\max}} = \frac{2,90\times10^{-3}}{966\times10^{-9}} = 3,00\times10^{3}\ \mathrm{K} \approx 3000\ \mathrm{K}, \]
soit $t = T - 273 \approx 2727\ ^\circ\mathrm{C}$.
\item $\lambda_{\max} = 966\ \mathrm{nm} > 800\ \mathrm{nm}$ : le maximum d'émission se situe dans le domaine \textbf{infrarouge}.
\item Le profil spectral d'un corps chaud est une courbe en cloche centrée sur $\lambda_{\max}$. Ici, ce maximum étant dans l'infrarouge, la plus grande partie de la puissance rayonnée par le filament l'est sous forme de photons infrarouges (chaleur), et seule une faible fraction (la « queue » de la courbe côté visible) correspond à de la lumière visible. D'où le très faible rendement lumineux des lampes à incandescence, remplacées au Cameroun comme ailleurs par les lampes fluocompactes et les LED.
\end{enumerate}
\end{corrige}

\courssub{Je me prépare au Probatoire}

\begin{corrige}[Exercice 9 — La lampe à vapeur de sodium de l'éclairage public]
\begin{enumerate}
\item Énergie d'un photon de fréquence $\nu$ : $E = h\nu$ ; relation entre longueur d'onde et fréquence : $\lambda = \dfrac{c}{\nu}$. On en déduit $E = \dfrac{hc}{\lambda}$.
\item Énergie du photon jaune ($\lambda = 589\ \mathrm{nm}$) :
\[ E = \frac{hc}{\lambda} = \frac{6,63\times10^{-34}\times3,00\times10^8}{589\times10^{-9}} = 3,38\times10^{-19}\ \mathrm{J} = \frac{3,38\times10^{-19}}{1,60\times10^{-19}} = 2,11\ \mathrm{eV}. \]
\item Calculons les écarts entre niveaux : $E_1 - E_0 = -3,03 + 5,14 = 2,11\ \mathrm{eV}$. L'énergie du photon émis est exactement égale à $E_1 - E_0$ : la radiation jaune correspond à la \textbf{transition $E_1 \to E_0$} (désexcitation du premier niveau excité vers le fondamental).
\begin{popfigure}
\begin{tikzpicture}[scale=1]
\draw[thick,popdark] (0,0) -- (4,0) node[right]{$E_0 = -5,14\ \mathrm{eV}$};
\draw[thick,popblue] (0,1.4) -- (4,1.4) node[right]{$E_1 = -3,03\ \mathrm{eV}$};
\draw[thick,popblue] (0,2.2) -- (4,2.2) node[right]{$E_2 = -1,94\ \mathrm{eV}$};
\draw[thick,popblue] (0,2.7) -- (4,2.7) node[right]{$E_3 = -1,52\ \mathrm{eV}$};
\draw[-{Stealth},very thick,popgold] (2,1.4) -- (2,0.05) node[midway,left]{$\lambda = 589\ \mathrm{nm}$};
\end{tikzpicture}
\legende{Transition responsable de la lumière jaune du sodium.}
\end{popfigure}
\item $E_2 - E_0 = -1,94 + 5,14 = 3,20\ \mathrm{eV}$ : le photon d'énergie $3,20\ \mathrm{eV}$ correspond exactement à l'écart $E_2 - E_0$. Il \textbf{peut donc être absorbé}, l'atome passant du niveau fondamental $E_0$ au niveau excité $E_2$.
\item Dans la vapeur de sodium, les atomes sont quasiment indépendants et leurs niveaux d'énergie sont quantifiés : ils ne peuvent émettre que des photons d'énergies égales aux écarts entre leurs niveaux, c'est-à-dire des radiations de longueurs d'onde discrètes et bien précises. Le spectre est donc un \textbf{spectre de raies}. Au contraire, dans le filament solide d'une lampe à incandescence, les atomes sont liés et leurs niveaux d'énergie s'élargissent en bandes quasi continues : toutes les longueurs d'onde sont émises, le spectre est \textbf{continu} et dépend essentiellement de la température du filament.
\item Nombre de photons émis par seconde : la puissance est l'énergie rayonnée par unité de temps, $P = N\times E_{\text{photon}} / \Delta t$ avec $\Delta t = 1\ \mathrm{s}$ :
\[ N = \frac{P}{E_{\text{photon}}} = \frac{10}{3,38\times10^{-19}} = 3,0\times10^{19}\ \text{photons par seconde}. \]
\end{enumerate}
\textbf{Barème indicatif (sur 10)} : relations (1 pt) ; calcul de $E$ en J et eV (2 pts) ; identification de la transition avec justification par l'égalité des énergies (2 pts) ; absorption justifiée (1,5 pt) ; explication spectre de raies vs continu (2 pts) ; question 6 (1,5 pt).

\textbf{Points de vigilance} : citer explicitement la condition $\Delta E = h\nu$ pour justifier une émission ou une absorption ; convertir les eV en joules avant d'appliquer $\lambda = hc/\Delta E$ ; flécher la transition dans le bon sens (vers le bas pour une émission) ; ne pas confondre puissance (W) et énergie (J) dans la question 6.
\end{corrige}

\begin{corrige}[Exercice 10 — Étude spectrale d'une étoile]
\begin{enumerate}
\item Le spectre d'une étoile est constitué d'un \textbf{fond continu} (courbe en cloche) surmonté de \textbf{raies sombres} (raies d'absorption). Le fond continu est émis par la surface chaude de l'étoile (la photosphère), dont le profil spectral dépend de sa température. Les raies sombres sont produites lorsque cette lumière traverse l'atmosphère de l'étoile, plus froide : les atomes qu'elle contient absorbent les photons dont l'énergie égale exactement l'écart entre deux de leurs niveaux.
\item Loi de Wien :
\[ T = \frac{2,90\times10^{-3}}{\lambda_{\max}} = \frac{2,90\times10^{-3}}{580\times10^{-9}} = 5,0\times10^{3}\ \mathrm{K} = 5000\ \mathrm{K}. \]
\item $T = 5000\ \mathrm{K} < T_S = 5800\ \mathrm{K}$ : l'étoile est \textbf{plus froide} que le Soleil. Son maximum d'émission ($580\ \mathrm{nm}$) est décalé vers les grandes longueurs d'onde par rapport à celui du Soleil ($500\ \mathrm{nm}$) : sa teinte est donc plutôt \textbf{orangée/rouge} (alors qu'une étoile plus chaude aurait un maximum décalé vers le bleu et paraîtrait bleutée).
\item Les niveaux d'énergie de l'atome d'hydrogène sont quantifiés et caractéristiques de cet élément : il ne peut absorber que des photons d'énergies $h\nu = |E_f - E_i|$, donc des longueurs d'onde précises, exactement celles des raies qu'il émet en laboratoire. Observer des raies d'absorption aux longueurs d'onde propres de l'hydrogène prouve que des atomes d'hydrogène, présents dans l'atmosphère de l'étoile, ont absorbé ces photons : c'est la signature de la présence d'hydrogène.
\item Longueur d'onde de la raie :
\[ \lambda = \frac{hc}{\Delta E} = \frac{6,63\times10^{-34}\times3,00\times10^8}{1,89\times1,60\times10^{-19}} = \frac{1,989\times10^{-25}}{3,024\times10^{-19}} = 6,58\times10^{-7}\ \mathrm{m} = 658\ \mathrm{nm}. \]
$400\ \mathrm{nm} < 658\ \mathrm{nm} < 800\ \mathrm{nm}$ : la raie appartient au domaine \textbf{visible} (c'est la raie rouge H$\alpha$ de l'hydrogène).
\end{enumerate}
\textbf{Barème indicatif (sur 10)} : description du spectre et origine des deux composantes (2 pts) ; calcul de $T$ (1,5 pt) ; comparaison et teinte justifiées (2 pts) ; raisonnement sur la quantification et la signature de l'hydrogène (2,5 pts) ; calcul de $\lambda$ et domaine (2 pts).

\textbf{Points de vigilance} : dans la loi de Wien, $\lambda_{\max}$ doit être converti en mètres ; la teinte doit être justifiée par le décalage de $\lambda_{\max}$, pas affirmée gratuitement ; pour la question 4, il faut impérativement mentionner que les écarts de niveaux sont \emph{caractéristiques} de l'élément (signature spectrale).
\end{corrige}

\begin{corrige}[Exercice 11 — Le spectre du Soleil et l'énergie solaire au Cameroun]
\begin{enumerate}
\item La photosphère solaire, surface chaude et dense, se comporte comme un corps chaud : elle émet un rayonnement au \textbf{profil continu} dont le maximum dépend de sa température. Cette lumière traverse ensuite la chromosphère (atmosphère solaire, plus froide) : les atomes qu'elle contient absorbent les photons dont l'énergie $h\nu$ correspond exactement à un écart entre deux de leurs niveaux d'énergie. Ces longueurs d'onde manquent alors dans la lumière reçue : on observe des \textbf{raies sombres} (raies de Fraunhofer) sur le fond continu.
\item Température de surface :
\[ T = \frac{2,90\times10^{-3}}{\lambda_{\max}} = \frac{2,90\times10^{-3}}{500\times10^{-9}} = 5,8\times10^{3}\ \mathrm{K} = 5800\ \mathrm{K}. \]
\item Énergie d'un photon de $\lambda = 500\ \mathrm{nm}$ :
\[ E = \frac{hc}{\lambda} = \frac{6,63\times10^{-34}\times3,00\times10^8}{500\times10^{-9}} = 3,98\times10^{-19}\ \mathrm{J} = \frac{3,98\times10^{-19}}{1,60\times10^{-19}} = 2,49\ \mathrm{eV}. \]
\item Comparaison : $E = 2,49\ \mathrm{eV} > E_g = 1,12\ \mathrm{eV}$ : le photon de $500\ \mathrm{nm}$ possède une énergie supérieure au seuil, il est donc \textbf{utilisable} par la cellule au silicium. La longueur d'onde maximale exploitable correspond au photon d'énergie exactement égale à $E_g$ :
\[ \lambda_g = \frac{hc}{E_g} = \frac{1,989\times10^{-25}}{1,12\times1,60\times10^{-19}} = \frac{1,989\times10^{-25}}{1,792\times10^{-19}} = 1,11\times10^{-6}\ \mathrm{m} = 1110\ \mathrm{nm}. \]
\item L'énergie d'un photon décroît quand la longueur d'onde croît ($E = hc/\lambda$). Un photon infrarouge de $\lambda > \lambda_g$ a une énergie $E < E_g$ : or les échanges d'énergie se font photon par photon, et l'énergie d'un seul photon est alors insuffisante pour libérer un électron dans le silicium, quelle que soit l'intensité du rayonnement. Une fraction importante du rayonnement solaire (tout l'infrarouge au-delà de $1110\ \mathrm{nm}$) ne peut donc pas être convertie en électricité : c'est l'une des causes majeures de la limitation du rendement des panneaux photovoltaïques installés dans les zones rurales du Cameroun.
\end{enumerate}
\textbf{Barème indicatif (sur 10)} : formation du spectre (fond continu + raies) (2 pts) ; température (1,5 pt) ; énergie du photon en J et eV (2 pts) ; comparaison à $E_g$ et calcul de $\lambda_g$ (2,5 pts) ; explication par le modèle du photon et conclusion sur le rendement (2 pts).

\textbf{Points de vigilance} : bien distinguer les deux composantes du spectre solaire et leurs origines respectives (photosphère / chromosphère) ; pour le seuil, raisonner sur l'énergie d'\emph{un} photon et non sur la puissance totale ; vérifier l'homogénéité de $\lambda_g = hc/E_g$ ($\mathrm{J\,s}\times\mathrm{m\,s^{-1}}/\mathrm{J} = \mathrm{m}$) ; conclure en reliant explicitement le résultat à la limitation du rendement.
\end{corrige}

\newpage

\coursec{Fiche bilan}

\begin{popfigure}
\begin{center}\tcbox[colback=popblue,colframe=popblue,coltext=white,arc=9pt,boxsep=4pt,left=6pt,right=6pt]{\bfseries\small Interaction lumi\`ere--mati\`ere : le photon}\end{center}
\vspace{-2pt}
\begin{tcbraster}[raster columns=3,raster equal height=rows,raster column skip=6pt,raster row skip=6pt,enhanced,blankest]
\begin{tcolorbox}[enhanced,colback=white,colframe=poporange,boxrule=0.9pt,arc=3pt,title={Interactions},fonttitle=\bfseries\scriptsize,coltitle=white,colbacktitle=poporange,left=4pt,right=4pt,top=2pt,bottom=2pt,boxsep=1pt,toptitle=1.5pt,bottomtitle=1.5pt,halign=left]
\begin{itemize}[nosep,leftmargin=*,itemsep=1pt,font=\scriptsize]
\item \'Emission
\item Absorption
\item Diffusion
\end{itemize}
\end{tcolorbox}
\begin{tcolorbox}[enhanced,colback=white,colframe=popgreen,boxrule=0.9pt,arc=3pt,title={Quantification},fonttitle=\bfseries\scriptsize,coltitle=white,colbacktitle=popgreen,left=4pt,right=4pt,top=2pt,bottom=2pt,boxsep=1pt,toptitle=1.5pt,bottomtitle=1.5pt,halign=left]
\begin{itemize}[nosep,leftmargin=*,itemsep=1pt,font=\scriptsize]
\item Niveaux d'\'energie $E_n$
\item Diagramme d'\'energie
\end{itemize}
\end{tcolorbox}
\begin{tcolorbox}[enhanced,colback=white,colframe=poppurple,boxrule=0.9pt,arc=3pt,title={Le photon},fonttitle=\bfseries\scriptsize,coltitle=white,colbacktitle=poppurple,left=4pt,right=4pt,top=2pt,bottom=2pt,boxsep=1pt,toptitle=1.5pt,bottomtitle=1.5pt,halign=left]
\begin{itemize}[nosep,leftmargin=*,itemsep=1pt,font=\scriptsize]
\item Grain de lumi\`ere, masse nulle
\item Vitesse $c$
\end{itemize}
\end{tcolorbox}
\begin{tcolorbox}[enhanced,colback=white,colframe=popgold,boxrule=0.9pt,arc=3pt,title={Formules cl\'es},fonttitle=\bfseries\scriptsize,coltitle=white,colbacktitle=popgold,left=4pt,right=4pt,top=2pt,bottom=2pt,boxsep=1pt,toptitle=1.5pt,bottomtitle=1.5pt,halign=left]
\begin{itemize}[nosep,leftmargin=*,itemsep=1pt,font=\scriptsize]
\item $E=h\nu$
\item $\lambda=\dfrac{c}{\nu}$
\item $\Delta E=h\nu$
\end{itemize}
\end{tcolorbox}
\begin{tcolorbox}[enhanced,colback=white,colframe=poppink,boxrule=0.9pt,arc=3pt,title={Profil spectral},fonttitle=\bfseries\scriptsize,coltitle=white,colbacktitle=poppink,left=4pt,right=4pt,top=2pt,bottom=2pt,boxsep=1pt,toptitle=1.5pt,bottomtitle=1.5pt,halign=left]
\begin{itemize}[nosep,leftmargin=*,itemsep=1pt,font=\scriptsize]
\item Spectre continu : corps chaud
\item Loi de Wien $\lambda_m T$
\end{itemize}
\end{tcolorbox}
\begin{tcolorbox}[enhanced,colback=white,colframe=popblueL,boxrule=0.9pt,arc=3pt,title={Spectre solaire},fonttitle=\bfseries\scriptsize,coltitle=white,colbacktitle=popblueL,left=4pt,right=4pt,top=2pt,bottom=2pt,boxsep=1pt,toptitle=1.5pt,bottomtitle=1.5pt,halign=left]
\begin{itemize}[nosep,leftmargin=*,itemsep=1pt,font=\scriptsize]
\item Fond continu + raies
\item Raies de Fraunhofer
\end{itemize}
\end{tcolorbox}
\end{tcbraster}

\legende{Carte mentale de la le\c{c}on : les six id\'ees forces de l'interaction lumi\`ere--mati\`ere.}
\end{popfigure}

\begin{bilan}
\textbf{D\'efinitions cl\'es}
\begin{itemize}
  \item \cle{\'Emission} : passage de la mati\`ere d'un niveau d'\'energie $E_p$ \`a un niveau inf\'erieur $E_n$ avec \'emission d'un photon.
  \item \cle{Absorption} : un photon dispara\^it en transf\'erant toute son \'energie \`a un atome qui passe d'un niveau $E_n$ \`a un niveau sup\'erieur $E_p$.
  \item \cle{Diffusion} : renvoi de la lumi\`ere dans toutes les directions par la mati\`ere, sans changement de fr\'equence (diffusion Rayleigh : ciel bleu).
  \item \cle{Photon} : particule de lumi\`ere, de masse nulle, se d\'epla\c{c}ant \`a la c\'el\'erit\'e $c$, transportant l'\'energie $E=h\nu$.
  \item \cle{Quantification} : l'\'energie d'un atome ne peut prendre que des valeurs discr\`etes $E_1, E_2, \dots$ (niveaux d'\'energie).
  \item \cle{Profil spectral} : courbe donnant l'intensit\'e lumineuse \'emise en fonction de la longueur d'onde $\lambda$.
\end{itemize}

\textbf{Formules \`a conna\^itre par c\oe ur}
\begin{center}
\begin{tabularx}{\linewidth}{|l|X|l|}\hline
\rowcolor{popblueL}
\textbf{Relation} & \textbf{Signification} & \textbf{Unit\'es SI} \\ \hline
$E = h\nu$ & \'Energie d'un photon de fr\'equence $\nu$ ; $h=\SI{6,63e-34}{J.s}$ (constante de Planck) & $E$ en J, $\nu$ en Hz \\ \hline
$\lambda = \dfrac{c}{\nu}$ & Lien longueur d'onde--fr\'equence ; $c=\SI{3,00e8}{m.s^{-1}}$ & $\lambda$ en m \\ \hline
$\Delta E = |E_p - E_n| = h\nu = \dfrac{hc}{\lambda}$ & \'Echange d'\'energie entre deux niveaux (\'emission ou absorption) & $\Delta E$ en J \\ \hline
$\SI{1}{eV} = \SI{1,60e-19}{J}$ & Conversion d'unit\'es d'\'energie & -- \\ \hline
$\lambda_m \cdot T \approx \SI{2,9e-3}{m.K}$ & Loi de Wien : temp\'erature d'un corps chaud par analyse spectrale & $T$ en K \\ \hline
\end{tabularx}
\end{center}

\textbf{M\'ethodes express}
\begin{itemize}
  \item \textbf{Exploiter un diagramme de niveaux} : rep\'erer les deux niveaux concern\'es, calculer $\Delta E = |E_p-E_n|$ (en eV puis en J), en d\'eduire $\nu = \Delta E/h$ puis $\lambda = hc/\Delta E$. Si $\lambda \in [400\,\text{nm}\,;\,800\,\text{nm}]$, la radiation est visible.
  \item \textbf{\'Emission ou absorption ?} Si l'atome \emph{descend} de niveau, il \'emet un photon ($\Delta E$ fourni) ; s'il \emph{monte}, il absorbe un photon d'\'energie exactement \'egale \`a $\Delta E$.
  \item \textbf{Temp\'erature d'un corps chaud} : lire $\lambda_m$ (maximum du profil spectral), appliquer $\lambda_m T \approx \SI{2,9e-3}{m.K}$. Plus le corps est chaud, plus $\lambda_m$ est petite (d\'eplacement vers le bleu).
  \item \textbf{Spectre solaire} : fond continu (photosphere chaude, $T \approx \SI{5800}{K}$) travers\'e de raies sombres d'absorption (raies de Fraunhofer) dues aux atomes de la chromosph\`ere ; chaque raie identifie un \'el\'ement chimique.
\end{itemize}
\end{bilan}

\begin{aretenir}[Checklist avant le Probatoire]
\begin{itemize}
  \item[$\square$] Je sais d\'efinir \'emission, absorption et diffusion, et donner un exemple de chacune.
  \item[$\square$] Je sais que les niveaux d'\'energie de la mati\`ere sont quantifi\'es et je sais lire un diagramme de niveaux.
  \item[$\square$] Je conna\^is la valeur de la constante de Planck $h=\SI{6,63e-34}{J.s}$ et celle de $c=\SI{3,00e8}{m.s^{-1}}$.
  \item[$\square$] Je sais \'ecrire et utiliser $E=h\nu$ avec les bonnes unit\'es (J, Hz).
  \item[$\square$] Je sais passer de $\nu$ \`a $\lambda$ gr\^ace \`a $\lambda=c/\nu$ (attention : $\lambda$ en m\`etres !).
  \item[$\square$] Je sais convertir les \'electronvolts en joules : $\SI{1}{eV}=\SI{1,60e-19}{J}$.
  \item[$\square$] Je sais calculer $\Delta E = h\nu = hc/\lambda$ entre deux niveaux et dire s'il s'agit d'\'emission ou d'absorption.
  \item[$\square$] Je sais qu'un photon n'est absorb\'e que si son \'energie correspond \emph{exactement} \`a un \'ecart entre deux niveaux.
  \item[$\square$] Je sais d\'ecrire un profil spectral et utiliser la loi de Wien pour d\'eterminer une temp\'erature.
  \item[$\square$] Je sais expliquer la forme du spectre solaire : fond continu + raies sombres d'absorption (Fraunhofer).
  \item[$\square$] Je sais que les raies spectrales sont la << carte d'identit\'e >> des \'el\'ements chimiques.
  \item[$\square$] Je v\'erifie toujours l'homog\'en\'eit\'e de mes calculs et le nombre de chiffres significatifs.
\end{itemize}
\end{aretenir}

\findecours

\end{document}
